% First Order Differential Equations With Separable Variables — Pure Mathematics with Statistics Upper Sixth — Cours signé OnBuch+
% Official MINESEC syllabus. Compiler avec : tectonic PMS05-first-order-differential-equations-with-separable-variables.tex
\newcommand{\DOCMATIERE}{Pure Mathematics with Statistics}
\newcommand{\DOCNIVEAU}{Upper Sixth}
\newcommand{\DOCMODULE}{Upper Sixth — Pure mathematics with statistics}
\newcommand{\DOCLECON}{5}
\newcommand{\DOCTITRE}{First Order Differential Equations With Separable Variables}
% OnBuch+ — préambule « cours signés » ENRICHI (compilé avec tectonic / XeLaTeX)
% Système visuel « Pop » d'OnBuch+ : fond crème (jamais blanc), bordures encre
% épaisses, ombres « dures » décalées, titres Archivo Black en capitales.
% Version MATHÉMATIQUES : graphes (pgfplots/tikz), boîtes pédagogiques
% (définition, propriété, méthode, attention, le savais-tu, exercice résolu,
% exercice, TP, bilan), page de garde et signature OnBuch+ ancrée en bas.
%
% Macros attendues dans main.tex AVANT \input{preamble} :
%   \DOCMATIERE  \DOCNIVEAU  \DOCMODULE  \DOCLECON (numéro)  \DOCTITRE
\documentclass[11pt]{article}

\usepackage{fontspec}
\usepackage[english]{babel}
\usepackage[a4paper,margin=2cm,top=3.4cm,bottom=2.8cm,headheight=1.4cm,footskip=1.3cm]{geometry}
\usepackage{amsmath,amssymb,amsfonts}
\usepackage{mathtools} % \xmapsto, \coloneqq... souvent utilisés par les modèles
\usepackage{cancel} % simplifications barrées (bilans de forces)
% \abs{z} (module, valeur absolue) et \norm{u} (norme), utilisés par les modèles
\providecommand{\abs}{}\let\abs\relax\DeclarePairedDelimiter{\abs}{\lvert}{\rvert}
\providecommand{\norm}{}\let\norm\relax\DeclarePairedDelimiter{\norm}{\lVert}{\rVert}
\usepackage[table]{xcolor} % [table] : fournit aussi \rowcolors (alternance de lignes)
\usepackage{pagecolor}
\usepackage{tikz}
\usetikzlibrary{arrows.meta,positioning,calc,shapes.geometric,shapes.misc,shapes.arrows,decorations.pathmorphing,decorations.pathreplacing,decorations.markings,patterns,shadows,mindmap,trees,fit,backgrounds,matrix,intersections,angles,quotes}
\usepackage{pgfplots}
\usepackage[version=4]{mhchem} % équations nucléaires \ce{^{235}_{92}U}
\usepackage[european,siunitx]{circuitikz} % schémas électriques
\pgfplotsset{compat=1.18}
\usepgfplotslibrary{fillbetween,groupplots} % aires entre courbes (calcul intégral)
\usepackage{enumitem}
\usepackage{fancyhdr}
% [most] charge toutes les bibliothèques tcolorbox (listings, minted, raster,
% documentation...) et gonfle inutilement la mémoire TeX sur un document long
% (~300 boîtes) : on ne charge que ce qui est réellement utilisé.
\usepackage[skins,breakable]{tcolorbox}
\usepackage{graphicx}
\usepackage{colortbl}
\usepackage{booktabs}
\usepackage{tabularx}
\usepackage{diagbox} % case d'en-tête diagonale des tableaux à double entrée
% Colonnes extensibles centrée / gauche / droite (C, L, R), souvent utilisées par les modèles
\newcolumntype{C}{>{\centering\arraybackslash}X}
\newcolumntype{L}{>{\raggedright\arraybackslash}X}
\newcolumntype{R}{>{\raggedleft\arraybackslash}X}
\usepackage{array}
\usepackage{multicol}
\usepackage{siunitx}
% parse-numbers=false : accepte aussi \SI{2,5 \times 10^{-3}}{...}
\sisetup{locale=UK,output-decimal-marker={.},group-minimum-digits=4,parse-numbers=false}
\DeclareSIUnit\ohm{\ensuremath{\symup{\Omega}}} % U+2126 absent de la police
\DeclareSIUnit\division{div} % réglages d'oscilloscope (ms/div, V/div)
\DeclareSIUnit\tr{tr} % tours (vitesse de rotation, ex. \SI{1500}{\tr\per\minute})
\usepackage{qrcode}
% \degree : commande fréquemment utilisée par les modèles pour un angle ou une
% température en dehors de \SI{}{} (non fournie par les paquets chargés).
\providecommand{\degree}{^{\circ}}
% \qed (fin de démonstration), utilisé par les modèles sans amsthm
\providecommand{\qed}{\hfill$\square$}
% \barre : double barre des valeurs interdites dans un tableau de variations (tkz-tab)
\providecommand{\barre}{\ensuremath{\|}}
% environnement proof (amsthm non chargé) utilisé par les modèles
\ifdefined\proof\else\newenvironment{proof}[1][Proof]{\par\noindent\textit{#1.}\ }{\qed\par}\fi
\usepackage{lastpage}
\usepackage{hyperref}
\hypersetup{hidelinks}

% ============================================================
% Palette « Pop » OnBuch+ — mêmes hex que la classe P (plus_ui.dart)
% ============================================================
\definecolor{poporange}{HTML}{F59321}
\definecolor{popdark}{HTML}{E07A0C}
\definecolor{popcream}{HTML}{FFF6E8}
\definecolor{popink}{HTML}{1C1714}
\definecolor{poppurple}{HTML}{7A5AE0}
\definecolor{popgreen}{HTML}{2E9E6A}
\definecolor{poppink}{HTML}{E0457B}
\definecolor{popblue}{HTML}{2F7DE1}
\definecolor{popgold}{HTML}{C98A12}
\definecolor{popmuted}{HTML}{8A7F76}
\definecolor{popline}{HTML}{EADFD2}
\definecolor{popgray}{HTML}{8A7F76}
\colorlet{popgrey}{popgray}
% Alias tolérants (le modèle nomme parfois les couleurs différemment)
\colorlet{popwhite}{white}
\colorlet{popblack}{popink}
\colorlet{popred}{poppink}
\colorlet{popyellow}{popgold}
% Teintes claires pour fonds de tableaux / zones
\colorlet{popblueL}{popblue!10!white}
\colorlet{poporangeL}{poporange!14!white}
\colorlet{popgreenL}{popgreen!12!white}
\colorlet{poppurpleL}{poppurple!12!white}
\colorlet{poppinkL}{poppink!12!white}
\colorlet{popgoldL}{popgold!14!white}

\pagecolor{popcream}

% ============================================================
% Typographie
% ============================================================
\defaultfontfeatures{Ligatures=TeX}
\setmainfont{PlusJakartaSans-Medium.ttf}[Path=../../../fonts/,BoldFont=PlusJakartaSans-Bold.ttf]
% Maths en sans-serif (Fira Math) avec chiffres et lettres droites de Plus
% Jakarta, pour que les formules se fondent dans le texte.
\usepackage[math-style=ISO,bold-style=ISO]{unicode-math}
\setmathfont{FiraMath-Regular.otf}[Path=../../../fonts/]
\setmathfont{PlusJakartaSans-Medium.ttf}[Path=../../../fonts/,range={up/num,up/{Latin,latin},\mathup}]
% (icomma retiré : décimales avec point en anglais)
% Caractères mathématiques Unicode absents des polices de texte (Plus Jakarta,
% Archivo Black) : rendus par leur équivalent mathématique, en texte comme en
% formule — sinon ils s'affichent en carré vide (ex. « Arithmétique dans ℤ »).
\usepackage{newunicodechar}
\newunicodechar{ℝ}{\ensuremath{\mathbb{R}}}
\newunicodechar{ℤ}{\ensuremath{\mathbb{Z}}}
\newunicodechar{ℂ}{\ensuremath{\mathbb{C}}}
\newunicodechar{ℕ}{\ensuremath{\mathbb{N}}}
\newunicodechar{ℚ}{\ensuremath{\mathbb{Q}}}
\newunicodechar{𝔻}{\ensuremath{\mathbb{D}}}
\newunicodechar{α}{\ensuremath{\alpha}}
\newunicodechar{β}{\ensuremath{\beta}}
\newunicodechar{γ}{\ensuremath{\gamma}}
\newunicodechar{δ}{\ensuremath{\delta}}
\newunicodechar{ε}{\ensuremath{\varepsilon}}
\newunicodechar{θ}{\ensuremath{\theta}}
\newunicodechar{λ}{\ensuremath{\lambda}}
\newunicodechar{μ}{\ensuremath{\mu}}
\newunicodechar{σ}{\ensuremath{\sigma}}
\newunicodechar{τ}{\ensuremath{\tau}}
\newunicodechar{φ}{\ensuremath{\varphi}}
\newunicodechar{ψ}{\ensuremath{\psi}}
\newunicodechar{ω}{\ensuremath{\omega}}
\newunicodechar{Σ}{\ensuremath{\Sigma}}
% majuscules produites par \MakeUppercase dans les titres (α-aminés -> Α)
\newunicodechar{Α}{\ensuremath{\alpha}}
\newunicodechar{Β}{\ensuremath{\beta}}
\newunicodechar{Γ}{\ensuremath{\Gamma}}
\newunicodechar{Δ}{\ensuremath{\Delta}}
\newunicodechar{Ε}{\ensuremath{\varepsilon}}
\newunicodechar{Θ}{\ensuremath{\Theta}}
\newunicodechar{Λ}{\ensuremath{\Lambda}}
\newunicodechar{Μ}{\ensuremath{\mu}}
\newunicodechar{Τ}{\ensuremath{\tau}}
\newunicodechar{Φ}{\ensuremath{\Phi}}
\newunicodechar{Ψ}{\ensuremath{\Psi}}
\newunicodechar{Ω}{\ensuremath{\Omega}}
\newunicodechar{↦}{\ensuremath{\mapsto}}
\newunicodechar{∈}{\ensuremath{\in}}
\newunicodechar{∉}{\ensuremath{\notin}}
\newunicodechar{∧}{\ensuremath{\wedge}}
\newunicodechar{⇔}{\ensuremath{\Leftrightarrow}}
\newunicodechar{⟺}{\ensuremath{\Longleftrightarrow}}
\newunicodechar{⇒}{\ensuremath{\Rightarrow}}
\newunicodechar{⟹}{\ensuremath{\Longrightarrow}}
\newunicodechar{∘}{\ensuremath{\circ}}
\newunicodechar{∪}{\ensuremath{\cup}}
\newunicodechar{∩}{\ensuremath{\cap}}
\newunicodechar{≡}{\ensuremath{\equiv}}
\newunicodechar{⊂}{\ensuremath{\subset}}
\newunicodechar{⊥}{\ensuremath{\perp}}
\newunicodechar{∃}{\ensuremath{\exists}}
\newunicodechar{∀}{\ensuremath{\forall}}
\newunicodechar{∛}{\ensuremath{\sqrt[3]{\,}}}
\newunicodechar{∜}{\ensuremath{\sqrt[4]{\,}}}
\newunicodechar{⃗}{\ensuremath{{}^{\to}}}
\newunicodechar{ⁿ}{\ifmmode^{n}\else\textsuperscript{\ensuremath{n}}\fi}
\newunicodechar{ˣ}{\ifmmode^{x}\else\textsuperscript{\ensuremath{x}}\fi}
\newunicodechar{ᵇ}{\ifmmode^{b}\else\textsuperscript{\ensuremath{b}}\fi}
\newunicodechar{ʳ}{\ifmmode^{r}\else\textsuperscript{\ensuremath{r}}\fi}
\newunicodechar{ᵘ}{\ifmmode^{u}\else\textsuperscript{\ensuremath{u}}\fi}
\newunicodechar{ʸ}{\ifmmode^{y}\else\textsuperscript{\ensuremath{y}}\fi}
\newunicodechar{ᵃ}{\ifmmode^{a}\else\textsuperscript{\ensuremath{a}}\fi}
\newunicodechar{ᵉ}{\ifmmode^{e}\else\textsuperscript{\ensuremath{e}}\fi}
\newunicodechar{ᵏ}{\ifmmode^{k}\else\textsuperscript{\ensuremath{k}}\fi}
\newunicodechar{ᵖ}{\ifmmode^{p}\else\textsuperscript{\ensuremath{p}}\fi}
\newunicodechar{ᴺ}{\ifmmode^{N}\else\textsuperscript{\ensuremath{N}}\fi}
\newunicodechar{ᶜ}{\ifmmode^{c}\else\textsuperscript{\ensuremath{c}}\fi}
\newunicodechar{ʰ}{\ifmmode^{h}\else\textsuperscript{\ensuremath{h}}\fi}
\newunicodechar{ᵠ}{\ifmmode^{q}\else\textsuperscript{\ensuremath{q}}\fi}
\newunicodechar{ₙ}{\ifmmode_{n}\else\textsubscript{\ensuremath{n}}\fi}
\newunicodechar{ₐ}{\ifmmode_{a}\else\textsubscript{\ensuremath{a}}\fi}
\newunicodechar{ᵢ}{\ifmmode_{i}\else\textsubscript{\ensuremath{i}}\fi}
\newunicodechar{ᵣ}{\ifmmode_{r}\else\textsubscript{\ensuremath{r}}\fi}
\newunicodechar{ₖ}{\ifmmode_{k}\else\textsubscript{\ensuremath{k}}\fi}
\newunicodechar{ᵦ}{\ifmmode_{\beta}\else\textsubscript{\ensuremath{\beta}}\fi}
\newunicodechar{ₓ}{\ifmmode_{x}\else\textsubscript{\ensuremath{x}}\fi}
\newfontfamily\popdisplay{ArchivoBlack-Regular.ttf}[Path=../../../fonts/]
\newfontfamily\poplogo{SpaceGrotesk-Bold.ttf}[Path=../../../fonts/]
\newfontfamily\popsemi{PlusJakartaSans-SemiBold.ttf}[Path=../../../fonts/]
\newfontfamily\popxbold{PlusJakartaSans-ExtraBold.ttf}[Path=../../../fonts/]
\newfontfamily\popxboldls{PlusJakartaSans-ExtraBold.ttf}[Path=../../../fonts/,LetterSpace=3.0]

\setlist[enumerate]{leftmargin=1.7em,itemsep=5pt,topsep=4pt}
\setlist[itemize]{leftmargin=1.7em,itemsep=4pt,topsep=4pt,label={\color{poporange}\textbullet}}
\setlength{\parindent}{0pt}
\setlength{\parskip}{5pt}
\renewcommand{\arraystretch}{1.25}

% pgfplots : style Pop par défaut
\pgfplotsset{
  every axis/.append style={axis line style={popink,line width=1pt},tick style={popink},
    grid style={popline,line width=0.5pt},label style={font=\small},tick label style={font=\footnotesize}},
  popcurve/.style={poporange,line width=1.8pt,no markers,smooth},
}
% Même style disponible dans un \draw TikZ simple (hors pgfplots)
\tikzset{popcurve/.style={poporange,line width=1.8pt}}
\tikzset{popgrid/.style={popline,very thin}}
\colorlet{popcurve}{poporange} % le nom du style est parfois utilisé comme couleur

% ============================================================
% Pill / squiggle
% ============================================================
\newcommand{\poppill}[3]{%
  \tikz[baseline=-0.55ex]\node[draw=popink,line width=1.2pt,rounded corners=2.6pt,%
    fill=#2,inner xsep=6.5pt,inner ysep=3pt]{%
    \popxboldls\fontsize{7}{7}\selectfont\color{#3}\MakeUppercase{#1}};%
}
\newcommand{\popsquiggle}[1]{%
  \tikz[baseline=-0.4ex]{
    \draw[#1,line width=2.6pt,line cap=round]
      plot [smooth] coordinates {(0,0.09) (0.34,0.01) (0.68,0.21) (1.02,0.01) (1.36,0.10)};
  }%
}
% Mot-clé surligné (marqueur orange sous le texte)
\newcommand{\cle}[1]{{\popxbold\color{popdark}#1}}

% ============================================================
% En-tête / pied
% ============================================================
\pagestyle{fancy}
\fancyhf{}
\renewcommand{\headrulewidth}{0pt}
\renewcommand{\footrulewidth}{0pt}
\lhead{\raisebox{-0.15ex}{\poplogo\fontsize{13}{13}\selectfont\color{popink}OnBuch\color{poporange}+}}
\rhead{\poppill{\DOCMATIERE\ \textbullet\ \DOCNIVEAU\ \textbullet\ Chapter \DOCLECON}{white}{popink}}
\lfoot{\popxbold\fontsize{7.6}{7.6}\selectfont\color{popmuted}\MakeUppercase{Signed course OnBuch+}\ \textbullet\ {\popsemi\DOCTITRE}}
\rfoot{\tikz[baseline=-0.6ex]\node[rounded corners=8.5pt,fill=popink,minimum height=17pt,minimum width=17pt,inner xsep=5pt,inner sep=0]{\popxbold\fontsize{8}{8}\selectfont\color{popcream}\thepage\,/\,\pageref*{LastPage}};}

% ============================================================
% Titres
% ============================================================
\newcommand{\chaptitre}[2]{%
  \begin{tcolorbox}[enhanced,colback=poporange,colframe=popink,boxrule=2.6pt,arc=11pt,
      drop shadow=popink,left=15pt,right=15pt,top=12pt,bottom=11pt,before skip=3pt,after skip=18pt]
    \poppill{#2}{popink}{popcream}\par\vspace{9pt}
    {\raggedright\hyphenpenalty=10000\exhyphenpenalty=10000\popdisplay\fontsize{22}{24}\selectfont\color{popcream}\MakeUppercase{#1}\par}
  \end{tcolorbox}
}
% Section numérotée : pastille orange avec le numéro + titre Archivo
\newcounter{popsec}
\newcommand{\coursec}[1]{%
  \stepcounter{popsec}\setcounter{popsub}{0}%
  \par\vspace{16pt}\noindent
  \begin{minipage}{\linewidth}
  \tikz[baseline=-0.7ex]\node[draw=popink,line width=1.6pt,fill=poporange,rounded corners=4pt,minimum size=22pt,inner sep=0]{\popdisplay\fontsize{12}{12}\selectfont\color{popcream}\thepopsec};%
  \hspace{8pt}{\popdisplay\fontsize{15}{17}\selectfont\color{popink}\MakeUppercase{#1}}\par
  \vspace{4pt}\noindent{\color{poporange}\rule{36pt}{3.6pt}}
  \end{minipage}\par\nopagebreak\vspace{8pt}
}
\newcounter{popsub}
\newcommand{\courssub}[1]{%
  \stepcounter{popsub}%
  \par\vspace{9pt}\noindent{\popxbold\fontsize{11.5}{13}\selectfont\color{popink}{\color{poporange}\thepopsec.\thepopsub}\hspace{6pt}#1}\par\nopagebreak\vspace{4pt}
}

% ============================================================
% Boîtes pédagogiques — même recette : fond blanc, bordure épaisse colorée,
% ombre dure encre, badge plein. Titre optionnel [#1].
% ============================================================
\tcbset{popbase/.style={breakable,enhanced,colback=white,boxrule=2.3pt,arc=9pt,
  shadow={3pt}{-3pt}{0pt}{popink},left=14pt,right=14pt,top=11pt,bottom=11pt,before skip=11pt,after skip=15pt,
  fontupper=\popsemi\fontsize{10.6}{14.5}\selectfont\color{popink},
  fontlower=\popsemi\fontsize{10.6}{14.5}\selectfont\color{popink}}}
% Coche dessinée (le glyphe ✓ n'existe pas dans les polices du document)
\newcommand{\popcheck}[1]{\tikz[baseline=-0.5ex]\draw[#1,line width=1.6pt,line cap=round,line join=round](-0.13,0.0)--(-0.04,-0.09)--(0.13,0.1);}
\providecommand{\coche}{\popcheck{popgreen}}
\providecommand{\resultat}[1]{\par\smallskip\noindent\fcolorbox{popgreen}{popgreenL}{\parbox{\dimexpr\linewidth-2\fboxsep-2\fboxrule\relax}{\textbf{Result.} #1}}\par\smallskip}
\newcommand{\popboxhead}[3]{\poppill{#1}{#2}{white}\ifx\relax#3\relax\else\hspace{6pt}{\popxbold\fontsize{10.6}{12}\selectfont\color{popink}#3}\fi\par\vspace{7pt}}

\newtcolorbox{aretenir}[1][]{popbase,colframe=popgreen,before upper={\popboxhead{Key points}{popgreen}{#1}}}
\newtcolorbox{exemplebox}[1][]{popbase,colframe=poporange,before upper={\popboxhead{Exemple}{poporange}{#1}}}
\newtcolorbox{definition}[1][]{popbase,colframe=popblue,before upper={\popboxhead{Definition}{popblue}{#1}}}
\newtcolorbox{propriete}[1][]{popbase,colframe=poppurple,before upper={\popboxhead{Property}{poppurple}{#1}}}
\newtcolorbox{methode}[1][]{popbase,colframe=popgold,before upper={\popboxhead{Method}{popgold}{#1}}}
\newtcolorbox{attention}[1][]{popbase,colframe=poppink,before upper={\popboxhead{Warning}{poppink}{#1}}}
\newtcolorbox{experience}[1][]{popbase,colframe=popblue,colback=popblueL,before upper={\popboxhead{Experiment}{popblue}{#1}}}
% « Le savais-tu ? » : carte violette pleine (contexte, culture, Cameroun)
\newtcolorbox{savaistu}[1][]{popbase,colback=poppurple,colframe=popink,
  fontupper=\popsemi\fontsize{10.4}{14.2}\selectfont\color{white},
  before upper={\poppill{Did you know?}{white}{poppurple}\ifx\relax#1\relax\else\hspace{6pt}{\popxbold\color{white}#1}\fi\par\vspace{7pt}}}
% Exercice résolu : énoncé en haut, \tcblower puis la solution sur fond crème
\newcounter{popexo}\newcounter{popexr}
\newtcolorbox{exoresolu}[1][]{popbase,colframe=poporange,colbacklower=popcream,
  segmentation style={popink,line width=1.2pt,dashed},
  before upper={\stepcounter{popexr}\popboxhead{Worked example \thepopexr}{poporange}{#1}},
  before lower={\poppill{Solution}{popink}{popcream}\par\vspace{6pt}}}
% Exercice (non résolu en place) — la correction est dans la partie Corrigés
\newtcolorbox{exercice}[1][]{popbase,colframe=popink,
  before upper={\stepcounter{popexo}\popboxhead{Exercise \thepopexo}{popink}{#1}}}
% Corrigé
\newtcolorbox{corrige}[1][]{popbase,colframe=popgreen,colback=popgreenL,
  before upper={\popboxhead{Solution}{popgreen}{#1}}}
% Objectifs / prérequis / bilan
\newtcolorbox{objectifs}{popbase,colframe=popink,colback=poporangeL,
  before upper={\popboxhead{Chapter objectives}{popink}{}}}
\newtcolorbox{prerequis}{popbase,colframe=popink,colback=popblueL,
  before upper={\popboxhead{Prerequisites}{popblue}{}}}
\newtcolorbox{bilan}{popbase,colframe=popink,colback=white,boxrule=2.8pt,
  before upper={\popboxhead{Fiche bilan}{poporange}{L'essentiel à connaître pour l'examen}}}
% Figure encadrée avec légende
\newtcolorbox{popfigure}[1][]{enhanced,breakable=false,colback=white,colframe=popline,boxrule=1.4pt,arc=8pt,
  left=10pt,right=10pt,top=10pt,bottom=8pt,before skip=10pt,after skip=12pt,halign=center,
  lower separated=false,halign lower=center,
  fontlower=\popsemi\fontsize{9}{11.5}\selectfont\color{popmuted},
  #1}
\newcommand{\legende}[1]{\par\vspace{6pt}{\popsemi\fontsize{9}{11.5}\selectfont\color{popmuted}#1}}

% ============================================================
% Signature OnBuch+
% ============================================================
\newcommand{\poplogoinline}[1]{{\poplogo\fontsize{#1}{#1}\selectfont\color{popink}OnBuch\color{poporange}+}}
\newcommand{\signatureonbuch}{%
  \begin{tcolorbox}[enhanced,colback=popink,colframe=popink,boxrule=2.2pt,arc=10pt,
      drop shadow=poporange,left=14pt,right=14pt,top=10pt,bottom=10pt,before skip=0pt,after skip=0pt]
    \tikz[baseline=-0.7ex]\node[circle,fill=popgreen,draw=popcream,line width=1.3pt,minimum size=22pt,inner sep=0]{\popcheck{white}};%
    \hspace{9pt}\begin{minipage}[c]{0.62\linewidth}
      {\popxbold\fontsize{12}{13}\selectfont\color{popcream}Signed\ \textbullet\ {\poplogo OnBuch\color{poporange}+}}\par\vspace{2pt}
      {\raggedright\popsemi\fontsize{8}{10.5}\selectfont\color{popline}\MakeUppercase{OnBuch+ teaching team}\\\MakeUppercase{Follows the official MINESEC syllabus}\par}
    \end{minipage}\hfill
    {\poplogo\fontsize{22}{22}\selectfont\color{popcream}OnBuch\color{poporange}+}
  \end{tcolorbox}%
}

% ============================================================
% Page de garde
% #1 titre · #2 matière · #3 niveau · #4 module · #5 n° leçon · #6 durée ·
% #7 description / objectifs (texte ou itemize)
% ============================================================
\newcommand{\pagegarde}[7]{%
  \thispagestyle{empty}
  \newgeometry{a4paper,margin=1.7cm,top=1.6cm,bottom=1.4cm}%
  \begin{tikzpicture}[remember picture,overlay]
    \fill[poporange!18] ([xshift=-3.2cm,yshift=-3.6cm]current page.north east) circle (4.6cm);
    \fill[poppurple!14] ([xshift=2.4cm,yshift=7.6cm]current page.south west) circle (3.8cm);
  \end{tikzpicture}%
  {\poplogo\fontsize{30}{30}\selectfont\color{popink}OnBuch\color{poporange}+}\hfill
  \raisebox{0.6ex}{\poppill{Signed course \textbullet\ Premium}{popink}{popcream}}\par
  \vspace{1pt}\popsquiggle{poppurple}\par
  \vspace{8pt}
  \begin{tcolorbox}[enhanced,colback=poporange,colframe=popink,boxrule=2.8pt,arc=13pt,
      drop shadow=popink,left=18pt,right=18pt,top=12pt,bottom=13pt,after skip=10pt]
    \poppill{#2 \textbullet\ #3}{popink}{popcream}\hspace{5pt}\poppill{#4}{white}{popink}\par\vspace{8pt}
    {\popdisplay\fontsize{14}{14}\selectfont\color{popink}CHAPTER #5}\par\vspace{4pt}
    {\raggedright\hyphenpenalty=10000\exhyphenpenalty=10000\popdisplay\fontsize{25}{27}\selectfont\color{popcream}\MakeUppercase{#1}\par}
    \vspace{8pt}
    {\popxbold\fontsize{9.5}{9.5}\selectfont\color{popink}\MakeUppercase{Suggested time: #6}}
  \end{tcolorbox}
  \begin{tcolorbox}[enhanced,colback=white,colframe=popink,boxrule=2.2pt,arc=10pt,
      drop shadow=popink,left=16pt,right=16pt,top=10pt,bottom=10pt,after skip=8pt]
    \poppill{In this chapter}{poppurple}{white}\par\vspace{5pt}
    \popsemi\fontsize{9.3}{12.6}\selectfont\color{popink}#7
  \end{tcolorbox}
  \noindent\begin{minipage}[t]{0.487\linewidth}
    \begin{tcolorbox}[enhanced,colback=popgreen,colframe=popink,boxrule=2.2pt,arc=10pt,
        drop shadow=popink,left=11pt,right=11pt,top=10pt,bottom=10pt,height=3.1cm,valign=center]
      \tcbox[colback=white,colframe=popink,boxrule=1.3pt,arc=5pt,boxsep=0pt,nobeforeafter,
        left=3pt,right=3pt,top=3pt,bottom=3pt,valign=center]{\qrcode[height=1.9cm]{https://play.google.com/store/apps/details?id=com.luvvix.onbuch&hl=en}}%
      \hspace{8pt}\begin{minipage}[c]{0.5\linewidth}
        {\popdisplay\fontsize{11}{12.5}\selectfont\color{white}\MakeUppercase{Download OnBuch}}\par\vspace{4pt}
        {\raggedright\popsemi\fontsize{8.4}{11}\selectfont\color{white}Free \textbullet\ Google Play\\AI tutor, past papers, results\par}
      \end{minipage}
    \end{tcolorbox}
  \end{minipage}\hfill
  \begin{minipage}[t]{0.487\linewidth}
    \begin{tcolorbox}[enhanced,colback=poppurple,colframe=popink,boxrule=2.2pt,arc=10pt,
        drop shadow=popink,left=11pt,right=11pt,top=10pt,bottom=10pt,height=3.1cm,valign=center]
      \tikz[baseline=-0.7ex]\node[circle,fill=white,draw=popink,line width=1.3pt,minimum size=1.6cm,inner sep=0]{\popdisplay\fontsize{24}{24}\selectfont\color{poppurple}+};%
      \hspace{8pt}\begin{minipage}[c]{0.6\linewidth}
        {\popdisplay\fontsize{11}{12.5}\selectfont\color{white}\MakeUppercase{Go OnBuch+}}\par\vspace{4pt}
        {\raggedright\popsemi\fontsize{8.4}{11}\selectfont\color{white}All signed courses, unlimited AI tutor, PDF sheets — OnBuch+ tab in the app\par}
      \end{minipage}
    \end{tcolorbox}
  \end{minipage}
  \vspace{8pt}
  \popcontenu
  \vfill
  \signatureonbuch
  \restoregeometry
  \newpage
}

% Rangée « Ce cours contient » (page de garde)
\newcommand{\poptile}[3]{% #1 couleur #2 gros texte #3 légende
  \begin{tcolorbox}[enhanced,colback=#1,colframe=popink,boxrule=2pt,arc=9pt,shadow={2.5pt}{-2.5pt}{0pt}{popink},
      left=6pt,right=6pt,top=9pt,bottom=9pt,halign=center,valign=center,height=1.75cm,width=\linewidth,nobeforeafter]
    {\popdisplay\fontsize{12.5}{13}\selectfont\color{white}\MakeUppercase{#2}}\par\vspace{3pt}
    {\centering\popsemi\fontsize{7.8}{9.5}\selectfont\color{white}#3\par}
  \end{tcolorbox}}
\newcommand{\popcontenu}{%
  \par\noindent{\popxboldls\fontsize{7.5}{7.5}\selectfont\color{popmuted}THIS SIGNED COURSE CONTAINS}\par\vspace{7pt}
  \noindent\begin{minipage}[t]{0.235\linewidth}\poptile{popblue}{Course}{complete, illustrated, step by step}\end{minipage}\hfill
  \begin{minipage}[t]{0.235\linewidth}\poptile{poporange}{Methods}{and worked examples}\end{minipage}\hfill
  \begin{minipage}[t]{0.235\linewidth}\poptile{poppink}{Exercises}{exam style + solutions}\end{minipage}\hfill
  \begin{minipage}[t]{0.235\linewidth}\poptile{popgreen}{Summary sheet}{the essentials to remember}\end{minipage}\par}

% Sommaire
\newtcolorbox{sommaire}{enhanced,colback=white,colframe=popink,boxrule=2.2pt,arc=10pt,
  drop shadow=popink,left=18pt,right=18pt,top=13pt,bottom=13pt,before skip=6pt,after skip=20pt,
  before upper={\poppill{Contents}{poppurple}{white}\par\vspace{9pt}\popsemi\fontsize{10.6}{14.5}\selectfont\color{popink}}}

% Signature de fin de cours — poussée tout en bas de la dernière page
\newcommand{\findecours}{%
  \par\vfill\nopagebreak
  \begin{center}\popsquiggle{poporange}\end{center}\vspace{4pt}
  \signatureonbuch
}
\AtBeginDocument{\def\llbracket{\mathopen{[\mkern-3.5mu[}}\def\rrbracket{\mathclose{]\mkern-3.5mu]}}} % intervalles d'entiers (glyphe ⟦ absent de la police)
% Glyphes absents de Fira Math : redéfinis à partir de glyphes disponibles
\AtBeginDocument{%
  \def\setminus{\mathbin{\backslash}}\let\smallsetminus\setminus
  \def\vdots{\mathord{\vbox{\baselineskip3.2pt\lineskiplimit0pt\kern4pt\hbox{.}\hbox{.}\hbox{.}}}}%
  \def\ddots{\mathinner{\mkern1mu\raise6.5pt\hbox{.}\mkern2mu\raise3.6pt\hbox{.}\mkern2mu\raise0.7pt\hbox{.}\mkern1mu}}%
  \def\triangle{\mathord{\scalebox{0.9}{$\symup{\Delta}$}}}%
  \def\nabla{\mathord{\raisebox{\depth}{\scalebox{1}[-1]{$\symup{\Delta}$}}}}%
  \def\checkmark{\popcheck{popdark}}%
  \def\star{\mathbin{\ast}}\def\bigstar{\mathord{\ast}}%
  \def\diamond{\mathbin{\scalebox{0.6}{$\Diamond$}}}%
  \def\bigcup{\mathop{\mathchoice{\raisebox{-0.3em}{\scalebox{1.7}{$\cup$}}}{\scalebox{1.25}{$\cup$}}{\cup}{\cup}}\displaylimits}%
  \def\bigcap{\mathop{\mathchoice{\raisebox{-0.3em}{\scalebox{1.7}{$\cap$}}}{\scalebox{1.25}{$\cap$}}{\cap}{\cap}}\displaylimits}%
  \def\lesseqqgtr{\mathrel{\lessgtr}}\def\bigoplus{\mathop{\oplus}\displaylimits}%
}
\providecommand{\ssi}{if and only if} % abréviation fréquente
\AtBeginDocument{\providecommand{\wideparen}{\overparen}} % arc de cercle

% Fonctions usuelles que les modèles utilisent sans que amsmath les fournisse
\providecommand{\arsinh}{\operatorname{arsinh}}\providecommand{\arcosh}{\operatorname{arcosh}}
\providecommand{\artanh}{\operatorname{artanh}}\providecommand{\arsech}{\operatorname{arsech}}
\providecommand{\arcsch}{\operatorname{arcsch}}\providecommand{\arcoth}{\operatorname{arcoth}}
\providecommand{\arcsinh}{\operatorname{arsinh}}\providecommand{\arccosh}{\operatorname{arcosh}}
\providecommand{\arctanh}{\operatorname{artanh}}\providecommand{\sech}{\operatorname{sech}}
\providecommand{\csch}{\operatorname{csch}}\providecommand{\arcsec}{\operatorname{arcsec}}
\providecommand{\arccsc}{\operatorname{arccsc}}\providecommand{\arccot}{\operatorname{arccot}}
\providecommand{\sgn}{\operatorname{sgn}}\providecommand{\lcm}{\operatorname{lcm}}
\providecommand{\Var}{\operatorname{Var}}\providecommand{\Cov}{\operatorname{Cov}}

%% FIGFIX-BEGIN v5 — correctifs globaux des figures (docs/onbuch-generation/tools/figfix)
\usepackage{adjustbox}\usepackage{etoolbox}\tcbuselibrary{raster}
% toute figure TikZ/pgfplots plus large que la ligne est réduite à la largeur disponible
\BeforeBeginEnvironment{tikzpicture}{\begin{adjustbox}{max width=\linewidth}}
\AfterEndEnvironment{tikzpicture}{\end{adjustbox}}
% légendes pgfplots lisibles par-dessus les courbes
\pgfplotsset{every axis/.append style={legend style={fill=white,fill opacity=0.92,text opacity=1,draw=black!15,inner sep=3pt}}}
% étiquettes d'arêtes (syntaxe edge["..."]) avec fond blanc
\tikzset{every edge quotes/.append style={fill=white,inner sep=1.2pt}}
%% FIGFIX-END
\begin{document}

\pagegarde{First Order Differential Equations With Separable Variables}{Pure Mathematics with Statistics}{Upper Sixth}{Upper Sixth — Pure mathematics with statistics}{5}{6 periods}{This chapter introduces first‑order differential equations with separable variables, guiding students from the basic definition to real‑life applications. Mastery of separation, integration, and initial‑condition techniques is essential for success in the GCE Advanced Level Pure Mathematics with Statistics examination.
\vspace{4pt}
\begin{multicols}{2}\small
\begin{itemize}
\item Define a first‑order differential equation and identify its order and degree.
\item Recognise when a differential equation is separable.
\item Apply the method of separation of variables to rewrite the equation in integrable form.
\item Integrate both sides correctly to obtain the general solution, including the constant of integration.
\item Use given initial conditions to determine the particular solution.
\item Interpret the solution in the context of a real‑world problem (population growth, cooling, chemical reactions, etc.).
\end{itemize}\end{multicols}}

\begin{sommaire}
\begin{multicols}{2}
\begin{enumerate}
\item 1. Meaning and Classification of First‑Order Differential Equations
\item 2. Recognising Separable Equations
\item 3. Separation of Variables and Integration to Obtain General Solutions
\item 4. Particular Solutions Using Initial Conditions
\item 5. Applications of Separable Differential Equations and Exam Preparation
\item Integration activity
\item Methods and worked examples
\item Exercises
\item Solutions to the exercises
\item Summary sheet
\end{enumerate}
\end{multicols}
\end{sommaire}

\begin{objectifs}
\begin{itemize}
\item Define a first‑order differential equation and identify its order and degree.
\item Recognise when a differential equation is separable.
\item Apply the method of separation of variables to rewrite the equation in integrable form.
\item Integrate both sides correctly to obtain the general solution, including the constant of integration.
\item Use given initial conditions to determine the particular solution.
\item Interpret the solution in the context of a real‑world problem (population growth, cooling, chemical reactions, etc.).
\item Check the validity of a solution by differentiation and by verifying the initial condition.
\item Identify common algebraic and integration pitfalls and avoid them in exam answers.
\end{itemize}
\end{objectifs}

\begin{prerequis}
\begin{itemize}
\item Basic differentiation rules (power, exponential, trigonometric, logarithmic).
\item Indefinite integration techniques (substitution, partial fractions, standard integrals).
\item Algebraic manipulation of fractions and logarithms.
\item Understanding of functions, domains, and implicit vs explicit forms.
\item Experience with simple modelling problems from Lower Sixth (e.g., exponential growth/decay).
\end{itemize}
\end{prerequis}

\newpage

\coursec{Meaning and Classification of First‑Order Differential Equations}

Imagine a cocoa farmer in the South‑West region who wants to predict how fast his cocoa pods will ripen depending on temperature and humidity. The rate of change of the ripening process can be modeled by a first‑order differential equation. By learning to separate variables and integrate, you will be able to give the farmer a precise formula for the harvest time.

What is the farmer really asking? He is not asking \emph{how ripe} the pods are today; he is asking \emph{how fast} ripeness changes with time. In mathematics, a quantity whose \cle{rate of change} is known — but whose value is not — is described by an equation involving a derivative. Such an equation is called a \cle{differential equation}. This whole chapter is devoted to the simplest and most useful family of them: first‑order equations whose variables can be separated. But before solving anything, we must understand exactly what these objects are.

\courssub{What is a differential equation?}

\textbf{From a concrete situation to a definition.}\par
Suppose a biologist observes that a colony of bacteria in a laboratory in Yaoundé grows at a rate proportional to its size. If $N(t)$ denotes the number of bacteria at time $t$ (in hours), the observation translates into
\[ \frac{dN}{dt} = kN, \]
where $k$ is a positive constant. This is not an ordinary algebraic equation: the unknown is not a number but a \emph{function} $N(t)$, and the equation links the function to its own derivative. Solving it means finding \emph{all} functions whose derivative equals $k$ times themselves.

\begin{definition}[Differential equation]
A \cle{differential equation} is an equation that relates an unknown function, its independent variable(s), and one or more of its derivatives. A \emph{solution} of the equation is any function which, together with its derivatives, satisfies the equation identically on some interval.
\end{definition}

\begin{exemplebox}[Checking that a function is a solution]
Consider the equation $\dfrac{dy}{dx} = 2x$. The function $y = x^{2} + 5$ is a solution, because $\dfrac{dy}{dx} = 2x$ for every real $x$. But so are $y = x^{2} - 3$ and $y = x^{2} + \sqrt{2}$. In fact every function $y = x^{2} + C$, where $C$ is an arbitrary constant, is a solution. This simple observation is fundamental: a differential equation usually has \emph{infinitely many} solutions, indexed by constants of integration.
\end{exemplebox}

\courssub{Ordinary versus partial differential equations}

\begin{definition}[ODE and PDE]
A differential equation is called \cle{ordinary} (an \emph{ODE}) when the unknown function depends on a \emph{single} independent variable, so that only ordinary derivatives $\dfrac{dy}{dx}, \dfrac{d^{2}y}{dx^{2}}, \dots$ appear. It is called \cle{partial} (a \emph{PDE}) when the unknown function depends on \emph{several} independent variables and partial derivatives such as $\dfrac{\partial u}{\partial x}$ appear.
\end{definition}

For example, $\dfrac{dy}{dx} = x - y$ is an ODE (one variable $x$), whereas the heat equation $\dfrac{\partial u}{\partial t} = \alpha \dfrac{\partial^{2} u}{\partial x^{2}}$, which describes how temperature $u(x,t)$ spreads along a metal bar, is a PDE (two variables $x$ and $t$). \textbf{In this course we study only ordinary differential equations.}

\courssub{Order and degree of a differential equation}

Two numbers classify an ODE at a glance.

\begin{definition}[Order and degree]
The \cle{order} of a differential equation is the order of the \emph{highest derivative} appearing in it. The \cle{degree} of a differential equation is the \emph{power} to which the highest derivative is raised, \emph{after} the equation has been cleared of fractions and radicals involving the derivatives.
\end{definition}

\begin{exemplebox}[Finding order and degree]
\begin{enumerate}
\item $\dfrac{dy}{dx} + 3y = \sin x$: highest derivative is the first, raised to power $1$. Order $1$, degree $1$.
\item $\dfrac{d^{2}y}{dx^{2}} + 4\dfrac{dy}{dx} + y = 0$: order $2$, degree $1$.
\item $\left(\dfrac{dy}{dx}\right)^{3} + y = x$: order $1$, degree $3$.
\item $\sqrt{\dfrac{d^{2}y}{dx^{2}}} = y$: squaring gives $\dfrac{d^{2}y}{dx^{2}} = y^{2}$, so order $2$, degree $1$.
\item $\dfrac{d^{3}y}{dx^{3}} + \left(\dfrac{d^{2}y}{dx^{2}}\right)^{2} = x$: order $3$, degree $1$ (the highest derivative is $\frac{d^3y}{dx^3}$ to the first power).
\end{enumerate}
\end{exemplebox}

\begin{attention}[Order is not degree]
Do not confuse the \cle{order} (which derivative is the highest) with the \cle{degree} (the power of that highest derivative \emph{after} clearing fractions and roots). In $\left(\dfrac{d^{2}y}{dx^{2}}\right)^{5} + \dfrac{dy}{dx} = 0$, the order is $2$ and the degree is $5$ — not the other way round. Examiners regularly test this distinction with a one‑mark question.
\end{attention}

\begin{definition}[First‑order ordinary differential equation]
A \cle{first‑order ordinary differential equation} is an equation involving an independent variable $x$, an unknown function $y(x)$ and its first derivative $\dfrac{dy}{dx}$ (and no higher derivative). It can always be written in the form
\[ F\!\left(x,\, y,\, \frac{dy}{dx}\right) = 0, \]
and, when we can solve for the derivative, in the \emph{normal form}
\[ \frac{dy}{dx} = f(x, y). \]
\end{definition}

The normal form $\dfrac{dy}{dx} = f(x,y)$ is the one we shall use throughout this chapter. It says: \emph{the slope of the unknown curve at the point $(x,y)$ is a known number $f(x,y)$}.

\courssub{The concept of a solution: explicit and implicit}

\begin{propriete}[What it means to be a solution]
A function $\varphi$ is a solution of $\dfrac{dy}{dx} = f(x,y)$ on an interval $I$ if $\varphi$ is differentiable on $I$ and, for every $x \in I$,
\[ \varphi'(x) = f\bigl(x, \varphi(x)\bigr). \]
In words: substituting the function and its derivative into the equation produces a true statement \emph{at every point of the interval}. (This characterisation is used constantly in examinations; its statement is examinable.)
\end{propriete}

Solutions come in two presentations.

\begin{definition}[Explicit and implicit solutions]
A solution is \cle{explicit} when it is given in the form $y = \varphi(x)$, with $y$ isolated. It is \cle{implicit} when it is given by a relation $G(x, y) = C$ which defines $y$ as a function of $x$ without $y$ being isolated.
\end{definition}

\begin{exoresolu}[Verifying an implicit solution]
Show that the relation $x^{2} + y^{2} = C$ (with $C > 0$) defines implicit solutions of the differential equation $x + y\dfrac{dy}{dx} = 0$.
\tcblower
\textbf{Solution.} Differentiate both sides of $x^{2} + y^{2} = C$ with respect to $x$, remembering that $y$ is a function of $x$ (chain rule: $\dfrac{d}{dx}(y^{2}) = 2y\dfrac{dy}{dx}$):
\[ 2x + 2y\frac{dy}{dx} = 0 \quad\Longrightarrow\quad x + y\frac{dy}{dx} = 0. \]
This is exactly the given equation, and it holds wherever $y \neq 0$ (so that the circle has a well‑defined tangent slope). Hence $x^{2} + y^{2} = C$ is an implicit solution. Geometrically, the solutions are the circles centred at the origin: at each point of a circle, the tangent is perpendicular to the radius, which is precisely what $x + y\, y' = 0$ expresses.
\end{exoresolu}

\begin{attention}[A solution must be checked on an interval]
A formula is only a solution where it is differentiable and where the substitution makes sense. For instance $y = \sqrt{C - x^{2}}$ solves $x + yy' = 0$ only on $(-\sqrt{C}, \sqrt{C})$, not at the endpoints where the tangent is vertical. Always state the interval when asked to \emph{verify} a solution.
\end{attention}

\courssub{Where first‑order equations come from}

First‑order ODEs appear whenever a rate of change depends on the present state. Here are three classical sources, all of which we shall meet again in the applications section.

\textbf{Physics: Newton's law of cooling.}\par
A bottle of water taken from a fridge in Douala warms up at a rate proportional to the difference between its temperature $T(t)$ and the ambient temperature $T_{a}$:
\[ \frac{dT}{dt} = -k\,(T - T_{a}), \qquad k > 0. \]

\textbf{Biology: limited (logistic) growth.}\par
When a population $P(t)$ cannot exceed a carrying capacity $K$ (think of fish in a pond near Limbe), a refined model is
\[ \frac{dP}{dt} = rP\left(1 - \frac{P}{K}\right). \]

\textbf{Economics: continuous compound interest.}\par
If a savings account at a credit union earns interest continuously at rate $r$, the balance $M(t)$ satisfies
\[ \frac{dM}{dt} = rM. \]

Notice the common pattern: in each case \emph{nature gives us the slope}, and mathematics must recover the curve.

\courssub{Graphical interpretation: slope fields and integral curves}

Since $\dfrac{dy}{dx} = f(x,y)$ tells us the slope of the solution curve passing through any point $(x,y)$, we can \emph{see} all the solutions at once without solving anything: at each point of a grid, draw a short segment with slope $f(x,y)$. The resulting picture is called a \cle{slope field} (or direction field). A curve that follows the little segments everywhere is called an \cle{integral curve} — it is the graph of a solution.

\begin{methode}[Sketching a slope field by hand]
\begin{enumerate}
\item Choose a grid of points, for example $x, y \in \{-2, -1, 0, 1, 2\}$.
\item At each grid point $(x, y)$, compute the slope $m = f(x, y)$. Organise the values in a table of slopes (rows: values of $y$; columns: values of $x$).
\item At each point, draw a short line segment with the computed slope (horizontal if $m = 0$, steep if $|m|$ is large).
\item To draw an integral curve, start at the required point and sketch a smooth curve that stays tangent to the segments.
\end{enumerate}
\end{methode}

\begin{exemplebox}[Table of slopes for $\dfrac{dy}{dx} = x - y$]
The slope depends only on $x - y$, so it is constant along each line $x - y = c$ (such lines are called \emph{isoclines}).
\begin{center}
\begin{tabularx}{\linewidth}{|l|X|X|X|X|X|}
\hline
\rowcolor{popblueL} $y \backslash x$ & $-2$ & $-1$ & $0$ & $1$ & $2$ \\
\hline
$2$ & $-4$ & $-3$ & $-2$ & $-1$ & $0$ \\
\hline
$1$ & $-3$ & $-2$ & $-1$ & $0$ & $1$ \\
\hline
$0$ & $-2$ & $-1$ & $0$ & $1$ & $2$ \\
\hline
$-1$ & $-1$ & $0$ & $1$ & $2$ & $3$ \\
\hline
$-2$ & $0$ & $1$ & $2$ & $3$ & $4$ \\
\hline
\end{tabularx}
\end{center}
Along the line $y = x$ every slope is $0$: solutions cross this line horizontally. Below it slopes are positive, above it negative.
\end{exemplebox}

\begin{popfigure}
\begin{tikzpicture}[scale=0.85]
\draw[->, popmuted] (-3.4,0) -- (3.4,0) node[below left] {$x$};
\draw[->, popmuted] (0,-3.4) -- (0,3.4) node[below left] {$y$};
% slope field for y' = x - y
\foreach \x in {-3,-2.5,...,3}{
  \foreach \y in {-3,-2.5,...,3}{
    \pgfmathsetmacro{\m}{\x-\y}
    \pgfmathsetmacro{\dx}{0.16/sqrt(1+\m*\m)}
    \pgfmathsetmacro{\dy}{\m*0.16/sqrt(1+\m*\m)}
    \draw[popline] ({\x-\dx},{\y-\dy}) -- ({\x+\dx},{\y+\dy});
  }
}
% isocline y = x
\draw[popmuted, dashed] (-3,-3) -- (3,3) node[pos=0.9, above left, text=popmuted] {$y=x$};
% integral curves y = C e^{-x} + x - 1
\draw[popblue, thick, domain=-3:3, samples=60] plot (\x, {2*exp(-\x)+\x-1});
\draw[poporange, thick, domain=-3:3, samples=60] plot (\x, {0.5*exp(-\x)+\x-1});
\draw[popgreen, thick, domain=-3:3, samples=60] plot (\x, {\x-1});
\draw[poppurple, thick, domain=-3:2.2, samples=60] plot (\x, {-0.5*exp(-\x)+\x-1});
\node[text=popblue, align=center] at (-2.3,2.6) {$C=2$};
\node[text=popgreen, align=center] at (2.4,0.6) {$C=0$};
\end{tikzpicture}
\legende{Slope field of $\dfrac{dy}{dx} = x - y$ with four integral curves. Every segment has slope $x - y$; each curve stays tangent to the segments.}
\end{popfigure}

The equation $\dfrac{dy}{dx} = x - y$ can in fact be solved exactly (it is linear; you will meet such equations later). Its solutions are
\[ y = Ce^{-x} + x - 1, \qquad C \in \mathbb{R}. \]
Let us verify this claim directly, as the property above requires: if $y = Ce^{-x} + x - 1$, then $\dfrac{dy}{dx} = -Ce^{-x} + 1$, and
\[ x - y = x - \left(Ce^{-x} + x - 1\right) = -Ce^{-x} + 1 = \frac{dy}{dx}. \]
The equation is satisfied for every real $x$, so $y = Ce^{-x} + x - 1$ is indeed the \cle{general solution}: a family of solutions depending on one arbitrary constant $C$. Each value of $C$ gives one integral curve of the slope field above.

\begin{popfigure}
\begin{tikzpicture}
\begin{axis}[
  width=0.85\linewidth, height=7cm,
  axis lines=middle,
  xlabel={$x$}, ylabel={$y$},
  xmin=-3, xmax=3, ymin=-4, ymax=6,
  domain=-3:3, samples=80,
  legend style={at={(0.03,0.97)}, anchor=north west, font=\small},
  grid=both, grid style={popline!40},
]
\addplot[popblue, thick] {3*exp(-x)+x-1};
\addplot[poporange, thick] {1*exp(-x)+x-1};
\addplot[popgreen, thick] {x-1};
\addplot[poppurple, thick] {-1.5*exp(-x)+x-1};
\legend{$C=3$, $C=1$, $C=0$, $C=-1.5$}
\end{axis}
\end{tikzpicture}
\legende{Explicit solutions $y = Ce^{-x} + x - 1$ for several values of $C$. For large $x$ the term $Ce^{-x}$ dies out and every curve approaches the line $y = x - 1$.}
\end{popfigure}

\begin{aretenir}[Key points of this section]
\begin{itemize}
\item A differential equation links an unknown function to its derivatives; solving it means finding \emph{functions}.
\item The \cle{order} is the highest derivative present; the \cle{degree} is the power of that derivative after clearing fractions and radicals.
\item A first‑order ODE in normal form reads $\dfrac{dy}{dx} = f(x,y)$: the slope of the solution is prescribed at every point.
\item A solution may be given \emph{explicitly} ($y = \varphi(x)$) or \emph{implicitly} ($G(x,y) = C$); it must satisfy the equation on an interval where it is differentiable.
\item The general solution of a first‑order ODE contains \emph{one arbitrary constant}; geometrically it is a one‑parameter family of integral curves of the slope field.
\end{itemize}
\end{aretenir}

\begin{savaistu}[Why one constant?]
A first‑order equation involves one derivative, and undoing one derivative requires one integration — hence one constant of integration. Likewise, a second‑order equation (like those governing the motion of a pendulum or the vibration of a \emph{balafon} string) produces two constants. This is why the general solution of a first‑order ODE is always a \emph{family} of curves, and why one extra piece of information — an initial condition — will be needed to select a single curve, as we shall see in Section 4.
\end{savaistu}

\begin{exercice}[Classifying differential equations]
For each equation, give the order and the degree, and state whether it is ordinary or partial.
\begin{enumerate}
\item $\dfrac{dy}{dx} = \dfrac{x^{2}}{y}$
\item $\left(\dfrac{d^{2}y}{dx^{2}}\right)^{2} + \left(\dfrac{dy}{dx}\right)^{3} = y$
\item $\dfrac{\partial u}{\partial t} = 3\dfrac{\partial u}{\partial x}$
\item $1 + \left(\dfrac{dy}{dx}\right)^{2} = \dfrac{1}{y^{2}}$
\end{enumerate}
\end{exercice}

\begin{corrige}[Exercise 1]
\begin{enumerate}
\item Highest derivative: $\dfrac{dy}{dx}$, power $1$. Order $1$, degree $1$; ordinary.
\item Highest derivative: $\dfrac{d^{2}y}{dx^{2}}$, raised to the power $2$. Order $2$, degree $2$; ordinary.
\item Partial derivatives appear: it is a PDE, of order $1$ and degree $1$.
\item Highest derivative: $\dfrac{dy}{dx}$, power $2$ (no fractions in the derivatives to clear). Order $1$, degree $2$; ordinary.
\end{enumerate}
\end{corrige}

\begin{exercice}[Verifying a solution]
Show that $y = \dfrac{1}{1 - x}$ is a solution of $\dfrac{dy}{dx} = y^{2}$, and state the largest interval containing $x = 0$ on which this solution is valid.
\end{exercice}

\begin{corrige}[Exercise 2]
If $y = (1 - x)^{-1}$, then $\dfrac{dy}{dx} = (1 - x)^{-2} = y^{2}$, so the equation is satisfied. The function is differentiable everywhere except at $x = 1$ where it is undefined. The largest interval containing $0$ on which it is a solution is therefore $(-\infty, 1)$.
\end{corrige}

Now that we can recognise, classify and visualise first‑order differential equations, the next step is to identify those that can be solved by the simplest and most powerful technique of this chapter: \emph{separation of variables}. That is the object of the next section.

\coursec{Recognising Separable Equations}

In Section 1 we learned to recognise a first-order differential equation: an equation involving $y$, $x$ and $\dfrac{dy}{dx}$, and we classified such equations by order and degree. We also saw that there is no universal formula that solves every first-order equation — each \emph{type} of equation has its own method. The very first question a GCE candidate must therefore ask when facing a differential equation is: \textbf{which type is it?} In this section we study the simplest and most frequently examined type at Advanced Level: equations whose variables can be \cle{separated}. By the end of the section you will be able to decide, quickly and with justification, whether a given equation is separable, and to carry out the algebraic manipulation that puts it into separated form, ready for the integration of Section 3.

\courssub{The Idea of Separation: an Intuitive Approach}

Think of the equation $\dfrac{dy}{dx} = 2xy$. The right-hand side is a \emph{product}: a piece that depends only on $x$ (namely $2x$) times a piece that depends only on $y$ (namely $y$). Because the two variables are ``unmixed'', we can send everything involving $y$ to the left with $dy$, and everything involving $x$ to the right with $dx$:
\[ \frac{1}{y}\,dy = 2x\,dx. \]
Each side now contains \emph{one single variable}, and each side can be integrated on its own. This is the whole philosophy of separation of variables.

Contrast this with $\dfrac{dy}{dx} = x + y$. Here the right-hand side is a \emph{sum} in which $x$ and $y$ are glued together by the $+$ sign. No amount of division or factorisation will produce a product of an $x$-part and a $y$-part: the variables are genuinely mixed. This equation is \emph{not} separable (it will be solved later by other methods).

\begin{definition}[Separable differential equation]
A first-order differential equation is said to be \cle{separable} (or to have \cle{separable variables}) if it can be written in the form
\[ \frac{dy}{dx} = H(x)\,G(y), \]
that is, the right-hand side is a \textbf{product} of a function of $x$ alone and a function of $y$ alone. Equivalently, it can be rearranged into the \emph{separated form}
\[ g(y)\,dy = h(x)\,dx, \]
where $g$ depends only on $y$ and $h$ depends only on $x$.
\end{definition}

\begin{attention}[The criterion is a PRODUCT, not a sum]
Separation is possible precisely when $x$ and $y$ appear as \textbf{factors} on the right-hand side. Sums and differences such as $x + y$, $x^2 - y^2$ (which does not factor as a product of single-variable functions because $(x-y)(x+y)$ still mixes the variables), or $e^{x+y}$ handled carelessly, are the classic traps. Note that $e^{x+y} = e^x e^y$ \emph{is} a product, so $\dfrac{dy}{dx} = e^{x+y}$ \textbf{is} separable, whereas $\dfrac{dy}{dx} = x + y$ is not.
\end{attention}

\courssub{The Key Property Justifying Separation}

\begin{propriete}[Separation of variables]
If $\dfrac{dy}{dx} = f(x)\,g(y)$ and $g(y) \neq 0$ on the interval considered, then the equation is equivalent to
\[ \frac{1}{g(y)}\,dy = f(x)\,dx. \]
(The proof that integrating both sides yields the solution is examinable and will be given in Section 3; here we focus on the algebraic step.)
\end{propriete}

\textbf{Why division by $g(y)$ is legitimate.} On any interval where $g(y) \neq 0$, the function $\dfrac{1}{g(y)}$ is well defined, so we may multiply both sides of $\dfrac{dy}{dx} = f(x)g(y)$ by $\dfrac{dx}{g(y)}$. Formally, using the chain rule, $\dfrac{1}{g(y)}\dfrac{dy}{dx} = f(x)$ says that the antiderivative of $\dfrac{1}{g(y)}$ with respect to $y$ and the antiderivative of $f(x)$ with respect to $x$ have the same derivative with respect to $x$; hence they differ only by a constant. This is exactly what justifies integrating the two sides separately.

\courssub{Algebraic Manipulation to Achieve Separation}

Most examination equations are not handed to you already in product form: you must \emph{manufacture} the product. The three standard tools are \textbf{factorisation}, \textbf{division} and \textbf{multiplication}. Let us see them at work.

\begin{exemplebox}[Factorising to reveal the product]
Consider $\dfrac{dy}{dx} = xy + x$. At first sight the right-hand side is a sum. But factorising gives
\[ xy + x = x(y+1), \]
a product of a function of $x$ and a function of $y$. Hence the equation is separable, and for $y \neq -1$:
\[ \frac{dy}{y+1} = x\,dx. \]
\end{exemplebox}

\begin{exemplebox}[Dividing both sides]
Consider $x\dfrac{dy}{dx} = y(1 + x^2)$. Divide both sides by $x$ (for $x \neq 0$):
\[ \frac{dy}{dx} = \frac{y(1+x^2)}{x} = y\cdot\frac{1+x^2}{x}, \]
a product of a function of $y$ and a function of $x$. Separating (for $y \neq 0$):
\[ \frac{dy}{y} = \frac{1+x^2}{x}\,dx. \]
\end{exemplebox}

\begin{exemplebox}[Cross-multiplication from a quotient form]
Consider $\dfrac{dy}{dx} = \dfrac{x\,e^{x}}{y^2}$. The right-hand side is already $H(x)G(y)$ with $H(x) = xe^x$ and $G(y) = \dfrac{1}{y^2}$. Multiplying through by $y^2\,dx$ gives the separated form directly:
\[ y^2\,dy = x e^{x}\,dx. \]
\end{exemplebox}

\begin{methode}[Checklist to test separability]
To decide whether $\dfrac{dy}{dx} = F(x,y)$ is separable:
\begin{enumerate}
\item \textbf{Isolate} $\dfrac{dy}{dx}$ on one side if necessary (clear brackets, divide by the coefficient of $\dfrac{dy}{dx}$).
\item \textbf{Try to write} the other side as a product $H(x)\,G(y)$: factorise, use index laws ($e^{x+y} = e^x e^y$), or split quotients.
\item \textbf{If a product is obtained}, the equation is separable: divide by $G(y)$ and multiply by $dx$ to reach $g(y)\,dy = h(x)\,dx$.
\item \textbf{If no factorisation is possible} (typically because of a sum $x+y$ or a genuinely mixed expression such as $\sin(xy)$), the equation is \emph{not} separable — do not force it.
\item \textbf{Record any values excluded} by your divisions (e.g.\ $y = 0$): they may be lost solutions.
\end{enumerate}
\end{methode}

\begin{popfigure}
\begin{center}
\begin{tikzpicture}[
  font=\small,
  box/.style={draw=popblue, thick, rounded corners=2pt, fill=popblueL, text width=4.6cm, align=center, minimum height=1.1cm},
  dec/.style={draw=poporange, thick, diamond, aspect=2.2, fill=white, text width=3.4cm, align=center, inner sep=1pt},
  res/.style={draw=popgreen, thick, rounded corners=2pt, fill=popgreenL, text width=4.2cm, align=center, minimum height=1.1cm},
  arr/.style={-{Stealth}, thick, popdark}
]
\node[box] (start) at (0,0) {Write the equation as $\dfrac{dy}{dx} = F(x,y)$};
\node[dec] (q1) at (0,-2.6) {Can $F(x,y)$ be factorised as $H(x)\,G(y)$?};
\node[res] (yes) at (-4.6,-5.4) {\textbf{Separable:} write $\dfrac{dy}{G(y)} = H(x)\,dx$, then integrate};
\node[box] (no) at (4.6,-5.4) {\textbf{Not separable:} try another method (integrating factor, substitution, \dots)};
\draw[arr] (start) -- (q1);
\draw[arr] (q1) -| node[pos=0.25, above]{Yes} (yes);
\draw[arr] (q1) -| node[pos=0.25, above]{No} (no);
\end{tikzpicture}
\end{center}
\legende{Decision flowchart: testing whether a first-order equation is separable.}
\end{popfigure}

\courssub{Recognising Non-Separable Forms}

Knowing when separation \emph{fails} is as valuable as knowing when it works: it saves precious examination time.

\begin{exemplebox}[Equations that are NOT separable]
\begin{itemize}
\item $\dfrac{dy}{dx} = x + y$: a sum; no factorisation $H(x)G(y)$ exists.
\item $\dfrac{dy}{dx} = x^2 + y^2$: a sum of squares; not factorisable into single-variable factors over $\mathbb{R}$.
\item $\dfrac{dy}{dx} = \sin(x + y)$ or $\dfrac{dy}{dx} = e^{xy}$: the variables are mixed \emph{inside} another function.
\item $\dfrac{dy}{dx} = \dfrac{x + y}{x - y}$: this is a \emph{homogeneous} equation (right-hand side depends only on $\dfrac{y}{x}$); it is solved by the substitution $v = \dfrac{y}{x}$, \textbf{not} by direct separation.
\end{itemize}
\end{exemplebox}

\begin{attention}[Homogeneous does not mean separable]
A frequent confusion at GCE: an equation like $\dfrac{dy}{dx} = \dfrac{y}{x}$ \emph{is} both homogeneous and separable (it equals $y \cdot \dfrac{1}{x}$), but $\dfrac{dy}{dx} = \dfrac{x+y}{x-y}$ is homogeneous \emph{without} being separable. Test for the product form first; never assume.
\end{attention}

\begin{attention}[Lost solutions when dividing by $g(y)$]
When we divide by $g(y)$ we silently assume $g(y) \neq 0$. Any constant $y = c$ with $g(c) = 0$ is in fact a \textbf{solution} of the original equation (since then $\dfrac{dy}{dx} = 0 = f(x)\cdot 0$), but it is \emph{lost} in the division. Example: separating $\dfrac{dy}{dx} = 2xy$ by dividing by $y$ loses the solution $y = 0$. \textbf{Always check} whether the excluded constant solutions satisfy the equation and the initial condition.
\end{attention}

\courssub{Worked Example and Classification Practice}

\begin{exoresolu}[Separating a disguised equation]
Show that $xy\,\dfrac{dy}{dx} = y^2 + 1$ is separable, and write it in the form $g(y)\,dy = h(x)\,dx$.
\tcblower
\textbf{Solution.} Divide both sides by $x$ (valid for $x \neq 0$):
\[ \frac{dy}{dx} = \frac{y^2+1}{xy} = \frac{1}{x}\cdot\frac{y^2+1}{y}. \]
This is a product $H(x)G(y)$ with $H(x) = \dfrac{1}{x}$ and $G(y) = \dfrac{y^2+1}{y}$, so the equation \textbf{is separable}. Multiplying by $\dfrac{y}{y^2+1}\,dx$ (note $y^2 + 1 \neq 0$ for all real $y$, so no solution is lost):
\[ \frac{y}{y^2+1}\,dy = \frac{1}{x}\,dx. \]
This is the required separated form, ready for integration.
\end{exoresolu}

\begin{exercice}[Classify ten equations]
For each equation, state whether it is \textbf{separable} or \textbf{not separable}, and justify briefly. When separable, give the separated form.
\begin{enumerate}
\item $\dfrac{dy}{dx} = 3x^2 y$
\item $\dfrac{dy}{dx} = x + 2y$
\item $\dfrac{dy}{dx} = e^{x - y}$
\item $\dfrac{dy}{dx} = \dfrac{x}{y} + \dfrac{y}{x}$
\item $(1+x^2)\dfrac{dy}{dx} = 1 + y^2$
\item $\dfrac{dy}{dx} = \sin x \cos y$
\item $\dfrac{dy}{dx} = \sin(x y)$
\item $y\,\dfrac{dy}{dx} = x(y^2 - 4)$
\item $\dfrac{dy}{dx} = \dfrac{x^2 + x}{y}$
\item $\dfrac{dy}{dx} = (x + y)^2$
\end{enumerate}
\end{exercice}

\begin{corrige}[Exercise 1]
\begin{enumerate}
\item \textbf{Separable:} $3x^2 y = (3x^2)(y)$. For $y \neq 0$: $\dfrac{dy}{y} = 3x^2\,dx$. (Check: $y = 0$ is also a solution.)
\item \textbf{Not separable:} $x + 2y$ is a sum and cannot be written as $H(x)G(y)$.
\item \textbf{Separable:} $e^{x-y} = e^x e^{-y}$, a product. $e^{y}\,dy = e^{x}\,dx$.
\item \textbf{Not separable:} $\dfrac{x}{y} + \dfrac{y}{x} = \dfrac{x^2+y^2}{xy}$ cannot be factored into single-variable functions (it is homogeneous).
\item \textbf{Separable:} $\dfrac{dy}{dx} = \dfrac{1+y^2}{1+x^2} = \dfrac{1}{1+x^2}(1+y^2)$. Hence $\dfrac{dy}{1+y^2} = \dfrac{dx}{1+x^2}$.
\item \textbf{Separable:} product $\sin x \cdot \cos y$. For $\cos y \neq 0$: $\dfrac{dy}{\cos y} = \sin x\,dx$.
\item \textbf{Not separable:} the variables are mixed inside the sine.
\item \textbf{Separable:} $\dfrac{dy}{dx} = \dfrac{x(y^2-4)}{y} = x\cdot\dfrac{y^2-4}{y}$. Hence $\dfrac{y}{y^2-4}\,dy = x\,dx$ (for $y \neq \pm 2$; check $y = \pm 2$ separately).
\item \textbf{Separable:} $\dfrac{dy}{dx} = (x^2+x)\cdot\dfrac{1}{y}$. Hence $y\,dy = (x^2+x)\,dx$.
\item \textbf{Not separable:} $(x+y)^2 = x^2 + 2xy + y^2$ contains the mixed term $2xy$; no product form exists (a substitution $u = x+y$ is needed).
\end{enumerate}
\end{corrige}

\courssub{A Brief Link with Exact Equations}

\begin{savaistu}[Separable implies exact — but not conversely]
An equation written as $M(x)\,dx + N(y)\,dy = 0$ is called \cle{exact} when $M$ and $N$ are the partial derivatives of some potential function. Every separable equation $g(y)\,dy - h(x)\,dx = 0$ is automatically exact, because the ``cross-derivatives'' are both zero: $\dfrac{\partial}{\partial y}(-h(x)) = 0 = \dfrac{\partial}{\partial x}(g(y))$. The converse is false: for example $2xy\,dx + x^2\,dy = 0$ is exact (it is the differential of $x^2 y$) yet its separated structure is hidden. At GCE Advanced Level you only need the separable case, but it is reassuring to know that separation of variables sits inside a larger, elegant theory.
\end{savaistu}

\begin{aretenir}[Key points of this section]
\begin{itemize}
\item An equation is \cle{separable} if it can be written $\dfrac{dy}{dx} = H(x)G(y)$, equivalently $g(y)\,dy = h(x)\,dx$.
\item The test is: can the right-hand side be \textbf{factored} into an $x$-part times a $y$-part? Use factorisation, division, multiplication and index laws.
\item Sums ($x + y$), mixed arguments ($\sin(xy)$, $e^{xy}$) and most homogeneous equations are \textbf{not} separable.
\item Dividing by $g(y)$ assumes $g(y) \neq 0$: always check the constant solutions $g(y) = 0$ that may have been lost.
\item Once separated, each side contains a single variable — the stage is set for integration, which is the object of the next section.
\end{itemize}
\end{aretenir}

\coursec{Separation of Variables and Integration to Obtain General Solutions}

In Section 2 we learned to \emph{recognise} a separable equation: one that can be written in the form
\[
\frac{dy}{dx} = f(x)\,g(y).
\]
Recognition, however, is only the passport. We now learn the journey itself: how to \emph{solve} such an equation completely. The strategy is beautifully simple — put all the $y$'s with $dy$ on one side, all the $x$'s with $dx$ on the other, and integrate. The skill lies in doing this \emph{correctly}: handling the constant of integration, the absolute values that appear in logarithms, and the final step of making $y$ the subject when the question demands it. This section is the computational heart of the chapter, and it is where most GCE marks on differential equations are won or lost.

\courssub{The Principle of Separation}

Consider the separable equation $\dfrac{dy}{dx} = f(x)\,g(y)$, with $g(y) \neq 0$ on the interval of interest. Dividing both sides by $g(y)$ gives
\[
\frac{1}{g(y)}\,\frac{dy}{dx} = f(x).
\]
Now integrate both sides with respect to $x$:
\[
\int \frac{1}{g(y)}\,\frac{dy}{dx}\,dx = \int f(x)\,dx.
\]
By the rule of integration by substitution (with $y$ as the new variable), the left-hand side is exactly $\displaystyle\int \frac{1}{g(y)}\,dy$. Hence the equation is equivalent to
\[
\boxed{\ \int \frac{1}{g(y)}\,dy = \int f(x)\,dx + C\ }
\]
This is the \cle{formal separation} of variables. Informally, students often remember it as "multiply through by $dx$ and divide by $g(y)$":
\[
\frac{dy}{g(y)} = f(x)\,dx,
\]
then attach integral signs. Although $dy$ and $dx$ are not ordinary numbers, this shorthand is justified rigorously by the substitution rule above, so you may use it confidently in the examination.

\begin{propriete}[One Constant of Integration Only]
When both sides of a separated equation are integrated, \textbf{only one} arbitrary constant is needed. If the left integral gives $G(y) + C_1$ and the right gives $F(x) + C_2$, then
\[
G(y) + C_1 = F(x) + C_2 \quad\Longrightarrow\quad G(y) = F(x) + (C_2 - C_1).
\]
Since $C_2 - C_1$ is itself just an arbitrary constant, we rename it $C$ and write $G(y) = F(x) + C$. \emph{Writing two constants is not wrong, but it is untidy and wastes time; examiners expect a single constant $C$.}
\end{propriete}

\begin{methode}[The Separation–Integration Algorithm]
To solve a separable first-order differential equation $\dfrac{dy}{dx} = f(x)\,g(y)$:
\begin{enumerate}
\item \textbf{Separate.} Rearrange so that all terms in $y$ multiply $dy$ and all terms in $x$ multiply $dx$: $\dfrac{dy}{g(y)} = f(x)\,dx$.
\item \textbf{Integrate both sides.} Write $\displaystyle\int \frac{dy}{g(y)} = \int f(x)\,dx$ and evaluate each integral using the appropriate technique (power rule, exponential, logarithmic, trigonometric, or partial fractions).
\item \textbf{Insert one constant.} Add a single arbitrary constant $C$ on one side only.
\item \textbf{Simplify.} Handle absolute values (e.g.\ $\ln|y|$), combine logarithms, and exponentiate where useful.
\item \textbf{Make $y$ the subject if required.} Give the general solution explicitly as $y = \varphi(x)$ when possible; otherwise leave it in implicit form.
\end{enumerate}
\end{methode}

\courssub{Integration Techniques You Must Command}

Separation reduces every problem in this chapter to two ordinary integrals. The GCE examiner will therefore test your integration fluency. The techniques required are:

\begin{center}
\begin{tabularx}{\linewidth}{|l|X|X|}
\hline
\rowcolor{popblueL} \textbf{Technique} & \textbf{Typical occurrence} & \textbf{Result} \\
\hline
Power rule & $\int x^n\,dx$, $\int y^n\,dy$ & $\dfrac{x^{n+1}}{n+1}$ \ ($n \neq -1$) \\
\hline
Logarithmic & $\int \dfrac{dy}{y}$, $\int \dfrac{f'(x)}{f(x)}\,dx$ & $\ln|y|$, $\ln|f(x)|$ \\
\hline
Exponential & $\int e^{kx}\,dx$ & $\dfrac{1}{k}e^{kx}$ \\
\hline
Trigonometric & $\int \sin x\,dx$, $\int \sec^2 y\,dy$ & $-\cos x$, $\tan y$ \\
\hline
Partial fractions & $\int \dfrac{dy}{y(a-y)}$ & split, then logarithms \\
\hline
\end{tabularx}
\end{center}

The last row deserves special attention: equations such as $\dfrac{dy}{dx} = ky(a-y)$ (the logistic-type equations that appear in population models) separate into an integral that \emph{only} partial fractions can unlock. We shall meet this in Section 5.

\courssub{The Absolute Value in $\ln|y|$: a Non-Negotiable Detail}

When separation produces $\displaystyle\int \frac{dy}{y}$, the answer is $\ln|y|$, \emph{not} $\ln y$. The reason is fundamental: the logarithm $\ln y$ is only defined for $y > 0$, whereas the solutions of a differential equation may well be negative. The antiderivative of $1/y$ that is valid on \emph{every} interval not containing $0$ is $\ln|y|$.

Let us see what happens when we integrate carefully. Suppose separation leads to
\[
\ln|y| = x + C.
\]
Exponentiating: $|y| = e^{x+C} = e^{C}e^{x}$. Writing $A = e^{C}$ (so $A > 0$), we get $|y| = Ae^{x}$, that is,
\[
y = \pm A e^{x}.
\]
Since $A$ is an arbitrary \emph{positive} constant, $\pm A$ is an arbitrary \emph{non-zero} constant. Renaming it $K$, the general solution is $y = Ke^{x}$ with $K \in \mathbb{R}\setminus\{0\}$ (and $y = 0$ can usually be absorbed by allowing $K = 0$). The negative branch $y = -Ae^{x}$ would have been \textbf{lost} had we written $\ln y$ instead of $\ln|y|$.

\begin{attention}[Forgetting the Absolute Value]
Writing $\ln y$ instead of $\ln|y|$ silently discards all negative solutions. In an examination, if the initial condition later gives $y < 0$, a candidate who dropped the absolute value will find no valid constant and lose the marks. \textbf{Rule:} every time you integrate $1/y$ (or $1/(ay+b)$), write the modulus signs immediately, and remove them only when exponentiating, absorbing the $\pm$ into the arbitrary constant.
\end{attention}

\courssub{Explicit versus Implicit General Solutions}

After integration we obtain a relation $G(y) = F(x) + C$. Two situations arise.

\textbf{Case 1: we can solve for $y$.} If $G$ is a simple expression (a power, a logarithm, an exponential), we isolate $y$ and present the \cle{general solution} explicitly as $y = \varphi(x, C)$. This is the preferred final form when the question says "express $y$ in terms of $x$".

\textbf{Case 2: we cannot (or need not) solve for $y$.} Some integrations yield relations such as
\[
\frac{y^2}{2} + \ln|y| = \frac{x^3}{3} + C,
\]
which cannot be solved for $y$ in closed form. The relation itself \emph{is} the general solution, given in \cle{implicit form}. This is perfectly acceptable — indeed, attempting to force an explicit form where none exists in elementary functions is a waste of examination time.

\begin{aretenir}[Tip: Read the Question's Final Instruction]
When the integral yields an implicit relation, solve for $y$ \textbf{only if the question requires it} ("give $y$ in terms of $x$", "show that $y = \ldots$"). Otherwise, a clean implicit relation with one arbitrary constant earns full marks. Never leave the answer as an unevaluated integral.
\end{aretenir}

\courssub{A Complete Worked Example, Step by Step}

\begin{exemplebox}[Solving $\dfrac{dy}{dx} = \dfrac{2x}{y}$]
\textbf{Step 1 — Separate.} Multiply both sides by $y$ and by $dx$:
\[
y\,dy = 2x\,dx.
\]
\textbf{Step 2 — Integrate.}
\[
\int y\,dy = \int 2x\,dx.
\]
\textbf{Step 3 — One constant.}
\[
\frac{y^2}{2} = x^2 + C_1.
\]
\textbf{Step 4 — Simplify.} Multiply by 2 and rename $2C_1$ as $C$:
\[
y^2 = 2x^2 + C.
\]
\textbf{Step 5 — Explicit form (optional).} $y = \pm\sqrt{2x^2 + C}$. The implicit form $y^2 = 2x^2 + C$ describes a \emph{family of hyperbolas} (and the pair of straight lines $y = \pm\sqrt{2}\,x$ when $C = 0$); each value of $C$ gives one curve of the family.
\end{exemplebox}

The figure below shows several members of this family. Notice the two branches (upper and lower) for each value of $C$: this is exactly the $\pm$ that the absolute-value discussion protects.

\begin{popfigure}
\begin{center}
\begin{tikzpicture}
\begin{axis}[
  width=11cm, height=8.5cm,
  axis lines=middle,
  xlabel={$x$}, ylabel={$y$},
  xmin=-3.2, xmax=3.2, ymin=-4.2, ymax=4.2,
  xtick={-2,-1,1,2}, ytick={-3,-2,-1,1,2,3},
  samples=120, domain=-3:3,
  legend style={at={(0.02,0.98)}, anchor=north west, font=\small}
]
\addplot[popblue, thick] {sqrt(max(0,2*x^2+2))};
\addplot[popblue, thick, forget plot] {-sqrt(max(0,2*x^2+2))};
\addlegendentry{$C=2$}
\addplot[poporange, thick] {sqrt(max(0,2*x^2+0.5))};
\addplot[poporange, thick, forget plot] {-sqrt(max(0,2*x^2+0.5))};
\addlegendentry{$C=0.5$}
\addplot[popgreen, thick] {sqrt(max(0,2))*x};
\addplot[popgreen, thick, forget plot] {-sqrt(max(0,2))*x};
\addlegendentry{$C=0$}
\addplot[poppurple, thick, domain=1.02:3] {sqrt(max(0,2*x^2-2))};
\addplot[poppurple, thick, domain=-3:-1.02, forget plot] {sqrt(max(0,2*x^2-2))};
\addplot[poppurple, thick, domain=1.02:3, forget plot] {-sqrt(max(0,2*x^2-2))};
\addplot[poppurple, thick, domain=-3:-1.02, forget plot] {-sqrt(max(0,2*x^2-2))};
\addlegendentry{$C=-2$}
\end{axis}
\end{tikzpicture}
\end{center}
\legende{Family of solution curves $y^2 = 2x^2 + C$ of $\dfrac{dy}{dx} = \dfrac{2x}{y}$ for several values of $C$. Each value of the constant selects one curve; the general solution is the whole family. For $C = -2$ the hyperbola exists only where $2x^2 - 2 \geq 0$, i.e.\ $|x| \geq 1$.}
\end{popfigure}

\courssub{Verification by Differentiation}

A powerful habit — and one that costs thirty seconds — is to \emph{differentiate your answer} and check that it satisfies the original equation. This is done by \cle{implicit differentiation} when the solution is implicit.

For the example above, differentiate $y^2 = 2x^2 + C$ with respect to $x$:
\[
2y\,\frac{dy}{dx} = 4x \quad\Longrightarrow\quad \frac{dy}{dx} = \frac{2x}{y},
\]
which is exactly the given equation. The arbitrary constant $C$ disappears on differentiation, confirming that the whole family satisfies the equation. \textbf{Always perform this check} before moving on in an examination: it catches sign errors, lost factors of 2, and forgotten constants instantly.

\begin{exoresolu}[Worked Exercise: an Equation Leading to a Logarithm]
Solve the differential equation $\dfrac{dy}{dx} = \dfrac{x}{x^2 + 1}\,y$, giving $y$ explicitly in terms of $x$.
\tcblower
\textbf{Solution.} The equation is separable with $f(x) = \dfrac{x}{x^2+1}$ and $g(y) = y$.

\textbf{Step 1 — Separate} (for $y \neq 0$):
\[
\frac{dy}{y} = \frac{x}{x^2+1}\,dx.
\]
\textbf{Step 2 — Integrate.} On the right, the numerator is (half of) the derivative of the denominator, so we use $\displaystyle\int \frac{f'(x)}{f(x)}\,dx = \ln|f(x)|$:
\[
\int \frac{dy}{y} = \frac{1}{2}\int \frac{2x}{x^2+1}\,dx
\quad\Longrightarrow\quad
\ln|y| = \frac{1}{2}\ln(x^2+1) + C,
\]
where the modulus is unnecessary on $x^2 + 1$ since it is always positive.

\textbf{Step 3 — Exponentiate and handle the modulus.}
\[
|y| = e^{C}\,(x^2+1)^{1/2}.
\]
Writing $K = \pm e^{C}$ (an arbitrary non-zero constant), and noting that $y = 0$ is recovered with $K = 0$:
\[
\boxed{\,y = K\sqrt{x^2 + 1}, \qquad K \in \mathbb{R}.\,}
\]
\textbf{Verification.} Differentiate: $\dfrac{dy}{dx} = K\cdot\dfrac{x}{\sqrt{x^2+1}} = \dfrac{x}{x^2+1}\cdot K\sqrt{x^2+1} = \dfrac{x}{x^2+1}\,y.$ \checkmark
\end{exoresolu}

\begin{exoresolu}[Worked Exercise: Trigonometric Separation, Implicit Answer]
Find the general solution of $\dfrac{dy}{dx} = \dfrac{\cos x}{\sin y}$, where $\sin y \neq 0$.
\tcblower
\textbf{Solution.} \textbf{Step 1 — Separate:} $\sin y\,dy = \cos x\,dx$.

\textbf{Step 2 — Integrate:}
\[
\int \sin y\,dy = \int \cos x\,dx
\quad\Longrightarrow\quad
-\cos y = \sin x + C_1.
\]
\textbf{Step 3 — Tidy:} multiplying by $-1$ and renaming $-C_1$ as $C$:
\[
\cos y + \sin x = C.
\]
Here it is not natural to write $y$ explicitly (one could write $y = \arccos(C - \sin x)$, but this introduces sign and branch complications). The \textbf{implicit form} $\cos y + \sin x = C$ is the expected general solution.

\textbf{Verification.} Differentiate implicitly: $-\sin y\,\dfrac{dy}{dx} + \cos x = 0$, so $\dfrac{dy}{dx} = \dfrac{\cos x}{\sin y}$. \checkmark
\end{exoresolu}

\courssub{Common Mistakes and How to Avoid Them}

\begin{attention}[Traps in the Separation–Integration Process]
\begin{itemize}
\item \textbf{Separating incorrectly.} In $\dfrac{dy}{dx} = x + y$, you \emph{cannot} separate — this equation is not separable (see Section 2). Never invent a separation that does not exist.
\item \textbf{Two constants.} Writing $\frac{y^2}{2} + C_1 = x^2 + C_2$ and leaving both: merge them into one constant immediately.
\item \textbf{Dropping the modulus.} $\int \frac{dy}{y} = \ln|y|$, always — until you exponentiate.
\item \textbf{Losing solutions when dividing.} Dividing by $g(y)$ assumes $g(y) \neq 0$; check separately whether $g(y) = 0$ gives a (constant) solution, e.g.\ $y = 0$ above.
\item \textbf{Stopping at the integral.} $\int \frac{dy}{y} = \int 2x\,dx$ is \emph{not} an answer; evaluate both integrals.
\item \textbf{Forgetting to verify.} One line of implicit differentiation protects all your marks.
\end{itemize}
\end{attention}

\begin{exercice}[Practice]
\begin{enumerate}
\item Find the general solution of $\dfrac{dy}{dx} = \dfrac{3x^2}{2y}$, expressing $y$ explicitly in terms of $x$.
\item Find the general solution of $\dfrac{dy}{dx} = e^{x}\,y^2$, leaving your answer in the form $y = \varphi(x)$. What restriction must be placed on the constant?
\item Find the general solution of $\dfrac{dy}{dx} = \dfrac{1 + y^2}{1 + x^2}$ in implicit form. (Hint: $\displaystyle\int\frac{du}{1+u^2} = \arctan u$.)
\end{enumerate}
\end{exercice}

\begin{corrige}[Exercise 1]
\textbf{1.} Separate: $2y\,dy = 3x^2\,dx$. Integrate: $y^2 = x^3 + C$. Hence $y = \pm\sqrt{x^3 + C}$.

\textbf{2.} Separate ($y \neq 0$): $\dfrac{dy}{y^2} = e^{x}\,dx$. Integrate: $-\dfrac{1}{y} = e^{x} + C$, so $y = -\dfrac{1}{e^{x} + C}$. We need $e^x + C \neq 0$ on the domain; also $y = 0$ is a separate constant solution lost when dividing by $y^2$.

\textbf{3.} Separate: $\dfrac{dy}{1+y^2} = \dfrac{dx}{1+x^2}$. Integrate: $\arctan y = \arctan x + C$. (Explicitly, $y = \tan(\arctan x + C)$, but the implicit form is the cleaner answer.)
\end{corrige}

\begin{aretenir}[Key Points of Section 3]
\begin{itemize}
\item Separation gives $\displaystyle\int \frac{dy}{g(y)} = \int f(x)\,dx + C$ — justified by integration by substitution.
\item \textbf{One} constant of integration only, placed on one side.
\item $\displaystyle\int \frac{dy}{y} = \ln|y|$: the modulus protects the negative branches; absorb $\pm$ into the constant when exponentiating.
\item Give $y$ explicitly if the question asks; otherwise a clean implicit relation is a complete answer.
\item \textbf{Always verify} by differentiating your solution and recovering the original equation.
\end{itemize}
\end{aretenir}

We can now produce the \emph{general} solution of any separable equation: an entire family of curves indexed by $C$. But real problems — a cooling cup of tea in Bamenda, a rumour spreading through a school — pick out \emph{one} curve from the family. Selecting that curve from a given starting value is the business of Section 4.

\coursec{Particular Solutions Using Initial Conditions}

In the previous section we learnt how to separate the variables of a first order differential equation and integrate both sides to obtain the \cle{general solution}. That general solution always contains an arbitrary constant $C$: it is not one curve but a whole \emph{family} of curves. In real problems, however, we usually know one extra piece of information — for instance the population of a town \emph{at the moment we start counting}, or the temperature of a pot of soup \emph{when it is taken off the fire}. This extra information, called an \cle{initial condition}, allows us to pick out exactly one curve from the family. That single curve is the \cle{particular solution}. This section teaches you how to find it, how to choose between several candidate constants, and how to state the domain on which your answer is valid — all essential skills for Paper 2 of the GCE Advanced Level.

\courssub{From a Family of Curves to a Single Curve: the Initial Condition}

\begin{definition}[Initial condition, initial value problem, particular solution]
Let $\dfrac{dy}{dx}=f(x)\,g(y)$ be a separable first order differential equation with general solution $F(x,y)=C$ (or $y=\varphi(x,C)$).
\begin{itemize}
\item An \cle{initial condition} is a pair of values $(x_0,y_0)$ such that the required solution satisfies $y(x_0)=y_0$. We often write ``$y=y_0$ when $x=x_0$''.
\item A differential equation together with an initial condition is called an \cle{initial value problem} (IVP).
\item The \cle{particular solution} of the IVP is the unique member of the family of general solutions whose graph passes through the point $(x_0,y_0)$. It is obtained by substituting $x=x_0$, $y=y_0$ into the general solution and solving for $C$.
\end{itemize}
\end{definition}

\textbf{Intuition.}\par
Think of the general solution as a bus timetable listing \emph{all} buses from Douala to Bamenda (infinitely many parallel curves). The initial condition tells you \emph{which} bus you actually boarded: the one passing through your town at your time. Once $C$ is fixed, there is no more freedom — the future and the past of the quantity are completely determined.

\begin{methode}[Finding the particular solution]
\begin{enumerate}
\item Solve the equation by separation of variables and write the \textbf{general solution}, keeping the constant $C$ visible.
\item \textbf{Substitute} the initial condition: replace $x$ by $x_0$ and $y$ by $y_0$.
\item \textbf{Solve for $C$} (or for the related constant $A=e^{C}$). If several values of the constant appear (for example $\pm$ from a square root), keep all candidates for the moment.
\item \textbf{Select the branch} compatible with the initial condition (sign, positivity, etc.).
\item \textbf{State the final answer} explicitly, in the form $y=\ldots$ (or as an implicit relation if $y$ cannot be isolated), and \textbf{mention the domain}: an interval containing $x_0$ on which the solution is defined and satisfies the equation.
\end{enumerate}
\end{methode}

\begin{propriete}[Existence and uniqueness for separable equations]
Consider the initial value problem $\dfrac{dy}{dx}=f(x)\,g(y)$, $y(x_0)=y_0$. If $f$ is continuous on an open interval containing $x_0$, and $g$ is continuous with $g(y_0)\neq 0$ on an open interval containing $y_0$, then the IVP has \textbf{exactly one} solution, defined on some interval around $x_0$. Geometrically: through the point $(x_0,y_0)$ passes one and only one solution curve, and two distinct solution curves never cross.
\emph{(Statement only — the proof is not examinable at this level.)}
\end{propriete}

\begin{attention}[When $g(y_0)=0$: the lost constant solutions]
If $g(y_0)=0$, then $y=y_0$ is a \textbf{constant solution} of the equation (both sides of $\frac{dy}{dx}=f(x)g(y)$ vanish). Recall that constant solutions are often \emph{lost} during separation because we divide by $g(y)$. Therefore: \textbf{if the initial condition lies on a constant solution, that constant solution \emph{is} the answer.} For example, the IVP $\frac{dy}{dx}=ky$, $y(0)=0$ has the solution $y=0$ for all $x$ — do not try to force $0=Ae^{kx}$ (which would require $A=0$, giving the same answer, but for equations like $\frac{dy}{dx}=y^{2/3}$ the constant solution can be missed entirely). Always check the constant solutions \emph{before} separating.
\end{attention}

\courssub{Worked Examples: Substituting the Initial Condition}

\begin{exoresolu}[A first initial value problem]
Find the particular solution of $\dfrac{dy}{dx}=\dfrac{x}{y}$ satisfying $y(0)=2$.
\tcblower
\textbf{Solution.} Separating variables: $y\,\dfrac{dy}{dx}=x$, so
\[ \int y\,dy=\int x\,dx \quad\Longrightarrow\quad \frac{y^{2}}{2}=\frac{x^{2}}{2}+C, \]
i.e. $y^{2}=x^{2}+2C$. Writing $A=2C$, the general solution is $y^{2}=x^{2}+A$.

Substitute $x=0$, $y=2$: $\;4=0+A$, so $A=4$.

Hence $y^{2}=x^{2}+4$, and taking the square root gives $y=\pm\sqrt{x^{2}+4}$. Since $y(0)=2>0$, we select the \textbf{positive branch}:
\[ \boxed{\,y=\sqrt{x^{2}+4}\,},\qquad x\in\mathbb{R}. \]
Check: $y(0)=\sqrt{4}=2$ \checkmark\ and $\frac{dy}{dx}=\frac{x}{\sqrt{x^{2}+4}}=\frac{x}{y}$ \checkmark.
\end{exoresolu}

\begin{exoresolu}[Population growth with given initial population]
A bacterial culture grows according to $\dfrac{dP}{dt}=0.3P$, where $P$ is the population (in thousands) and $t$ the time in hours. Given that the initial population is $P(0)=5$ (i.e. $5000$ bacteria), find $P(t)$ and the population after $10$ hours.
\tcblower
\textbf{Solution.} Note $P=0$ is a constant solution, but $P(0)=5\neq 0$, so we separate:
\[ \int\frac{dP}{P}=\int 0.3\,dt \quad\Longrightarrow\quad \ln|P|=0.3t+C. \]
Since $P>0$, $P=e^{0.3t+C}=Ae^{0.3t}$ with $A=e^{C}>0$.

Substitute $t=0$, $P=5$: $\;5=Ae^{0}=A$. Hence
\[ \boxed{\,P(t)=5e^{0.3t}\,},\qquad t\geq 0. \]
After $10$ hours: $P(10)=5e^{3}\approx 5\times 20.0855\approx 100.4$ thousand, i.e. about $100\,400$ bacteria.
\end{exoresolu}

\begin{exoresolu}[Newton's law of cooling with initial temperature]
A pot of water is removed from the fire at temperature $95^{\circ}$C and placed in a kitchen at $25^{\circ}$C. By Newton's law of cooling, $\dfrac{dT}{dt}=-k(T-25)$, where $T$ is the temperature after $t$ minutes. Given $k=0.05$, find $T(t)$ and the temperature after $20$ minutes.
\tcblower
\textbf{Solution.} The constant solution $T=25$ does not satisfy $T(0)=95$, so we separate:
\[ \int\frac{dT}{T-25}=\int -0.05\,dt \quad\Longrightarrow\quad \ln|T-25|=-0.05t+C. \]
Since $T>25$, $T-25=Ae^{-0.05t}$, i.e. $T(t)=25+Ae^{-0.05t}$.

Substitute $t=0$, $T=95$: $\;95=25+A$, so $A=70$. Hence
\[ \boxed{\,T(t)=25+70e^{-0.05t}\,},\qquad t\geq 0. \]
After $20$ minutes: $T(20)=25+70e^{-1}\approx 25+70(0.3679)\approx 50.8^{\circ}$C.
\end{exoresolu}

\courssub{Choosing the Correct Branch and Stating the Domain}

When the integration produces $\ln|y|$, a square, or a square root, the constant may not be uniquely determined by algebra alone. The initial condition is the referee that decides between the candidates.

\textbf{The $\pm$ ambiguity.}\par
If the general solution is $y^{2}=G(x)+A$, then $y=\pm\sqrt{G(x)+A}$. The sign is chosen so that $y(x_0)=y_0$: positive branch if $y_0>0$, negative branch if $y_0<0$. The two branches are \emph{different} solutions and never meet (uniqueness theorem).

\textbf{The sign of $A=e^{C}$.}\par
From $\ln|y|=h(x)+C$ we get $|y|=e^{C}e^{h(x)}$, so $y=Ae^{h(x)}$ where $A$ can be \emph{any nonzero real} (positive or negative), and $A=0$ recovers the constant solution $y=0$. The initial condition fixes both the value and the sign of $A$.

\textbf{Domain considerations.}\par
A particular solution is only valid on an interval that
\begin{itemize}
\item \textbf{contains $x_0$} (the solution must pass through the given point);
\item avoids values where the formula breaks down (division by zero, logarithm of a non-positive number, square root of a negative number);
\item is as large as possible subject to the above (the \emph{maximal interval} of the solution).
\end{itemize}

\begin{exoresolu}[Branch selection and maximal domain]
Solve the IVP $\dfrac{dy}{dx}=\dfrac{1}{2y}$, $y(1)=-3$.
\tcblower
\textbf{Solution.} Note $y=0$ is excluded (division by $2y$), and $y(1)=-3\neq 0$. Separating:
\[ \int 2y\,dy=\int dx \quad\Longrightarrow\quad y^{2}=x+C. \]
Substitute $x=1$, $y=-3$: $\;9=1+C$, so $C=8$, giving $y^{2}=x+8$.

Since $y(1)=-3<0$, take the \textbf{negative branch}: $y=-\sqrt{x+8}$.

Domain: we need $x+8>0$ (strictly, since $y\neq 0$), i.e. $x>-8$, and this interval contains $x_0=1$. Hence
\[ \boxed{\,y=-\sqrt{x+8},\qquad x\in\,]{-8},\,+\infty[\,}. \]
Check: $y(1)=-\sqrt{9}=-3$ \checkmark, $\frac{dy}{dx}=-\frac{1}{2\sqrt{x+8}}=\frac{1}{2y}$ \checkmark.
\end{exoresolu}

\begin{exoresolu}[Implicit solution and domain restriction]
Solve the IVP $\dfrac{dy}{dx}=\dfrac{2x}{y}$, $y(0)=-4$.
\tcblower
\textbf{Solution.} $y=0$ is excluded; $y(0)=-4\neq 0$. Separate: $y\,dy=2x\,dx$, integrate: $\frac{y^2}{2}=x^2+C$, so $y^2=2x^2+A$.
Substitute $x=0$, $y=-4$: $16=A$. Thus $y^2=2x^2+16$.
Since $y(0)=-4<0$, we take $y=-\sqrt{2x^2+16}$.
The expression under the square root is always positive, so the maximal domain is $\mathbb{R}$.
\[ \boxed{\,y=-\sqrt{2x^2+16},\qquad x\in\mathbb{R}\,}. \]
\end{exoresolu}

\begin{attention}[Common mistakes with initial conditions]
\begin{itemize}
\item \textbf{Forgetting to substitute before simplifying.} Substitute $(x_0,y_0)$ as soon as the general solution is found; do not carry $C$ through extra manipulations.
\item \textbf{Taking the wrong branch.} $y=\pm\sqrt{\ldots}$: the sign is dictated by the \emph{sign of $y_0$}, not by convenience.
\item \textbf{Writing $C$ instead of a number.} The final answer must contain \emph{no} arbitrary constant.
\item \textbf{Ignoring the domain.} A formula like $y=-\sqrt{x+8}$ stated for all $x$ is wrong for $x<-8$.
\item \textbf{Missing the constant solution.} If $y_0$ makes $g(y_0)=0$, the answer is the constant function $y=y_0$; separation of variables is unnecessary.
\item \textbf{Forgetting the absolute value in $\ln|y|$.} This leads to missing the negative branch when $y_0<0$.
\end{itemize}
\end{attention}

\begin{popfigure}
\begin{center}
\begin{tikzpicture}
\begin{axis}[
  width=12cm, height=8cm,
  axis lines=middle,
  xlabel={$x$}, ylabel={$y$},
  xmin=-2.2, xmax=2.2, ymin=-0.5, ymax=8.5,
  xtick={-2,-1,0,1,2}, ytick={2,4,6,8},
  domain=-2.2:2.2, samples=80,
  legend style={at={(0.03,0.97)}, anchor=north west, font=\small},
  grid=major, grid style={popline}
]
\addplot[popmuted, thin] {0.5*exp(0.8*x)}; \addlegendentry{$A=0.5$}
\addplot[popmuted, thin] {1.0*exp(0.8*x)}; \addlegendentry{$A=1$}
\addplot[popblue, very thick] {2.0*exp(0.8*x)}; \addlegendentry{$A=2$ (particular)}
\addplot[popmuted, thin] {4.0*exp(0.8*x)}; \addlegendentry{$A=4$}
\addplot[only marks, mark=*, poporange, mark size=2.5pt] coordinates {(0,2)};
\node[poporange, anchor=west, font=\small] at (axis cs:0.08,2.0) {$(0,P_0)=(0,2)$};
\end{axis}
\end{tikzpicture}
\end{center}
\legende{Family of solutions $y=Ae^{0.8x}$ of $\frac{dy}{dx}=0.8y$. The initial condition $y(0)=P_0=2$ selects the single curve $y=2e^{0.8x}$ (thick blue) passing through $(0,2)$.}
\end{popfigure}

\begin{savaistu}[Why uniqueness matters]
The uniqueness theorem is the mathematical reason why prediction is possible. Because exactly one curve passes through $(x_0,y_0)$, knowing the present state of a system (population, temperature, charge) and its law of change determines its entire future. This is the principle behind weather forecasting models, the design of cooling systems, and the modelling of epidemic spread studied by Cameroonian health researchers.
\end{savaistu}

\begin{aretenir}[Key points]
\begin{itemize}
\item An initial condition $y(x_0)=y_0$ fixes the constant $C$: substitute, solve, rewrite.
\item If algebra gives several candidates ($\pm$, sign of $A$), the initial condition selects the correct branch.
\item If $g(y_0)=0$, the particular solution is the constant solution $y=y_0$.
\item Always finish with an explicit answer $y=\ldots$ (or a clean implicit relation) \textbf{and} the domain — an interval containing $x_0$.
\item Always \textbf{check}: does $y(x_0)=y_0$? Does the function satisfy the differential equation?
\end{itemize}
\end{aretenir}

\courssub{Practice: Exam-Style Initial Value Problems}

\begin{exercice}[Exercise 1 — direct substitution]
Find the particular solution of $\dfrac{dy}{dx}=3x^{2}e^{-y}$ given that $y=0$ when $x=0$.
\end{exercice}

\begin{corrige}[Exercise 1]
Separate: $\int e^{y}\,dy=\int 3x^{2}\,dx$, so $e^{y}=x^{3}+C$.
Substitute $x=0$, $y=0$: $e^{0}=1=0+C$, so $C=1$.
Hence $e^{y}=x^{3}+1$, i.e. $y=\ln(x^{3}+1)$, valid where $x^{3}+1>0$, i.e. $x>-1$. Since $x_0=0\in\,]{-1},+\infty[$, the solution is $y=\ln(x^{3}+1)$, $x>-1$.
\end{corrige}

\begin{exercice}[Exercise 2 — exponential decay]
The mass $m$ (in grams) of a radioactive substance satisfies $\dfrac{dm}{dt}=-0.02m$, with $t$ in days. Initially the mass is $80$ g. Find $m(t)$ and the mass remaining after $30$ days.
\end{exercice}

\begin{corrige}[Exercise 2]
$m=0$ is a constant solution but $m(0)=80\neq 0$. Separate: $\ln|m|=-0.02t+C$, so $m=Ae^{-0.02t}$.
Substitute $t=0$, $m=80$: $A=80$. Hence $m(t)=80e^{-0.02t}$, $t\geq 0$.
After $30$ days: $m(30)=80e^{-0.6}\approx 80(0.5488)\approx 43.9$ g.
\end{corrige}

\begin{exercice}[Exercise 3 — branch selection]
Solve the IVP $\dfrac{dy}{dx}=\dfrac{x}{y^{2}}$, $y(0)=-1$.
\end{exercice}

\begin{corrige}[Exercise 3]
$y=0$ is excluded by the equation itself; $y(0)=-1\neq 0$. Separate: $\int y^{2}\,dy=\int x\,dx$, so $\dfrac{y^{3}}{3}=\dfrac{x^{2}}{2}+C$, i.e. $y^{3}=\frac{3}{2}x^{2}+A$.
Substitute $x=0$, $y=-1$: $(-1)^{3}=A$, so $A=-1$.
Hence $y^{3}=\frac{3}{2}x^{2}-1$ and $y=\sqrt[3]{\frac{3}{2}x^{2}-1}$ (the cube root needs no sign choice and automatically gives $y(0)=-1$). Domain: all real $x$, since the cube root is defined everywhere; the solution is $y=\sqrt[3]{\frac{3x^{2}}{2}-1}$, $x\in\mathbb{R}$.
\end{corrige}

\begin{exercice}[Exercise 4 — trigonometric initial condition]
Find the particular solution of $\dfrac{dy}{dx}=y\cos x$ with $y\!\left(\frac{\pi}{2}\right)=e$.
\end{exercice}

\begin{corrige}[Exercise 4]
$y=0$ is a constant solution but $y(\pi/2)=e\neq 0$. Separate: $\int\frac{dy}{y}=\int\cos x\,dx$, so $\ln|y|=\sin x+C$, i.e. $y=Ae^{\sin x}$.
Substitute $x=\frac{\pi}{2}$, $y=e$: $e=Ae^{\sin(\pi/2)}=Ae^{1}$, so $A=1$.
Hence $y=e^{\sin x}$, $x\in\mathbb{R}$. Check: $y(\pi/2)=e^{1}=e$ \checkmark.
\end{corrige}

\begin{exercice}[Exercise 5 — cooling problem with two conditions]
A body cools according to Newton's law $\dfrac{dT}{dt}=-k(T-20)$, where the room temperature is $20^{\circ}$C. The body's initial temperature is $T(0)=100^{\circ}$C, and after $10$ minutes it has fallen to $60^{\circ}$C. Determine $k$, write $T(t)$, and find when the temperature reaches $30^{\circ}$C.
\end{exercice}

\begin{corrige}[Exercise 5]
Separate: $\ln|T-20|=-kt+C$, so $T=20+Ae^{-kt}$.
First condition $T(0)=100$: $100=20+A$, so $A=80$, giving $T=20+80e^{-kt}$.
Second condition $T(10)=60$: $60=20+80e^{-10k}$, so $e^{-10k}=\frac{40}{80}=\frac{1}{2}$, hence $-10k=\ln\frac{1}{2}=-\ln 2$ and $k=\frac{\ln 2}{10}\approx 0.0693$.
Thus $T(t)=20+80e^{-\frac{\ln 2}{10}t}=20+80\left(\frac{1}{2}\right)^{t/10}$, $t\geq 0$.
For $T=30$: $30=20+80e^{-kt}$, so $e^{-kt}=\frac{10}{80}=\frac{1}{8}=\left(\frac{1}{2}\right)^{3}$, giving $\frac{t}{10}=3$, i.e. $t=30$ minutes.
\end{corrige}

\begin{exercice}[Exercise 6 — rational function with restricted domain]
Solve the IVP $\dfrac{dy}{dx}=\dfrac{y^2}{x}$, $y(1)=2$.
\end{exercice}

\begin{corrige}[Exercise 6]
Note $y=0$ is a constant solution, but $y(1)=2\neq 0$. Separate: $\int\frac{dy}{y^2}=\int\frac{dx}{x}$, so $-\frac{1}{y}=\ln|x|+C$.
Substitute $x=1$, $y=2$: $-\frac{1}{2}=\ln 1 + C = C$. Hence $-\frac{1}{y}=\ln|x|-\frac{1}{2}$.
Solve for $y$: $\frac{1}{y}=\frac{1}{2}-\ln|x|$, so $y=\frac{1}{\frac{1}{2}-\ln|x|}=\frac{2}{1-2\ln|x|}$.
Domain: we need $x\neq 0$ (from $\ln|x|$) and $1-2\ln|x|\neq 0 \Rightarrow \ln|x|\neq \frac{1}{2} \Rightarrow |x|\neq \sqrt{e}$.
The initial point is $x_0=1>0$, so we take the interval $(0,\sqrt{e})$ containing $1$.
\[ \boxed{\,y=\frac{2}{1-2\ln x},\qquad x\in\,]0,\sqrt{e}[\,}. \]
\end{corrige}

\coursec{Applications of Separable Differential Equations and Exam Preparation}

In the previous section we learnt how to pin down the constant of integration using an initial condition, so that a whole family of curves is reduced to the single curve passing through a given point. We now put this skill to work where it truly matters: real situations. Almost every differential equation you will meet at the GCE Advanced Level comes wrapped in a story --- a population of bacteria, a cooling cup of tea, a radioactive sample, a tank of brine. The examination does not only test whether you can integrate; it tests whether you can \emph{translate} a story into mathematics, solve it, and \emph{interpret} the answer. This section trains exactly that skill, model by model, and ends with examination strategy.

\courssub{Exponential Growth and Decay: $\dfrac{dP}{dt}=kP$}

\textbf{Intuition.} Many quantities grow (or shrink) at a rate proportional to their current size: a bacterial colony doubles faster when it is larger, because more cells are dividing; a radioactive sample decays faster when more unstable nuclei remain; money in a savings account earns interest on the interest already earned. In each case the rate of change is proportional to the quantity itself. This single observation produces the most important differential equation of the whole syllabus.

\begin{propriete}[Exponential growth/decay model]
If a quantity $P(t)$ satisfies
\[
\frac{dP}{dt}=kP,
\]
where $k$ is a constant, then its general solution is
\[
P(t)=P_{0}\,e^{kt},\qquad P_{0}=P(0).
\]
If $k>0$ we have \cle{exponential growth}; if $k<0$ we have \cle{exponential decay}. Writing $k=-\lambda$ with $\lambda>0$ for decay, the \cle{half-life} (time for the quantity to be halved) is
\[
t_{1/2}=\frac{\ln 2}{\lambda}.
\]
For growth ($k>0$), the \cle{doubling time} is $t_{2}=\dfrac{\ln 2}{k}$.
\end{propriete}

\textbf{Derivation (examinable).} Separating variables:
\[
\frac{dP}{P}=k\,dt
\;\Longrightarrow\;
\ln|P|=kt+C
\;\Longrightarrow\;
P=Ae^{kt},\quad A=\pm e^{C}.
\]
The initial condition $P(0)=P_{0}$ gives $A=P_{0}$.

\textbf{Derivation of the half-life.} For decay, $P=P_{0}e^{-\lambda t}$. The half-life is defined by $P(t_{1/2})=\dfrac{P_{0}}{2}$:
\[
\frac{P_{0}}{2}=P_{0}e^{-\lambda t_{1/2}}
\;\Longrightarrow\;
e^{-\lambda t_{1/2}}=\frac12
\;\Longrightarrow\;
-\lambda t_{1/2}=-\ln 2
\;\Longrightarrow\;
t_{1/2}=\frac{\ln 2}{\lambda}.
\]
Note that $P_{0}$ cancels: the half-life \emph{does not depend on the initial amount}. This is the signature of exponential decay, and it is why half-life is such a useful concept.

\begin{popfigure}
\begin{center}
\begin{tikzpicture}
\begin{axis}[
  width=11cm, height=7cm,
  axis lines=left,
  xlabel={$t$}, ylabel={$P$},
  xmin=0, xmax=10, ymin=0, ymax=9,
  xtick={0,2,4,6,8,10},
  ytick={0,2,4,6,8},
  grid=both, grid style={popline, very thin},
  legend style={at={(0.03,0.97)}, anchor=north west, font=\small}
]
\addplot[domain=0:10, samples=80, thick, popgreen] {4*exp(0.2*x)};
\addlegendentry{growth: $4e^{0.2t}$}
\addplot[domain=0:10, samples=80, thick, poporange, dashed] {8*exp(-0.3*x)};
\addlegendentry{decay: $8e^{-0.3t}$}
\end{axis}
\end{tikzpicture}
\end{center}
\legende{Exponential growth ($k>0$) curves away from the axis ever faster; exponential decay ($k<0$) approaches zero asymptotically, losing the same fraction in every equal time interval.}
\end{popfigure}

\begin{exoresolu}[Radioactive decay]
A sample contains $N_{0}=8.0\times 10^{6}$ radioactive nuclei and decays according to $\dfrac{dN}{dt}=-\lambda N$. After $10$ days, only half remain. Find $\lambda$ and the number of nuclei after $30$ days.
\tcblower
\textbf{Model.} The solution is $N=N_{0}e^{-\lambda t}$.

\textbf{Find $\lambda$.} The half-life condition gives
\[
\frac{N_{0}}{2}=N_{0}e^{-10\lambda}
\;\Longrightarrow\;
\lambda=\frac{\ln 2}{10}\approx 0.0693\ \text{day}^{-1}.
\]
\textbf{Evaluate.} After $30$ days:
\[
N(30)=8.0\times 10^{6}\,e^{-30\lambda}=8.0\times 10^{6}\,e^{-3\ln 2}
=8.0\times 10^{6}\times \frac{1}{8}=1.0\times 10^{6}.
\]
\textbf{Check.} $30$ days is three half-lives, so the amount is divided by $2^{3}=8$: the answer is consistent.
\end{exoresolu}

\begin{exoresolu}[Bacterial growth]
A culture of bacteria grows according to $\dfrac{dP}{dt}=kP$. Initially there are $2000$ bacteria; after $3$ hours there are $16000$. Find $k$ and the size of the culture after $5$ hours.
\tcblower
\textbf{Model.} $P(t)=2000e^{kt}$.

\textbf{Find $k$.} At $t=3$: $16000=2000e^{3k}$, so $e^{3k}=8$, giving
\[
k=\frac{\ln 8}{3}=\frac{3\ln 2}{3}=\ln 2\approx 0.693\ \text{h}^{-1}.
\]
So the population doubles every hour.

\textbf{Evaluate.} $P(5)=2000e^{5\ln 2}=2000\times 2^{5}=64000$ bacteria.

\textbf{Check.} After $5$ hours there have been $5$ doublings from $2000$: $2000\to 4000\to 8000\to 16000\to 32000\to 64000$. Consistent.
\end{exoresolu}

\courssub{Newton's Law of Cooling}

\textbf{Intuition.} A hot bowl of \emph{achu} left on the table does not cool at a constant rate: it cools quickly at first (when it is much hotter than the room) and ever more slowly as its temperature approaches room temperature. Newton observed that the rate of cooling is proportional to the \emph{difference} between the object's temperature and the ambient temperature.

\begin{propriete}[Newton's law of cooling]
If $T(t)$ is the temperature of a body placed in an environment at constant temperature $T_{\mathrm{env}}$, then
\[
\frac{dT}{dt}=-k\bigl(T-T_{\mathrm{env}}\bigr),\qquad k>0,
\]
with solution
\[
T(t)=T_{\mathrm{env}}+\bigl(T_{0}-T_{\mathrm{env}}\bigr)e^{-kt}.
\]
\end{propriete}

\textbf{Derivation.} Let $u=T-T_{\mathrm{env}}$. Since $T_{\mathrm{env}}$ is constant, $\dfrac{du}{dt}=\dfrac{dT}{dt}=-ku$, which is exponential decay of $u$. Hence $u=u_{0}e^{-kt}$, i.e.\ $T-T_{\mathrm{env}}=(T_{0}-T_{\mathrm{env}})e^{-kt}$.

\begin{attention}[Cool down \emph{towards}, not \emph{to}, zero]
The body never cools below room temperature, and its temperature tends to $T_{\mathrm{env}}$, not to $0$. The equation must therefore contain $T-T_{\mathrm{env}}$, not $T$ alone. Writing $\dfrac{dT}{dt}=-kT$ would predict that your tea eventually freezes in a warm room --- a clear sign that the model is wrong.
\end{attention}

\begin{popfigure}
\begin{center}
\begin{tikzpicture}
\begin{axis}[
  width=11cm, height=7cm,
  axis lines=left,
  xlabel={$t$ (min)}, ylabel={$T$ ($^\circ$C)},
  xmin=0, xmax=60, ymin=15, ymax=100,
  xtick={0,10,20,30,40,50,60},
  ytick={20,40,60,80,100},
  grid=both, grid style={popline, very thin},
  legend style={at={(0.97,0.97)}, anchor=north east, font=\small}
]
\addplot[domain=0:60, samples=80, thick, popblue] {20+80*exp(-0.05*x)};
\addlegendentry{$k=0.05$}
\addplot[domain=0:60, samples=80, thick, poporange, dashed] {20+80*exp(-0.12*x)};
\addlegendentry{$k=0.12$}
\addplot[domain=0:60, samples=2, thin, popmuted] {20};
\addlegendentry{$T_{\mathrm{env}}=20$}
\end{axis}
\end{tikzpicture}
\end{center}
\legende{Cooling curves $T(t)=T_{\mathrm{env}}+(T_{0}-T_{\mathrm{env}})e^{-kt}$ with $T_{0}=100$, $T_{\mathrm{env}}=20$. Larger $k$ means faster cooling; both curves approach $T_{\mathrm{env}}$ asymptotically.}
\end{popfigure}

\begin{exoresolu}[Cooling of a liquid]
A liquid at $90^\circ\mathrm{C}$ is placed in a room at $25^\circ\mathrm{C}$. After $10$ minutes its temperature is $60^\circ\mathrm{C}$. Find its temperature after $25$ minutes.
\tcblower
\textbf{Model.} $T(t)=25+65e^{-kt}$ since $T_{0}-T_{\mathrm{env}}=90-25=65$.

\textbf{Find $k$.} At $t=10$: $60=25+65e^{-10k}$, so $e^{-10k}=\dfrac{35}{65}=\dfrac{7}{13}$, giving
\[
k=\frac{1}{10}\ln\frac{13}{7}\approx 0.0619\ \text{min}^{-1}.
\]
\textbf{Evaluate.} At $t=25$:
\[
T(25)=25+65e^{-25k}=25+65\left(\frac{7}{13}\right)^{2.5}\approx 25+65(0.2127)\approx 38.8^\circ\mathrm{C}.
\]
\textbf{Check.} The temperature lies between $60^\circ\mathrm{C}$ and $25^\circ\mathrm{C}$ and is still decreasing --- consistent with the physics.
\end{exoresolu}

\courssub{Chemical Reaction Rates}

In a reaction $A\to$ products, the \cle{rate law} often takes the form
\[
\frac{dC}{dt}=-kC^{n},
\]
where $C(t)$ is the concentration of reactant $A$ and $n$ is the \cle{order} of the reaction.

\textbf{First-order kinetics ($n=1$).} $\dfrac{dC}{dt}=-kC$ gives $C=C_{0}e^{-kt}$: exponential decay, with half-life $\dfrac{\ln 2}{k}$ independent of $C_{0}$.

\textbf{Second-order kinetics ($n=2$).} $\dfrac{dC}{dt}=-kC^{2}$. Separating:
\[
-\frac{dC}{C^{2}}=k\,dt
\;\Longrightarrow\;
\frac{1}{C}=kt+C'
\;\Longrightarrow\;
\frac{1}{C}=kt+\frac{1}{C_{0}},
\]
using $C(0)=C_{0}$. So a plot of $\dfrac{1}{C}$ against $t$ is a straight line of slope $k$ --- a standard way of testing second-order behaviour experimentally.

\begin{exemplebox}[Second-order reaction]
A reactant obeys $\dfrac{dC}{dt}=-0.02\,C^{2}$ (concentration in $\mathrm{mol\,L^{-1}}$, time in minutes) with $C_{0}=0.5\ \mathrm{mol\,L^{-1}}$. Then
\[
\frac{1}{C}=0.02t+\frac{1}{0.5}=0.02t+2
\;\Longrightarrow\;
C(t)=\frac{1}{2+0.02t}.
\]
At $t=50$ min: $C=\dfrac{1}{3}\approx 0.333\ \mathrm{mol\,L^{-1}}$.
\end{exemplebox}

\begin{attention}[Units of the rate constant]
The units of $k$ depend on the order: for $n=1$, $k$ is in $\mathrm{time}^{-1}$; for $n=2$, $k$ is in $\mathrm{concentration}^{-1}\,\mathrm{time}^{-1}$. Examiners award marks for correct units.
\end{attention}

\courssub{Mixing Problems}

\textbf{The set-up.} A tank contains a solution (e.g.\ brine). A solution flows \emph{in} at a given rate with a given concentration, the mixture is kept uniform by stirring, and flows \emph{out} at a given rate. We want the mass $x(t)$ of solute (salt) in the tank.

The governing principle is a simple balance:
\[
\frac{dx}{dt}=\text{(rate of solute in)}-\text{(rate of solute out)}.
\]
The outflow concentration is $\dfrac{x(t)}{V(t)}$, where $V(t)$ is the volume of liquid in the tank at time $t$, because the stirring keeps the mixture uniform.

\begin{popfigure}
\begin{center}
\begin{tikzpicture}[scale=1.0, >=stealth]
% tank
\draw[thick, popink] (0,0) -- (0,3) (3,0) -- (3,3) (0,0) -- (3,0);
\draw[thick, popink] (0,3) -- (-0.3,3.4) (3,3) -- (3.3,3.4);
% liquid
\fill[popblueL, opacity=0.6] (0.05,0.05) rectangle (2.95,2.1);
\draw[popblue] (0.05,2.1) -- (2.95,2.1);
\node[popdark] at (1.5,1.0) {$x(t)$ g of salt};
\node[popdark] at (1.5,0.5) {volume $V(t)$ L};
% inflow
\draw[->, very thick, popgreen] (1.5,4.6) -- (1.5,3.5);
\node[popgreen, align=center, font=\small] at (1.5,5.0) {inflow: $r_{\mathrm{in}}$ L/min,\\ $c_{\mathrm{in}}$ g/L};
% outflow
\draw[->, very thick, poporange] (3,0.7) -- (4.3,0.7);
\node[poporange, align=center, font=\small] at (5.6,0.7) {outflow: $r_{\mathrm{out}}$ L/min,\\ conc.\ $x/V$ g/L};
% stirrer
\draw[thick, popmuted] (1.5,3.4) -- (1.5,2.3);
\draw[thick, popmuted] (1.1,2.3) -- (1.9,2.3);
\node[popmuted, font=\small] at (2.6,2.6) {stirrer};
\end{tikzpicture}
\end{center}
\legende{Mixing tank: solute enters at rate $r_{\mathrm{in}}c_{\mathrm{in}}$ and leaves at rate $r_{\mathrm{out}}\,x/V(t)$.}
\end{popfigure}

\begin{propriete}[Mixing equation]
With the notation above,
\[
\frac{dx}{dt}=r_{\mathrm{in}}c_{\mathrm{in}}-r_{\mathrm{out}}\frac{x}{V(t)},
\qquad V(t)=V_{0}+\bigl(r_{\mathrm{in}}-r_{\mathrm{out}}\bigr)t.
\]
When $r_{\mathrm{in}}=r_{\mathrm{out}}$ (constant volume) and $c_{\mathrm{in}}$ is constant, this equation is \cle{separable}: it has the form $\dfrac{dx}{dt}=a-bx$. When the volume varies, the equation generally requires an integrating factor --- beyond this chapter --- but the balance equation itself must still be set up correctly.
\end{propriete}

\begin{attention}[Volume may change]
If $r_{\mathrm{in}}\neq r_{\mathrm{out}}$, the volume is \emph{not} constant. Always write $V(t)=V_{0}+(r_{\mathrm{in}}-r_{\mathrm{out}})t$ \textbf{before} attempting to separate variables. Using $V_{0}$ throughout when the volume changes is the most common error in mixing questions.
\end{attention}

\begin{exoresolu}[Brine tank, constant volume]
A tank contains $100$ L of brine with $20$ g of dissolved salt. Brine containing $0.5$ g/L flows in at $4$ L/min, and the well-stirred mixture flows out at $4$ L/min. Find the mass of salt after $25$ minutes.
\tcblower
\textbf{Set up.} Volume constant: $V=100$ L. Rate in: $4\times 0.5=2$ g/min. Rate out: $4\times\dfrac{x}{100}=\dfrac{x}{25}$ g/min. Hence
\[
\frac{dx}{dt}=2-\frac{x}{25}=\frac{50-x}{25},\qquad x(0)=20.
\]
\textbf{Separate and integrate.}
\[
\frac{dx}{50-x}=\frac{dt}{25}
\;\Longrightarrow\;
-\ln|50-x|=\frac{t}{25}+C
\;\Longrightarrow\;
50-x=Ae^{-t/25}.
\]
\textbf{Initial condition.} $x(0)=20\Rightarrow A=30$, so $x(t)=50-30e^{-t/25}$.

\textbf{Evaluate.} $x(25)=50-30e^{-1}\approx 50-11.04\approx 38.96$ g.

\textbf{Interpret.} As $t\to\infty$, $x\to 50$ g, i.e.\ the concentration tends to $\dfrac{50}{100}=0.5$ g/L, the inflow concentration --- exactly as expected physically.
\end{exoresolu}

\begin{methode}[Translating a word problem into a separable DE]
\begin{enumerate}
\item \textbf{Identify the unknown quantity} and its variable (e.g.\ $P(t)$, $T(t)$, $x(t)$), with units.
\item \textbf{Identify the rate statement} in the text: ``proportional to'', ``rate of change'', ``half of'' etc., and name the proportionality constant $k$.
\item \textbf{Write the balance or rate equation} $\dfrac{dy}{dt}=\dots$ and check it is separable: the right-hand side must be a product (or quotient) of a function of $y$ and a function of $t$.
\item \textbf{Separate, integrate, and apply the initial condition} to find the constant.
\item \textbf{Answer the question asked} (a value, a time, a limit) and \textbf{interpret} it in the context, with units.
\end{enumerate}
\end{methode}

\courssub{Exam Strategies and Common Mistakes}

\begin{aretenir}[GCE exam strategy for DE questions]
\begin{itemize}
\item \textbf{Read twice.} Underline the rate statement and the initial condition; note all units.
\item \textbf{Identify the model}: growth/decay ($y'=ky$), cooling ($T'=-k(T-T_{\mathrm{env}})$), reaction ($C'=-kC^{n}$), mixing (balance equation).
\item \textbf{Show the separation step explicitly.} Marks are awarded for correct separation even if the subsequent integration contains a minor slip.
\item \textbf{Use the initial condition immediately} after integrating to fix the constant.
\item \textbf{Interpret and check}: does the sign make sense? Does the long-term limit agree with the physics? Are the units right?
\end{itemize}
\end{aretenir}

\begin{attention}[Frequent traps]
\begin{itemize}
\item Forgetting the constant of integration, or applying the initial condition before integrating.
\item Sign errors in decay problems: write $\dfrac{dP}{dt}=-kP$ with $k>0$ stated explicitly, rather than leaving the sign of $k$ ambiguous.
\item In cooling problems, writing $\dfrac{dT}{dt}=-kT$ instead of $-k(T-T_{\mathrm{env}})$.
\item In mixing problems, using the inflow concentration instead of $x/V$ for the outflow.
\item Giving a final answer without units or without rounding sensibly.
\end{itemize}
\end{attention}

\begin{exercice}[Practice]
\begin{enumerate}
\item A population grows according to $\dfrac{dP}{dt}=0.03P$ with $P(0)=5000$. Find the time needed for the population to double.
\item A body cools from $80^\circ\mathrm{C}$ to $50^\circ\mathrm{C}$ in $15$ minutes in a room at $20^\circ\mathrm{C}$. When will it reach $30^\circ\mathrm{C}$?
\item A tank of $200$ L contains $10$ g of salt. Pure water enters at $5$ L/min and the mixture leaves at $5$ L/min. Find $x(t)$ and the mass of salt after $40$ minutes.
\end{enumerate}
\end{exercice}

\begin{corrige}[Exercise 1]
$P=5000e^{0.03t}$. Doubling: $e^{0.03t}=2\Rightarrow t=\dfrac{\ln 2}{0.03}\approx 23.1$ time units.
\end{corrige}

\begin{corrige}[Exercise 2]
$T=20+60e^{-kt}$; $50=20+60e^{-15k}\Rightarrow e^{-15k}=\tfrac12$, so $k=\tfrac{\ln 2}{15}$. Then $30=20+60e^{-kt}\Rightarrow e^{-kt}=\tfrac16\Rightarrow t=\dfrac{\ln 6}{k}=\dfrac{15\ln 6}{\ln 2}\approx 38.8$ min.
\end{corrige}

\begin{corrige}[Exercise 3]
Rate in $=0$; rate out $=\dfrac{5x}{200}=\dfrac{x}{40}$. So $\dfrac{dx}{dt}=-\dfrac{x}{40}$, giving $x=10e^{-t/40}$. At $t=40$: $x=10e^{-1}\approx 3.68$ g.
\end{corrige}

\begin{savaistu}[Why this chapter matters beyond the exam]
Separable equations are the workhorses of quantitative science: pharmacists use them for drug elimination from the bloodstream, engineers for capacitor discharge, demographers for population projections in Cameroon's development planning, and nuclear physicists for carbon dating. Mastering the translation from words to equations is a skill you will reuse throughout higher education.
\end{savaistu}

\newpage

\coursec{Integration activity}

\courssub{Aims of the activity}

This activity brings together all the skills developed in this chapter in a single, authentic modelling task. Working in groups of three or four, you will:

\begin{itemize}
\item \textbf{Recognise} a first order differential equation with separable variables and justify why the method of separation applies;
\item \textbf{Separate the variables} correctly, taking care with a right-hand side that depends on time through given data;
\item \textbf{Integrate} both sides to obtain a \cle{general solution}, handling a piecewise-constant coefficient by splitting the integral over sub-intervals;
\item \textbf{Use an initial condition} to determine the constant of integration and hence a \cle{particular solution};
\item \textbf{Interpret} the solution: predict a harvest date, test the sensitivity of the prediction to a parameter, and discuss the limitations of the model.
\end{itemize}

By the end of the activity you should be able to move fluently through the complete modelling cycle: \emph{equation} $\rightarrow$ \emph{separation} $\rightarrow$ \emph{integration} $\rightarrow$ \emph{initial condition} $\rightarrow$ \emph{interpretation and critique}. This cycle is exactly what the GCE Advanced Level examiner rewards in structured modelling questions.

\courssub{Situation}

\textbf{Designing a Sustainable Harvest Plan for a Cocoa Farm.}\par
\medskip
The Nkongsamba area, in the Littoral Region of Cameroon, is one of the country's important cocoa-producing zones. A cooperative of smallholder farmers wants to plan the main harvest so that picking starts when the pods are ripe but before over-ripening and pod disease reduce the quality of the beans. The cooperative's agricultural adviser proposes the following simplified model for the \cle{ripeness index} $R(t)$ of an average pod, where $t$ is measured in months from the start of the growing season:
\[
\frac{dR}{dt}=k\bigl(T(t)-T_{\mathrm{opt}}\bigr)\,R,
\]
where:
\begin{itemize}
\item $R(t)$ is the ripeness index, a number between $0$ (green pod) and $1$ (fully ripe);
\item $T(t)$ is the average monthly temperature, in ${}^{\circ}$C;
\item $T_{\mathrm{opt}}=22\ {}^{\circ}$C is a reference temperature below which ripening effectively stops;
\item $k=0.06$ per month per ${}^{\circ}$C is a growth constant estimated from previous seasons.
\end{itemize}

The adviser explains that the model says: \emph{the rate of ripening is proportional to the current ripeness and to the excess of temperature above the reference value.}

Records from the cooperative's weather station give the following average monthly temperatures for the six-month growing season:

\begin{center}
\begin{tabular}{|l|c|c|c|c|c|c|}
\hline
\rowcolor{popblueL} Month $n$ & 1 & 2 & 3 & 4 & 5 & 6 \\
\hline Period & $0\le t<1$ & $1\le t<2$ & $2\le t<3$ & $3\le t<4$ & $4\le t<5$ & $5\le t\le 6$ \\
\hline $T$ (${}^{\circ}$C) & 24 & 25 & 26 & 26 & 25 & 24 \\
\hline
\end{tabular}
\end{center}

At the start of the season ($t=0$), field measurements give $R(0)=0.1$. The cooperative considers pods \textbf{harvest-ready} when $R=0.9$.

Your group's mission: produce a ripeness formula $R(t)$, predict the harvest month, and advise the cooperative on how reliable the prediction is.

\courssub{Tasks}

Work through the following questions in your group. Keep a written record: your group will present its solution to the class.

\begin{enumerate}
\item \textbf{Recognition.} Write the model in the form $\dfrac{dR}{dt}=f(t)\,g(R)$, identifying $f(t)$ and $g(R)$ explicitly. Hence explain why the equation is a first order differential equation with \cle{separable variables}. What is the order of the equation, and why?

\item \textbf{Separation.} Assuming $R>0$, separate the variables and write down the equation you obtain after integrating both sides with respect to their own variable. State clearly where the constant of integration appears.

\item \textbf{Integration with piecewise data.} Because $T(t)$ is constant on each month, show that for $t$ in month $n$,
\[
\ln R = k\sum_{j=1}^{n-1}(T_j-T_{\mathrm{opt}}) + k(T_n-T_{\mathrm{opt}})(t-(n-1)) + C,
\]
where $T_j$ is the temperature in month $j$. (Hint: split $\int_0^t k(T(s)-T_{\mathrm{opt}})\,ds$ month by month.)

\item \textbf{Initial condition.} Use $R(0)=0.1$ to find $C$, and hence write the particular solution in the form
\[
R(t)=0.1\exp\!\left(k\int_0^t \bigl(T(s)-T_{\mathrm{opt}}\bigr)\,ds\right).
\]

\item \textbf{Computing the accumulated heat.} Complete the following table of the accumulated exponent $A(t)=k\displaystyle\int_0^t (T(s)-T_{\mathrm{opt}})\,ds$ at the end of each month:

\begin{center}
\begin{tabular}{|l|c|c|c|c|c|c|c|}
\hline
\rowcolor{popblueL} $t$ (months) & 0 & 1 & 2 & 3 & 4 & 5 & 6 \\
\hline $A(t)$ & & & & & & & \\
\hline $R(t)$ & & & & & & & \\
\hline
\end{tabular}
\end{center}

\item \textbf{Harvest prediction.} Determine the value of $t$ for which $R(t)=0.9$. In which month should the cooperative begin the harvest? Give the answer as a number of months (to 2 decimal places) and interpret it in words.

\item \textbf{Sensitivity.} The constant $k=0.06$ is only an estimate. Repeat the harvest-time calculation with $k=0.05$ and $k=0.07$. Comment on how sensitive the harvest date is to $k$, and what this means for the cooperative's planning.

\item \textbf{Critical reflection.} State at least \textbf{three} assumptions or weaknesses of the model (think about what happens when $T<T_{\mathrm{opt}}$, about rainfall, pests, variation between pods, and the meaning of $R>1$). Suggest one improvement.

\item \textbf{Presentation.} Prepare a five-minute group presentation: one member justifies separability, one presents the integration, one the harvest prediction and sensitivity, and one the critique.
\end{enumerate}

\courssub{Model answer}

\textbf{1. Recognition.}\par
The equation $\dfrac{dR}{dt}=k\bigl(T(t)-T_{\mathrm{opt}}\bigr)R$ has the form $\dfrac{dR}{dt}=f(t)\,g(R)$ with
\[
f(t)=k\bigl(T(t)-T_{\mathrm{opt}}\bigr),\qquad g(R)=R.
\]
The highest derivative present is the \emph{first} derivative, so the equation is of \cle{first order}. Since the right-hand side is a product of a function of $t$ alone and a function of $R$ alone, the variables are \cle{separable}.

\textbf{2. Separation.}\par
Since $R>0$ we may divide both sides by $R$:
\[
\frac{1}{R}\,\frac{dR}{dt}=k\bigl(T(t)-T_{\mathrm{opt}}\bigr)
\quad\Longrightarrow\quad
\int\frac{dR}{R}=\int k\bigl(T(t)-T_{\mathrm{opt}}\bigr)\,dt.
\]
Integrating the left side gives $\ln R$ (as $R>0$, no absolute value is needed); the constant $C$ is placed on the right:
\[
\ln R = \int_0^t k\bigl(T(s)-T_{\mathrm{opt}}\bigr)\,ds + C.
\]

\textbf{3. Integration with piecewise data.}\par
On month $j$, $T(s)=T_j$ is constant, so over that month the integrand contributes $k(T_j-T_{\mathrm{opt}})$ per unit time. For $t$ in month $n$, split the integral into the completed months $1,\dots,n-1$ (each of length $1$) plus the partial current month of length $t-(n-1)$:
\[
\int_0^t k(T(s)-T_{\mathrm{opt}})\,ds
= k\sum_{j=1}^{n-1}(T_j-T_{\mathrm{opt}})\cdot 1 + k(T_n-T_{\mathrm{opt}})\bigl(t-(n-1)\bigr),
\]
which gives the stated formula.

\textbf{4. Initial condition.}\par
Setting $t=0$: the integral from $0$ to $0$ is $0$, so $\ln R(0)=C$, giving $C=\ln 0.1$. Therefore
\[
\ln R = \ln 0.1 + k\int_0^t\bigl(T(s)-T_{\mathrm{opt}}\bigr)\,ds
\quad\Longrightarrow\quad
R(t)=0.1\exp\!\left(k\int_0^t \bigl(T(s)-T_{\mathrm{opt}}\bigr)\,ds\right).
\]

\textbf{5. Accumulated heat.}\par
With $k=0.06$ and $T_{\mathrm{opt}}=22$: the monthly excesses $T_j-22$ are $2,3,4,4,3,2$. Multiplying by $k$ gives monthly increments $0.12,\ 0.18,\ 0.24,\ 0.24,\ 0.18,\ 0.12$. Accumulating and computing $R=0.1e^{A(t)}$:

\begin{center}
\begin{tabular}{|l|c|c|c|c|c|c|c|}
\hline
\rowcolor{popblueL} $t$ (months) & 0 & 1 & 2 & 3 & 4 & 5 & 6 \\
\hline $A(t)$ & 0 & 0.12 & 0.30 & 0.54 & 0.78 & 0.96 & 1.08 \\
\hline $R(t)$ & 0.100 & 0.113 & 0.135 & 0.172 & 0.218 & 0.261 & 0.294 \\
\hline
\end{tabular}
\end{center}

For example, $R(3)=0.1e^{0.54}\approx 0.1\times 1.716 = 0.172$.

\begin{popfigure}
\begin{center}
\begin{tikzpicture}
\begin{axis}[
width=0.85\linewidth, height=6cm,
xlabel={$t$ (months)}, ylabel={$R(t)$},
xmin=0, xmax=6.2, ymin=0, ymax=1,
xtick={0,1,2,3,4,5,6},
ytick={0,0.1,0.3,0.5,0.7,0.9},
grid=both, grid style={popline, very thin},
legend style={at={(0.03,0.97)}, anchor=north west, font=\small},
]
\addplot[popblue, very thick, domain=0:1, samples=20] {0.1*exp(0.12*x)};
\addplot[popblue, very thick, domain=1:2, samples=20] {0.1*exp(0.12+0.18*(x-1))};
\addplot[popblue, very thick, domain=2:3, samples=20] {0.1*exp(0.30+0.24*(x-2))};
\addplot[popblue, very thick, domain=3:4, samples=20] {0.1*exp(0.54+0.24*(x-3))};
\addplot[popblue, very thick, domain=4:5, samples=20] {0.1*exp(0.78+0.18*(x-4))};
\addplot[popblue, very thick, domain=5:6, samples=20] {0.1*exp(0.96+0.12*(x-5))};
\addlegendentry{$R(t)=0.1e^{A(t)}$}
\addplot[poporange, dashed, thick, domain=0:6.2] {0.9};
\addlegendentry{harvest threshold $R=0.9$}
\end{axis}
\end{tikzpicture}
\end{center}
\legende{The ripeness index grows exponentially within each month, but stays far below the harvest threshold $R=0.9$ throughout the six-month season.}
\end{popfigure}

\textbf{6. Harvest prediction.}\par
We need $R(t)=0.9$, i.e. $0.1e^{A(t)}=0.9$, so
\[
A(t)=\ln 9 \approx 2.1972.
\]
But the table shows $A(6)=1.08 < 2.1972$: within the six-month season the accumulated exponent is far too small, and $R$ only reaches about $0.29$. \textbf{Conclusion: under this model with $k=0.06$, the pods do not reach harvest ripeness within six months.} Solving formally, with the seasonal pattern repeating, the average monthly increment of $A$ is $\frac{1.08}{6}=0.18$, so $A(t)\approx 0.18\,t$ and
\[
t \approx \frac{\ln 9}{0.18}\approx 12.21\ \text{months}.
\]
The model therefore predicts harvest-readiness after roughly \textbf{12.2 months}, i.e. only in the \emph{following} growing season — a signal that the model (or the value of $k$) needs revision, since real cocoa pods mature within a single season.

\textbf{7. Sensitivity.}\par
Repeating with the average monthly increment scaled by $k$:
\begin{itemize}
\item $k=0.05$: increment $=\dfrac{0.05}{0.06}\times 0.18=0.15$, so $t\approx \dfrac{2.1972}{0.15}\approx 14.65$ months;
\item $k=0.07$: increment $=\dfrac{0.07}{0.06}\times 0.18=0.21$, so $t\approx \dfrac{2.1972}{0.21}\approx 10.46$ months.
\end{itemize}
A change of about $17\%$ in $k$ shifts the predicted harvest time by roughly two months in either direction: the prediction is \textbf{quite sensitive} to $k$. The cooperative should therefore treat $k$ as the key quantity to measure more accurately, e.g. by recording ripeness over one full season and fitting the model.

\textbf{8. Critical reflection.}\par
Weaknesses of the model include:
\begin{itemize}
\item If $T<T_{\mathrm{opt}}$, the model gives $\frac{dR}{dt}<0$: pods would ``un-ripen'', which is biologically unrealistic — ripening should at worst stop;
\item rainfall, humidity, sunlight and soil fertility are ignored, though they strongly affect cocoa development;
\item all pods are described by a single index $R(t)$, ignoring variation between trees and pods;
\item nothing prevents $R$ from exceeding $1$, although the index is defined only up to $1$;
\item $k$ and $T_{\mathrm{opt}}$ are rough estimates, and monthly average temperatures hide daily fluctuations.
\end{itemize}
One improvement: replace the factor $(T-T_{\mathrm{opt}})$ by $\max(T-T_{\mathrm{opt}},0)$ so ripening never reverses, and use weekly temperature data for a finer integration.

\textbf{9. Presentation.}\par
Groups present their work; the class compares harvest predictions and discusses which assumption is most damaging to the model's credibility.

\begin{attention}[Common traps in this activity]
Do not integrate $k(T(t)-T_{\mathrm{opt}})$ as if $T$ were a single constant — you must split the integral month by month. Do not forget that $\ln R$ requires $R>0$ (justified here by $R(0)=0.1>0$). Finally, when the model contradicts reality (harvest after 12 months), say so honestly: recognising a model's failure is a valued modelling skill, not a mistake.
\end{attention}

\newpage

\coursec{Methods and worked examples}

\courssub{Skill 1: Verifying that a Function is a Solution of a Differential Equation}

A \cle{differential equation} is an equation that involves an unknown function $y$ and one or more of its derivatives. When only the first derivative $\dfrac{dy}{dx}$ appears, the equation is said to be of \cle{first order}. A \cle{solution} of such an equation is a function $y = f(x)$ which, together with its derivative, satisfies the equation identically on some interval. Before learning how to \emph{find} solutions, we must be able to \emph{check} a proposed solution — a skill that the GCE tests directly and that rewards careful, methodical work.

\begin{methode}[Verifying a Proposed Solution]
To check whether a given function $y = f(x)$ is a solution of the differential equation $F(x, y, y') = 0$:
\begin{enumerate}
\item Compute the derivative $y' = f'(x)$ (and $y''$ if needed).
\item Substitute $y$ and $y'$ into the \textbf{left-hand side} of the equation.
\item Simplify carefully and compare with the right-hand side.
\item If the equality holds for \textbf{all} $x$ in the domain, conclude that $f$ is a solution; otherwise it is not.
\item If an initial condition $y(x_0) = y_0$ is given, check it separately by direct substitution.
\end{enumerate}
\textbf{Reflex:} never ``solve'' anything here — you only \emph{verify}. This is a frequent source of easy marks at the GCE.
\end{methode}

\begin{exoresolu}[Verifying a Solution]
Show that the function $y = 3e^{2x}$ is a solution of the differential equation $\dfrac{dy}{dx} - 2y = 0$, and that it satisfies the initial condition $y(0) = 3$.
\tcblower
\textbf{Solution.}\par
\textbf{Step 1: differentiate.} Given $y = 3e^{2x}$, we obtain
\[ \frac{dy}{dx} = 3 \times 2e^{2x} = 6e^{2x}. \]
\textbf{Step 2: substitute into the left-hand side.}
\[ \frac{dy}{dx} - 2y = 6e^{2x} - 2\bigl(3e^{2x}\bigr) = 6e^{2x} - 6e^{2x} = 0. \]
\textbf{Step 3: conclude.} The left-hand side equals the right-hand side for every real $x$, so $y = 3e^{2x}$ \textbf{is a solution} of the equation on $\mathbb{R}$.\par
\textbf{Step 4: check the initial condition.} $y(0) = 3e^{0} = 3 \times 1 = 3$. The condition $y(0) = 3$ is satisfied.\par
Hence $y = 3e^{2x}$ is the particular solution of $\dfrac{dy}{dx} - 2y = 0$ satisfying $y(0) = 3$.
\end{exoresolu}

\begin{exoresolu}[Verifying an Implicitly Defined Solution]
Show that the relation $x^2 + y^2 = 25$ defines a solution of the differential equation $y\dfrac{dy}{dx} + x = 0$.
\tcblower
\textbf{Solution.}\par
\textbf{Step 1: differentiate implicitly} with respect to $x$:
\[ \frac{d}{dx}\bigl(x^2 + y^2\bigr) = \frac{d}{dx}(25) \quad\Longrightarrow\quad 2x + 2y\frac{dy}{dx} = 0. \]
\textbf{Step 2: simplify.} Dividing by $2$:
\[ y\frac{dy}{dx} + x = 0, \]
which is exactly the given equation. Hence every function $y$ defined implicitly by $x^2 + y^2 = 25$ (i.e.\ $y = \pm\sqrt{25 - x^2}$ on $-5 < x < 5$) is a solution. Geometrically, the solution curves are semicircles centred at the origin.
\end{exoresolu}

\courssub{Skill 2: Recognising a Separable Equation}

Not every first-order equation can be solved by the techniques of this chapter. The crucial first step is to recognise those equations whose right-hand side splits into a part depending only on $x$ and a part depending only on $y$.

\begin{definition}[Separable First-Order Equation]
A first-order differential equation is called \cle{separable} (or ``with separable variables'') if it can be written in the form
\[ \frac{dy}{dx} = g(x)\,h(y), \]
where $g$ is a function of $x$ alone and $h$ is a function of $y$ alone.
\end{definition}

\begin{methode}[Testing for Separability]
Practical reflexes:
\begin{enumerate}
\item Try to \textbf{factorise} the right-hand side: look for common factors, e.g.\ $xy + x = x(y+1)$.
\item Powers multiply across variables: $e^{x+y} = e^{x}e^{y}$ is separable; $e^{x} + y$ is \textbf{not}.
\item Sums that cannot be factorised, e.g.\ $\dfrac{dy}{dx} = x + y$, are \textbf{not} separable.
\item Equations of the form $\dfrac{dy}{dx} = g(x)$ (no $y$) and $\dfrac{dy}{dx} = h(y)$ (no $x$) are trivially separable.
\end{enumerate}
\end{methode}

\begin{exoresolu}[Recognising Separable Equations]
For each equation, state whether it is separable. If it is, write it in the form $\dfrac{dy}{dx} = g(x)\,h(y)$.
\begin{enumerate}
\item $\dfrac{dy}{dx} = xy^2 + x$ \qquad (b)\; $\dfrac{dy}{dx} = e^{x - y}$ \qquad (c)\; $\dfrac{dy}{dx} = x^2 + y^2$ \qquad (d)\; $\dfrac{dy}{dx} = \dfrac{\sin x}{y}$
\end{enumerate}
\tcblower
\textbf{Solution.}\par
\textbf{(a)} Factorise the right-hand side: $xy^2 + x = x(y^2 + 1)$. This is a product of $g(x) = x$ and $h(y) = y^2 + 1$. \textbf{Separable.}\par
\textbf{(b)} Using the law of indices, $e^{x-y} = e^{x}\,e^{-y} = e^{x}\cdot\dfrac{1}{e^{y}}$. With $g(x) = e^{x}$ and $h(y) = e^{-y}$, the equation is \textbf{separable}.\par
\textbf{(c)} $x^2 + y^2$ cannot be written as a product of a function of $x$ alone and a function of $y$ alone (any attempted factorisation introduces cross terms in $xy$). \textbf{Not separable.}\par
\textbf{(d)} $\dfrac{\sin x}{y} = \sin x \cdot \dfrac{1}{y}$, a product of $g(x) = \sin x$ and $h(y) = \dfrac{1}{y}$. \textbf{Separable.}
\end{exoresolu}

\begin{attention}[Sums versus Products]
The single most common error is to treat $\dfrac{dy}{dx} = x + y$ as separable. A \textbf{sum} of an $x$-term and a $y$-term is \emph{not} separable unless it factorises. Always ask: ``Can I write the right-hand side as one factor in $x$ times one factor in $y$?''
\end{attention}

\courssub{Skill 3: Separating Variables and Finding a General Solution (Polynomial/Algebraic Cases)}

\begin{methode}[Solving a Separable Equation — General Solution]
\begin{enumerate}
\item Write the equation as $\dfrac{dy}{dx} = g(x)\,h(y)$.
\item \textbf{Separate}: divide by $h(y)$ (noting any constant solutions where $h(y) = 0$) and multiply by $dx$:
\[ \frac{dy}{h(y)} = g(x)\,dx. \]
\item \textbf{Integrate both sides}: $\displaystyle\int \frac{dy}{h(y)} = \int g(x)\,dx$.
\item Add \textbf{one} constant of integration only (a single $C$ covers both sides).
\item If possible, make $y$ the subject to obtain the \cle{general solution} explicitly; otherwise leave it implicit.
\end{enumerate}
\textbf{Key formulas:} $\displaystyle\int x^n\,dx = \frac{x^{n+1}}{n+1}\ (n \neq -1)$, $\displaystyle\int \frac{dy}{y} = \ln|y|$, $\displaystyle\int e^{ax}\,dx = \frac{1}{a}e^{ax}$.
\end{methode}

\begin{aretenir}[Why One Constant is Enough]
Integrating both sides formally gives $F(y) + C_1 = G(x) + C_2$. Subtracting, $F(y) = G(x) + (C_2 - C_1)$, and since $C_2 - C_1$ is just an arbitrary constant, we rename it $C$. Writing two constants is not wrong, but failing to merge them loses presentation marks.
\end{aretenir}

\begin{exoresolu}[General Solution of a Separable Equation]
Find the general solution of the differential equation
\[ \frac{dy}{dx} = \frac{x}{y}, \qquad y \neq 0. \]
\tcblower
\textbf{Solution.}\par
\textbf{Step 1: identify.} The right-hand side is $x \cdot \dfrac{1}{y}$, a product of a function of $x$ and a function of $y$: the equation is separable.\par
\textbf{Step 2: separate the variables.} Multiply both sides by $y\,dx$:
\[ y\,dy = x\,dx. \]
\textbf{Step 3: integrate both sides.}
\[ \int y\,dy = \int x\,dx \quad\Longrightarrow\quad \frac{y^2}{2} = \frac{x^2}{2} + C_1. \]
\textbf{Step 4: simplify.} Multiplying through by $2$ and writing $C = 2C_1$:
\[ y^2 = x^2 + C. \]
\textbf{Step 5: make $y$ the subject.}
\[ y = \pm\sqrt{x^2 + C}. \]
\textbf{Conclusion.} The general solution is $y^2 = x^2 + C$, i.e.\ $y = \pm\sqrt{x^2 + C}$, where $C$ is an arbitrary constant. Geometrically, the solution curves are rectangular hyperbolas $y^2 - x^2 = C$.\par
\textbf{Check:} differentiating $y^2 = x^2 + C$ implicitly gives $2y\dfrac{dy}{dx} = 2x$, i.e.\ $\dfrac{dy}{dx} = \dfrac{x}{y}$. $\checkmark$
\end{exoresolu}

\begin{popfigure}
\begin{center}
\begin{tikzpicture}
\begin{axis}[
  width=9.5cm, height=7cm,
  axis lines=middle,
  xlabel={$x$}, ylabel={$y$},
  xmin=-3.2, xmax=3.2, ymin=-3.2, ymax=3.2,
  xtick={-2,-1,1,2}, ytick={-2,-1,1,2},
  tick label style={font=\scriptsize},
  samples=100, domain=-3:3, clip=true,
]
\addplot[popblue, thick] {sqrt(max(0,x^2+1))};
\addplot[popblue, thick] {-sqrt(max(0,x^2+1))};
\addplot[poporange, thick] {sqrt(max(0,x^2+3))};
\addplot[poporange, thick] {-sqrt(max(0,x^2+3))};
\addplot[popgreen, thick, domain=0.35:3] {sqrt(max(0,x^2-0.12))};
\addplot[popgreen, thick, domain=0.35:3] {-sqrt(max(0,x^2-0.12))};
\addplot[popgreen, thick, domain=-3:-0.35] {sqrt(max(0,x^2-0.12))};
\addplot[popgreen, thick, domain=-3:-0.35] {-sqrt(max(0,x^2-0.12))};
\node[font=\scriptsize, text=popblue] at (axis cs:1.6,2.3) {$C=1$};
\node[font=\scriptsize, text=poporange] at (axis cs:1.15,2.35) {$C=3$};
\node[font=\scriptsize, text=popgreen] at (axis cs:2.6,1.6) {$C<0$};
\end{axis}
\end{tikzpicture}
\end{center}
\legende{The family of solution curves $y^2 = x^2 + C$ of $\dfrac{dy}{dx} = \dfrac{x}{y}$: one arbitrary constant, one family of curves.}
\end{popfigure}

\courssub{Skill 4: Equations Leading to Logarithms and Exponentials}

\begin{methode}[Handling $\int \frac{dy}{y}$ and Exponential Answers]
When separation produces $\displaystyle\int \frac{dy}{y}$:
\begin{enumerate}
\item Integrate to $\ln|y| = \int g(x)\,dx + C$.
\item \textbf{Exponentiate}: $|y| = e^{\int g(x)dx + C} = e^{C}\,e^{\int g(x)dx}$.
\item Absorb the sign and the constant: write $y = A\,e^{\int g(x)dx}$ where $A = \pm e^{C}$ is an arbitrary \textbf{non-zero} constant; the case $A = 0$ (the zero solution) is usually added by inspection.
\item For $\dfrac{dy}{dx} = ky$, memorise the model result: $y = Ae^{kx}$.
\end{enumerate}
\begin{attention}[Forgetting the Constant or the Absolute Value]
Do not write $\ln y$ instead of $\ln|y|$, and do not add a constant on \emph{both} sides — one constant is enough. At the GCE, a missing constant costs marks even if everything else is correct.
\end{attention}
\end{methode}

\begin{propriete}[The Exponential Model (proof examinable)]
The general solution of $\dfrac{dy}{dx} = ky$ (where $k$ is a constant) is
\[ y = A e^{kx}, \qquad A \in \mathbb{R}. \]
\textbf{Proof.} For $y \neq 0$, separation gives $\displaystyle\int \frac{dy}{y} = \int k\,dx$, so $\ln|y| = kx + C$. Exponentiating, $|y| = e^{C}e^{kx}$, hence $y = \pm e^{C} e^{kx}$. Setting $A = \pm e^{C}$ and adjoining the zero solution $y = 0$ (which corresponds to $A = 0$) gives $y = Ae^{kx}$, $A \in \mathbb{R}$. $\blacksquare$
\end{propriete}

\begin{exoresolu}[Exponential-Type General Solution]
Find the general solution of
\[ \frac{dy}{dx} = \frac{3x^2}{y}, \qquad y > 0, \]
and the general solution of $\dfrac{dy}{dx} = 2xy$.
\tcblower
\textbf{Solution.}\par
\textbf{First equation.} Separate: $y\,dy = 3x^2\,dx$. Integrate:
\[ \int y\,dy = \int 3x^2\,dx \quad\Longrightarrow\quad \frac{y^2}{2} = x^3 + C_1. \]
Hence $y^2 = 2x^3 + C$ (with $C = 2C_1$), and since $y > 0$:
\[ y = \sqrt{2x^3 + C}. \]
\textbf{Second equation.} $\dfrac{dy}{dx} = 2xy = 2x \cdot y$ is separable. Note $y = 0$ is a constant solution. For $y \neq 0$, divide by $y$ and multiply by $dx$:
\[ \frac{dy}{y} = 2x\,dx. \]
Integrate:
\[ \int \frac{dy}{y} = \int 2x\,dx \quad\Longrightarrow\quad \ln|y| = x^2 + C. \]
Exponentiate:
\[ |y| = e^{x^2 + C} = e^{C}e^{x^2} \quad\Longrightarrow\quad y = \pm e^{C} e^{x^2}. \]
Writing $A = \pm e^{C}$ (and allowing $A = 0$ to include the zero solution):
\[ y = A e^{x^2}, \quad A \in \mathbb{R}. \]
\textbf{Check:} $\dfrac{dy}{dx} = A \cdot 2x e^{x^2} = 2x\bigl(Ae^{x^2}\bigr) = 2xy$. $\checkmark$
\end{exoresolu}

\courssub{Skill 5: Trigonometric Separable Equations}

\begin{methode}[Separable Equations with Trigonometric Functions]
\begin{enumerate}
\item Separate as usual; typical integrals are
\[ \int \sin x\,dx = -\cos x, \qquad \int \sec^2 x\,dx = \tan x, \qquad \int \tan x\,dx = -\ln|\cos x| = \ln|\sec x|. \]
\item GCE questions usually arrange for $\displaystyle\int \cot y\,dy = \ln|\sin y|$ or $\displaystyle\int \tan y\,dy = -\ln|\cos y|$.
\item Remember $\displaystyle\int \frac{f'(y)}{f(y)}\,dy = \ln|f(y)|$ — the ``derivative over function'' pattern is the key reflex.
\end{enumerate}
\end{methode}

\begin{exoresolu}[A Trigonometric Equation]
Find the general solution of
\[ \frac{dy}{dx} = e^{x}\cos^2 y, \qquad \cos y \neq 0. \]
\tcblower
\textbf{Solution.}\par
\textbf{Step 1: separability.} $e^{x}\cos^2 y = e^{x} \cdot \cos^2 y$: product of a function of $x$ and a function of $y$. Separable.\par
\textbf{Step 2: separate.} Divide by $\cos^2 y$ (valid since $\cos y \neq 0$) and multiply by $dx$:
\[ \frac{dy}{\cos^2 y} = e^{x}\,dx \quad\Longleftrightarrow\quad \sec^2 y\,dy = e^{x}\,dx. \]
\textbf{Step 3: integrate.} Since $\displaystyle\int \sec^2 y\,dy = \tan y$ and $\displaystyle\int e^{x}\,dx = e^{x}$:
\[ \tan y = e^{x} + C. \]
\textbf{Step 4: make $y$ the subject.}
\[ y = \tan^{-1}\!\bigl(e^{x} + C\bigr). \]
\textbf{Conclusion.} The general solution is $y = \arctan\bigl(e^{x} + C\bigr)$, $C \in \mathbb{R}$.\par
\textbf{Check:} differentiate: $\dfrac{dy}{dx} = \dfrac{e^{x}}{1 + (e^{x}+C)^2}$. Also $\cos^2 y = \dfrac{1}{1 + \tan^2 y} = \dfrac{1}{1 + (e^{x}+C)^2}$, so $e^{x}\cos^2 y = \dfrac{e^{x}}{1+(e^{x}+C)^2} = \dfrac{dy}{dx}$. $\checkmark$
\end{exoresolu}

\begin{exoresolu}[An Equation Leading to $\ln|\sin y|$]
Find the general solution of $\dfrac{dy}{dx} = \dfrac{\cos y}{\sin y}\,x^2$, where $\sin y \neq 0$.
\tcblower
\textbf{Solution.}\par
\textbf{Step 1: separate.} The right-hand side is $x^2 \cdot \cot y$: separable. Divide by $\cot y$:
\[ \tan y\,dy = x^2\,dx. \]
\textbf{Step 2: integrate.} Using $\displaystyle\int \tan y\,dy = -\ln|\cos y|$:
\[ -\ln|\cos y| = \frac{x^3}{3} + C_1. \]
\textbf{Step 3: simplify.} Multiply by $-1$ and write $C = -C_1$:
\[ \ln|\cos y| = C - \frac{x^3}{3}. \]
Exponentiating, $|\cos y| = e^{C}e^{-x^3/3}$, so the general solution can be written
\[ \cos y = A e^{-x^3/3}, \qquad A \neq 0, \]
or equivalently $y = \arccos\bigl(Ae^{-x^3/3}\bigr)$ where defined.\par
\textbf{Check:} differentiating $\cos y = Ae^{-x^3/3}$ implicitly gives $-\sin y\,\dfrac{dy}{dx} = -x^2 Ae^{-x^3/3} = -x^2\cos y$, hence $\dfrac{dy}{dx} = x^2\dfrac{\cos y}{\sin y}$. $\checkmark$
\end{exoresolu}

\courssub{Skill 6: Finding Particular Solutions from Initial Conditions}

The general solution contains an arbitrary constant and therefore describes a whole \emph{family} of curves. An \cle{initial condition} $y(x_0) = y_0$ picks out the single member of the family passing through the point $(x_0, y_0)$; the result is called a \cle{particular solution}.

\begin{methode}[From General to Particular Solution]
\begin{enumerate}
\item Solve the equation to obtain the general solution (with constant $C$ or $A$).
\item \textbf{Substitute} $x = x_0$ and $y = y_0$ into the general solution.
\item Solve the resulting equation for the constant.
\item Rewrite the solution with this value: this is the \cle{particular solution}.
\item State the domain of validity if relevant (e.g.\ where a square root or logarithm is defined).
\end{enumerate}
\textbf{Reflex:} if the general solution is $y = A e^{f(x)}$ and $y(x_0) = y_0$, then $A = y_0 e^{-f(x_0)}$.
\end{methode}

\begin{exoresolu}[Particular Solution (GCE Style)]
\begin{enumerate}
\item Find the general solution of $\dfrac{dy}{dx} = \dfrac{y}{x}$, for $x > 0$, $y > 0$.
\item Hence find the particular solution for which $y = 6$ when $x = 2$.
\end{enumerate}
\tcblower
\textbf{Solution.}\par
\textbf{(a) General solution.} The equation is separable: $\dfrac{y}{x} = \dfrac{1}{x}\cdot y$. Separate:
\[ \frac{dy}{y} = \frac{dx}{x}. \]
Integrate both sides:
\[ \ln y = \ln x + C \qquad (\text{valid since } x>0,\ y>0). \]
Exponentiate:
\[ y = e^{\ln x + C} = e^{C}\,x. \]
Writing $A = e^{C} > 0$:
\[ y = Ax. \]
(The solution curves are straight lines through the origin.)\par
\textbf{(b) Particular solution.} Substitute $x = 2$, $y = 6$ into $y = Ax$:
\[ 6 = 2A \quad\Longrightarrow\quad A = 3. \]
\textbf{Conclusion.} The particular solution is
\[ y = 3x, \qquad x > 0. \]
\textbf{Check:} $\dfrac{dy}{dx} = 3$ and $\dfrac{y}{x} = \dfrac{3x}{x} = 3$. $\checkmark$ Also $y(2) = 6$. $\checkmark$
\end{exoresolu}

\begin{exoresolu}[Particular Solution with a Square Root]
Find the particular solution of $\dfrac{dy}{dx} = \dfrac{x}{y}$ ($y > 0$) for which $y = 5$ when $x = 0$, and state the largest interval on which it is defined.
\tcblower
\textbf{Solution.}\par
\textbf{Step 1: general solution.} From Skill 3, $y^2 = x^2 + C$, and since $y > 0$, $y = \sqrt{x^2 + C}$.\par
\textbf{Step 2: apply the condition.} $y(0) = 5$ gives $\sqrt{C} = 5$, so $C = 25$.\par
\textbf{Step 3: particular solution.}
\[ y = \sqrt{x^2 + 25}. \]
\textbf{Step 4: domain.} Since $x^2 + 25 > 0$ for every real $x$, the solution is defined on the whole of $\mathbb{R}$.\par
\textbf{Check:} $\dfrac{dy}{dx} = \dfrac{x}{\sqrt{x^2+25}} = \dfrac{x}{y}$. $\checkmark$
\end{exoresolu}

\courssub{Skill 7: Solving Application Problems (Modelling with Separable Equations)}

Many natural and economic quantities change at a rate that depends on the quantity itself: populations, radioactive masses, temperatures, the value of an investment. Such situations are modelled by separable equations, which is why this chapter matters far beyond the examination hall.

\begin{methode}[Modelling with Separable Differential Equations]
\begin{enumerate}
\item \textbf{Translate} the situation into an equation: ``rate of change proportional to the quantity present'' gives $\dfrac{dQ}{dt} = kQ$; ``proportional to the difference from a limiting value $L$'' gives $\dfrac{dQ}{dt} = k(L - Q)$.
\item Identify the \textbf{variables} and the \textbf{initial condition} $Q(0) = Q_0$.
\item Separate, integrate, find the general solution.
\item Use the initial condition to fix the constant; use any extra data point to fix $k$.
\item \textbf{Answer the question asked} (a time, a quantity, a percentage), giving units and sensible rounding.
\end{enumerate}
\textbf{Key models:} exponential growth/decay $Q = Q_0 e^{kt}$; Newton's law of cooling $T - T_{\text{env}} = (T_0 - T_{\text{env}})e^{-kt}$.
\end{methode}

\begin{exoresolu}[Radioactive Decay and Newton Cooling]
\textbf{(a)} A radioactive substance decays so that the mass $m$ (in grams) satisfies $\dfrac{dm}{dt} = -km$. Initially the mass is $80$ g, and after $5$ days it is $60$ g. Find the mass remaining after $10$ days.\par
\textbf{(b)} A bottle at $90\,^{\circ}$C is placed in a room at $30\,^{\circ}$C. Its temperature $T$ satisfies $\dfrac{dT}{dt} = -k(T - 30)$. After $10$ minutes the temperature is $60\,^{\circ}$C. Find $k$ and the temperature after $20$ minutes.
\tcblower
\textbf{Solution.}\par
\textbf{(a)} \emph{General solution.} Separate: $\dfrac{dm}{m} = -k\,dt$. Integrate: $\ln m = -kt + C$, so $m = Ae^{-kt}$.\par
\emph{Initial condition.} $m(0) = 80 \Rightarrow A = 80$, hence $m = 80e^{-kt}$.\par
\emph{Find $k$.} $m(5) = 60$: $\;60 = 80e^{-5k} \Rightarrow e^{-5k} = \dfrac{3}{4} \Rightarrow -5k = \ln\dfrac{3}{4}$, so $k = \dfrac{1}{5}\ln\dfrac{4}{3} \approx 0.0575$ per day.\par
\emph{Mass after 10 days.} Note $e^{-10k} = \bigl(e^{-5k}\bigr)^2 = \bigl(\tfrac34\bigr)^2 = \tfrac{9}{16}$:
\[ m(10) = 80 \times \frac{9}{16} = 45\ \text{g}. \]
\textbf{(b)} \emph{General solution.} Let $u = T - 30$; then $\dfrac{du}{dt} = -ku$, so $u = Ae^{-kt}$, i.e.
\[ T = 30 + Ae^{-kt}. \]
\emph{Initial condition.} $T(0) = 90 \Rightarrow A = 60$, hence $T = 30 + 60e^{-kt}$.\par
\emph{Find $k$.} $T(10) = 60$: $\;60 = 30 + 60e^{-10k} \Rightarrow e^{-10k} = \dfrac{30}{60} = \dfrac12$, so
\[ k = \frac{\ln 2}{10} \approx 0.0693\ \text{min}^{-1}. \]
\emph{Temperature after 20 minutes.} $e^{-20k} = \bigl(e^{-10k}\bigr)^2 = \tfrac14$:
\[ T(20) = 30 + 60 \times \frac{1}{4} = 30 + 15 = 45\,^{\circ}\mathrm{C}. \]
\textbf{Interpretation:} the temperature gap halves every $10$ minutes — the hallmark of exponential decay.
\end{exoresolu}

\begin{savaistu}[Why ``Half-Life'' is Constant]
In part (a) above, the mass fell from $80$ g to $60$ g in $5$ days, and from $60$ g to $45$ g in the next $5$ days: the \emph{ratio} $\tfrac34$ is the same over each equal time interval. This is the defining feature of exponential decay $m = m_0 e^{-kt}$: the time needed to halve the quantity (the \cle{half-life}) is always $T_{1/2} = \dfrac{\ln 2}{k}$, independent of the starting amount. This principle underlies radiocarbon dating and the management of radioactive waste.
\end{savaistu}

\courssub{Skill 8: Full GCE-Style Structured Question}

\begin{methode}[Tackling a Structured Examination Question]
\begin{enumerate}
\item Read the \textbf{whole} question first; part (a) usually feeds into part (b).
\item Show \textbf{every} line of the separation and integration — method marks are awarded for $\int \frac{dy}{h(y)} = \int g(x)\,dx$ even if a later slip occurs.
\item Keep the constant symbolic until the initial condition is applied.
\item In the final part, substitute numbers \emph{after} obtaining the exact formula, and give the final answer to the accuracy requested.
\end{enumerate}
\end{methode}

\begin{exoresolu}[Comprehensive GCE-Style Exercise]
The number of bacteria $N$ in a culture grows according to
\[ \frac{dN}{dt} = \frac{N}{2}\left(1 + \frac{1}{t+1}\right), \qquad t \geq 0, \]
where $t$ is measured in hours. Initially $N = 1000$.
\begin{enumerate}
\item Show that $\dfrac{1}{N}\,dN = \dfrac{t+2}{2(t+1)}\,dt$.
\item Hence show that $\ln N = \dfrac{t}{2} + \dfrac{1}{2}\ln(t+1) + C$ for some constant $C$.
\item Find the particular solution for $N$ in terms of $t$.
\item Calculate, to the nearest whole number, the value of $N$ when $t = 3$.
\end{enumerate}
\tcblower
\textbf{Solution.}\par
\textbf{(a)} The equation is separable. Divide by $N$ and multiply by $dt$:
\[ \frac{dN}{N} = \frac{1}{2}\left(1 + \frac{1}{t+1}\right)dt = \frac{1}{2}\cdot\frac{(t+1)+1}{t+1}\,dt = \frac{t+2}{2(t+1)}\,dt. \quad\blacksquare \]
\textbf{(b)} Write $\dfrac{t+2}{t+1} = \dfrac{(t+1)+1}{t+1} = 1 + \dfrac{1}{t+1}$. Integrate both sides:
\[ \int \frac{dN}{N} = \frac{1}{2}\int\left(1 + \frac{1}{t+1}\right)dt. \]
\[ \ln N = \frac{1}{2}\bigl(t + \ln(t+1)\bigr) + C = \frac{t}{2} + \frac{1}{2}\ln(t+1) + C. \quad\blacksquare \]
(Here $N > 0$ and $t + 1 > 0$, so no absolute-value signs are needed.)\par
\textbf{(c)} Apply $N = 1000$ when $t = 0$:
\[ \ln 1000 = 0 + \tfrac12\ln 1 + C \quad\Longrightarrow\quad C = \ln 1000. \]
Therefore
\[ \ln N = \frac{t}{2} + \frac{1}{2}\ln(t+1) + \ln 1000. \]
Exponentiating and using $e^{\frac12\ln(t+1)} = \sqrt{t+1}$:
\[ N = 1000\,e^{t/2}\sqrt{t+1}. \]
\textbf{(d)} At $t = 3$:
\[ N = 1000\,e^{1.5}\sqrt{4} = 2000\,e^{1.5} \approx 2000 \times 4.4817 \approx 8963. \]
\textbf{Conclusion:} after $3$ hours there are approximately $\mathbf{8963}$ bacteria.
\end{exoresolu}

\begin{attention}[Frequent Errors at the GCE]
\begin{itemize}
\item Dividing by $h(y)$ without checking whether $h(y) = 0$ gives extra (constant) solutions.
\item Writing $\ln y$ when $y$ could be negative: use $\ln|y|$.
\item Adding two constants $C_1$ and $C_2$ and forgetting to merge them.
\item Applying the initial condition \emph{before} integrating — the condition is used \emph{after} the general solution is found.
\item In applications, mixing up the proportionality constant sign: decay requires $-k$ with $k > 0$.
\end{itemize}
\end{attention}

\begin{exercice}[Practice — Core Skills]
\begin{enumerate}
\item Verify that $y = 2e^{-3x}$ satisfies $\dfrac{dy}{dx} + 3y = 0$ with $y(0) = 2$.
\item State which of the following are separable, giving $g(x)$ and $h(y)$ where possible:\par
(a)\; $\dfrac{dy}{dx} = x^2 y$ \qquad (b)\; $\dfrac{dy}{dx} = x + xy$ \qquad (c)\; $\dfrac{dy}{dx} = \ln(x + y)$ \qquad (d)\; $\dfrac{dy}{dx} = \dfrac{y}{x^2 + 1}$.
\item Find the general solution of $\dfrac{dy}{dx} = \dfrac{x^2}{y^2}$, $y \neq 0$.
\item Find the general solution of $\dfrac{dy}{dx} = 3x^2 y$, and the particular solution with $y(0) = 4$.
\item A quantity $P$ grows according to $\dfrac{dP}{dt} = 0.2\,P$ with $P(0) = 500$. Find $P$ when $t = 5$, giving your answer to the nearest whole number.
\end{enumerate}
\end{exercice}

\begin{corrige}[Practice — Core Skills]
\textbf{1.} $y' = -6e^{-3x}$, so $y' + 3y = -6e^{-3x} + 6e^{-3x} = 0$ for all $x$; and $y(0) = 2e^{0} = 2$. Both conditions hold. $\checkmark$\par
\textbf{2.} (a) Separable: $g(x) = x^2$, $h(y) = y$. (b) Separable: $x + xy = x(1 + y)$, so $g(x) = x$, $h(y) = 1 + y$. (c) Not separable (a sum inside the logarithm cannot be split into a product). (d) Separable: $g(x) = \dfrac{1}{x^2+1}$, $h(y) = y$.\par
\textbf{3.} Separate: $y^2\,dy = x^2\,dx$. Integrate: $\dfrac{y^3}{3} = \dfrac{x^3}{3} + C_1$, so $y^3 = x^3 + C$ and $y = \sqrt[3]{x^3 + C}$.\par
\textbf{4.} Separate: $\dfrac{dy}{y} = 3x^2\,dx$ (noting $y = 0$ is a constant solution). Integrate: $\ln|y| = x^3 + C$, so $y = Ae^{x^3}$. Condition $y(0) = 4$: $A = 4$. Particular solution: $y = 4e^{x^3}$.\par
\textbf{5.} Model solution: $P = 500e^{0.2t}$. At $t = 5$: $P = 500e^{1} \approx 500 \times 2.71828 \approx 1359$.
\end{corrige}

\begin{exercice}[GCE-Style Consolidation]
A tank contains water leaking out so that the volume $V$ (in litres) satisfies
\[ \frac{dV}{dt} = -\frac{V}{t + 4}, \qquad t \geq 0, \]
where $t$ is in minutes. Initially $V = 200$ litres.
\begin{enumerate}
\item Show that $\ln V = C - \ln(t + 4)$ for some constant $C$.
\item Hence show that $V = \dfrac{800}{t + 4}$.
\item Find the volume after $6$ minutes, and the time at which the volume falls to $40$ litres.
\end{enumerate}
\end{exercice}

\begin{corrige}[GCE-Style Consolidation]
\textbf{(a)} Separate: $\dfrac{dV}{V} = -\dfrac{dt}{t+4}$. Integrate: $\ln V = -\ln(t+4) + C$ (valid since $V > 0$, $t + 4 > 0$). $\blacksquare$\par
\textbf{(b)} Exponentiate: $V = e^{C}\,(t+4)^{-1} = \dfrac{A}{t+4}$ with $A = e^{C}$. Condition $V(0) = 200$: $200 = \dfrac{A}{4}$, so $A = 800$ and $V = \dfrac{800}{t+4}$. $\blacksquare$\par
\textbf{(c)} At $t = 6$: $V = \dfrac{800}{10} = 80$ litres. For $V = 40$: $40 = \dfrac{800}{t+4} \Rightarrow t + 4 = 20 \Rightarrow t = 16$ minutes.
\end{corrige}

\newpage

\coursec{Exercises}

\courssub{I check my knowledge}

\begin{exercice}[Multiple choice: identifying a separable equation]
Only one of the following differential equations is separable. Which one? Justify your answer by writing the right-hand side as a product (or quotient) of a function of $x$ alone and a function of $y$ alone.
\begin{enumerate}
\item $\dfrac{dy}{dx}=x+y$
\item $\dfrac{dy}{dx}=e^{x+y}$
\item $\dfrac{dy}{dx}=\sin(x+y)$
\item $\dfrac{dy}{dx}=xy+x+y+1$
\end{enumerate}
\end{exercice}

\begin{exercice}[True or false?]
State whether each assertion is true or false, justifying your answer carefully.
\begin{enumerate}
\item Every first-order differential equation of the form $\dfrac{dy}{dx}=f(x)\,g(y)$ is separable.
\item The equation $\dfrac{dy}{dx}=3y$ has general solution $y=Ce^{3x}$, where $C$ is an arbitrary constant.
\item The function $y=2e^{x^{2}}$ is a solution of the initial-value problem $\dfrac{dy}{dx}=2xy$, $y(0)=2$.
\item When separating variables in $\dfrac{dy}{dx}=y^{2}$, dividing by $y^{2}$ may cause us to lose the constant solution $y=0$.
\end{enumerate}
\end{exercice}

\begin{exercice}[Definitions]
\begin{enumerate}
\item Define a \cle{first-order differential equation} and give one example.
\item Explain what is meant by the \cle{general solution} and by a \cle{particular solution} of a first-order differential equation.
\item What is an \cle{initial condition}, and why does a first-order equation need exactly one initial condition to determine a unique particular solution?
\end{enumerate}
\end{exercice}

\begin{exercice}[Direct calculations]
Find the general solution of each of the following separable equations.
\begin{enumerate}
\item $\dfrac{dy}{dx}=\dfrac{x}{y}$
\item $\dfrac{dy}{dx}=e^{2x}\,y$
\item $\dfrac{dy}{dx}=\dfrac{1}{x^{2}y}$, \quad $x>0$, $y>0$
\item $x\,\dfrac{dy}{dx}=y+1$, \quad $y>-1$
\end{enumerate}
\end{exercice}

\courssub{I practise}

\begin{exercice}[An initial-value problem with an explicit solution]
Solve the initial-value problem
\[
\frac{dy}{dx}=\frac{3x^{2}}{2y},\qquad y(1)=2.
\]
Give the \emph{explicit} solution (i.e. $y$ expressed in terms of $x$), state its domain of validity, and explain how the initial condition determines the sign of the square root.
\end{exercice}

\begin{exercice}[Exponential growth of a bacterial culture]
A bacterial culture in a laboratory in Yaoundé grows according to the law
\[
\frac{dN}{dt}=0.05\,N,
\]
where $N$ is the number of bacteria (in thousands) and $t$ is the time in hours. Given that $N(0)=500$:
\begin{enumerate}
\item Find $N(t)$.
\item Calculate, to the nearest hour, the time required for the population to triple.
\item Show that this tripling time does not depend on the initial population.
\end{enumerate}
\end{exercice}

\begin{exercice}[Trigonometric and logarithmic separation]
\begin{enumerate}
\item Find the general solution of $\dfrac{dy}{dx}=y\cos x$ and sketch, on the same axes, the solution curves passing through $(0,1)$, $(0,-1)$ and $(0,0)$.
\item Solve the initial-value problem $\dfrac{dy}{dx}=\dfrac{y\ln x}{x}$, $y(1)=3$, for $x>0$.
\end{enumerate}
\end{exercice}

\begin{exercice}[Newton's law of cooling]
A cup of tea at $90^{\circ}$C is placed in a room kept at $20^{\circ}$C. According to Newton's law of cooling,
\[
\frac{dT}{dt}=-k(T-20),
\]
where $T$ is the temperature of the tea in degrees Celsius, $t$ the time in minutes and $k>0$ a constant. After $5$ minutes the tea is at $70^{\circ}$C.
\begin{enumerate}
\item Solve the differential equation and determine the value of $k$.
\item Find the temperature of the tea after $15$ minutes, giving your answer to the nearest degree.
\item What happens to $T(t)$ as $t\to+\infty$? Comment.
\end{enumerate}
\end{exercice}

\courssub{I prepare for the GCE}

\begin{exercice}[GCE-style structured question: curve through a point] \hfill (10 marks)
A curve passes through the point $(1,2)$ and is such that the gradient at every point $(x,y)$ of the curve is equal to $\dfrac{2y}{x}$.
\begin{enumerate}
\item[(a)] Write down a differential equation satisfied by $y$ and show that it is separable. \hfill (2 marks)
\item[(b)] Solve the differential equation, giving the general solution in the form $y=f(x)$. \hfill (4 marks)
\item[(c)] Use the given point to find the particular solution, and state the set of values of $x$ for which it is valid. \hfill (2 marks)
\item[(d)] Find the equation of the tangent to the curve at the point where $x=2$. \hfill (2 marks)
\end{enumerate}
\end{exercice}

\begin{exercice}[GCE-style structured question: a mixing problem] \hfill (12 marks)
A tank at a water-treatment station in Douala initially contains $100$ litres of pure water. Brine containing $0.2$ kg of salt per litre flows into the tank at a rate of $3$ litres per minute, and the well-stirred mixture flows out at a rate of $2$ litres per minute. Let $S(t)$ be the amount of salt, in kilograms, in the tank after $t$ minutes.
\begin{enumerate}
\item[(a)] Show that the volume of liquid in the tank after $t$ minutes is $V(t)=100+t$ litres. \hfill (1 mark)
\item[(b)] Explain why
\[
\frac{dS}{dt}=0.6-\frac{2S}{100+t},
\]
and show that this equation is \emph{not} separable as it stands. \hfill (3 marks)
\item[(c)] Now suppose instead that the mixture flows out at $3$ litres per minute (so the volume stays constant at $100$ litres). Write down the new differential equation for $S(t)$, show that it is separable, and solve it given that $S(0)=0$. \hfill (5 marks)
\item[(d)] For the model of part (c), find the limiting amount of salt in the tank as $t\to+\infty$, and interpret the result physically. \hfill (3 marks)
\end{enumerate}
\end{exercice}

\begin{exercice}[GCE-style essay question: beyond separability] \hfill (14 marks)
Consider the differential equation
\[
\frac{dy}{dx}=\frac{x^{2}+y^{2}}{xy},\qquad x>0,\ y>0.
\]
\begin{enumerate}
\item[(a)] Show that this equation is \emph{not} separable, i.e. that its right-hand side cannot be written as a product $f(x)\,g(y)$. \hfill (3 marks)
\item[(b)] Using the substitution $y=vx$, where $v$ is a function of $x$, show that
\[
x\frac{dv}{dx}=\frac{1}{v}.
\]
\hfill (4 marks)
\item[(c)] Solve the equation obtained in (b) and deduce the general solution of the original equation in the form $y^{2}=2x^{2}\ln x+Cx^{2}$. \hfill (4 marks)
\item[(d)] Find the particular solution satisfying $y(1)=2$ and state its domain. \hfill (3 marks)
\end{enumerate}
\emph{(Bonus, 2 marks: verify by direct differentiation that your solution of part (c) satisfies the original equation.)}
\end{exercice}

\newpage

\coursec{Solutions to the exercises}

\begin{corrige}[Exercise 1 — Multiple choice: identifying a separable equation]
The correct answer is \textbf{(b)}.
\begin{enumerate}
\item $\dfrac{dy}{dx}=x+y$: a sum of a function of $x$ and a function of $y$; it cannot be factored into $f(x)g(y)$. Not separable.
\item $\dfrac{dy}{dx}=e^{x+y}=e^{x}\cdot e^{y}$: a product of a function of $x$ alone and a function of $y$ alone. \textbf{Separable.}
\item $\dfrac{dy}{dx}=\sin(x+y)$: the argument $x+y$ cannot be split into a product $f(x)g(y)$ (the identity $\sin(x+y)=\sin x\cos y+\cos x\sin y$ gives a sum, not a product). Not separable.
\item $\dfrac{dy}{dx}=xy+x+y+1$: although this \emph{factors} as $(x+1)(y+1)$, note carefully: $xy+x+y+1=(x+1)(y+1)$, which \emph{is} a product of a function of $x$ and a function of $y$! So (d) is also separable. In a typical multiple-choice setting where only one answer is expected, (b) is the intended answer; however, a careful student should observe that (d) factorises as $(x+1)(y+1)$ and is therefore separable too. In an examination, mention both and justify.
\end{enumerate}
\textbf{Answer: (b)} (with the remark that (d) is also separable since $xy+x+y+1=(x+1)(y+1)$).
\end{corrige}

\begin{corrige}[Exercise 2 — True or false?]
\begin{enumerate}
\item \textbf{True.} By definition, $\dfrac{dy}{dx}=f(x)g(y)$ separates as $\dfrac{dy}{g(y)}=f(x)\,dx$ (for $g(y)\neq 0$), and both sides can be integrated.
\item \textbf{True.} Separating: $\dfrac{dy}{y}=3\,dx$, so $\ln|y|=3x+C_1$, hence $|y|=e^{C_1}e^{3x}$, i.e. $y=Ce^{3x}$ with $C\in\mathbb{R}$ (including $C=0$, the constant solution). Check: $y'=3Ce^{3x}=3y$. \checkmark
\item \textbf{True.} If $y=2e^{x^{2}}$, then $\dfrac{dy}{dx}=2e^{x^{2}}\cdot 2x=2x\bigl(2e^{x^{2}}\bigr)=2xy$, and $y(0)=2e^{0}=2$. Both conditions hold. \checkmark
\item \textbf{True.} Separating $\dfrac{dy}{dx}=y^{2}$ gives $\dfrac{dy}{y^{2}}=dx$, valid only for $y\neq 0$. Yet $y=0$ clearly satisfies the equation ($0'=0=0^{2}$). This constant solution is lost when dividing by $y^{2}$; it must be checked and quoted separately.
\end{enumerate}
\end{corrige}

\begin{corrige}[Exercise 3 — Definitions]
\begin{enumerate}
\item A \cle{first-order differential equation} is an equation involving an unknown function $y$ of one variable $x$, its derivative $\dfrac{dy}{dx}$, and possibly $x$ itself; it contains no derivative of order higher than one. General form: $F\!\left(x,y,\dfrac{dy}{dx}\right)=0$. Example: $\dfrac{dy}{dx}=3x^{2}y$.
\item The \cle{general solution} is the family of all solutions, expressed with one arbitrary constant $C$ (one constant because the equation is of first order, so one integration is needed). A \cle{particular solution} is obtained by assigning a specific value to $C$, usually determined by an initial condition. Example: $y'=2x$ has general solution $y=x^{2}+C$; with $y(0)=3$, the particular solution is $y=x^{2}+3$.
\item An \cle{initial condition} is a requirement of the form $y(x_0)=y_0$, specifying the value of the solution at one point. A first-order equation involves exactly one integration, producing exactly one arbitrary constant; one initial condition gives exactly one equation to determine that constant, hence a \emph{unique} particular solution (geometrically: one point selects one curve from the family of solution curves).
\end{enumerate}
\end{corrige}

\begin{corrige}[Exercise 4 — Direct calculations]
\begin{enumerate}
\item $\dfrac{dy}{dx}=\dfrac{x}{y}$. Separate: $y\,dy=x\,dx$. Integrate: $\dfrac{y^{2}}{2}=\dfrac{x^{2}}{2}+C_1$, i.e.
\[ y^{2}=x^{2}+C \qquad (C\in\mathbb{R}). \]
(Solution curves are hyperbolas; for $C=0$, the lines $y=\pm x$.)
\item $\dfrac{dy}{dx}=e^{2x}y$. Separate ($y\neq 0$): $\dfrac{dy}{y}=e^{2x}\,dx$. Integrate: $\ln|y|=\dfrac{1}{2}e^{2x}+C_1$, hence
\[ y=Ce^{\frac{1}{2}e^{2x}} \qquad (C\in\mathbb{R},\ C=0\text{ included}). \]
\item $\dfrac{dy}{dx}=\dfrac{1}{x^{2}y}$ with $x>0$, $y>0$. Separate: $y\,dy=\dfrac{dx}{x^{2}}$. Integrate: $\dfrac{y^{2}}{2}=-\dfrac{1}{x}+C_1$, hence
\[ y^{2}=C-\frac{2}{x},\qquad\text{i.e. } y=\sqrt{C-\frac{2}{x}} \quad (\text{positive root since } y>0). \]
\item $x\dfrac{dy}{dx}=y+1$ with $y>-1$. Separate: $\dfrac{dy}{y+1}=\dfrac{dx}{x}$. Integrate: $\ln(y+1)=\ln|x|+C_1$, so $y+1=Ax$ with $A>0$ if $x>0$; in general the family of solutions is
\[ y=Cx-1 \qquad (C\in\mathbb{R}). \]
The constant solution $y=-1$ corresponds to $C=0$. The condition $y>-1$ excludes $C=0$ and restricts the domain according to the sign of $Cx$. Check: $x\dfrac{dy}{dx}=xC=y+1$. \checkmark
\end{enumerate}
\end{corrige}

\begin{corrige}[Exercise 5 — An initial-value problem with an explicit solution]
\textbf{Separation.} $\dfrac{dy}{dx}=\dfrac{3x^{2}}{2y}$ gives $2y\,dy=3x^{2}\,dx$.

\textbf{Integration.} $y^{2}=x^{3}+C$.

\textbf{Initial condition.} $y(1)=2$: $4=1+C$, so $C=3$. Hence $y^{2}=x^{3}+3$.

\textbf{Sign of the root.} $y=\pm\sqrt{x^{3}+3}$. Since $y(1)=2>0$, we must take the \emph{positive} root:
\[ \boxed{\,y=\sqrt{x^{3}+3}\,}. \]

\textbf{Domain of validity.} We need $x^{3}+3>0$ (strictly, since $y$ appears in the denominator of the equation, $y\neq 0$), i.e. $x>-\sqrt[3]{3}$. The solution is valid on $\left]-\sqrt[3]{3},\ +\infty\right[$, the largest interval containing $x_0=1$ on which the solution is defined and differentiable.
\end{corrige}

\begin{corrige}[Exercise 6 — Exponential growth of a bacterial culture]
\begin{enumerate}
\item Separate: $\dfrac{dN}{N}=0.05\,dt$, so $\ln N=0.05t+C_1$ ($N>0$), i.e. $N=Ae^{0.05t}$. With $N(0)=500$: $A=500$. Hence
\[ \boxed{N(t)=500\,e^{0.05t}} \quad (\text{in thousands}). \]
\item Tripling: $N(t)=3N(0)$, i.e. $e^{0.05t}=3$, so $t=\dfrac{\ln 3}{0.05}=20\ln 3\approx 20\times 1.0986\approx 21.97$. To the nearest hour, the population triples after $\boxed{22\text{ hours}}$.
\item In general, with $N(0)=N_0$: $N(t)=N_0e^{0.05t}$. Tripling means $N_0e^{0.05t}=3N_0$; dividing by $N_0$ gives $e^{0.05t}=3$, so $t=20\ln 3$ \emph{independently of $N_0$}. This is a characteristic feature of exponential growth: the time to multiply by any fixed factor is constant.
\end{enumerate}
\end{corrige}

\begin{corrige}[Exercise 7 — Trigonometric and logarithmic separation]
\begin{enumerate}
\item $\dfrac{dy}{dx}=y\cos x$. Separate ($y\neq 0$): $\dfrac{dy}{y}=\cos x\,dx$, so $\ln|y|=\sin x+C_1$, hence
\[ y=Ce^{\sin x} \qquad (C\in\mathbb{R}). \]
Through $(0,1)$: $1=Ce^{0}\Rightarrow C=1$, curve $y=e^{\sin x}$.
Through $(0,-1)$: $C=-1$, curve $y=-e^{\sin x}$.
Through $(0,0)$: $C=0$, the constant solution $y=0$ (the $x$-axis).
\begin{popfigure}
\begin{center}
\begin{tikzpicture}
\begin{axis}[width=11cm, height=6.5cm, axis lines=middle, xlabel=$x$, ylabel=$y$,
xmin=-7, xmax=7, ymin=-3, ymax=3, domain=-6.5:6.5, samples=200,
xtick={-6.28,-3.14,0,3.14,6.28}, xticklabels={$-2\pi$,$-\pi$,$0$,$\pi$,$2\pi$}]
\addplot[popblue, thick] {exp(sin(deg(x)))};
\addplot[poporange, thick] {-exp(sin(deg(x)))};
\addplot[popgreen, thick] {0};
\end{axis}
\end{tikzpicture}
\end{center}
\legende{Curves $y=e^{\sin x}$ (blue), $y=-e^{\sin x}$ (orange), $y=0$ (green)}
\end{popfigure}
Note that solution curves never cross: $e^{\sin x}>0$ always, so the sign of $y$ is fixed by the initial condition.
\item $\dfrac{dy}{dx}=\dfrac{y\ln x}{x}$, $x>0$. Separate: $\dfrac{dy}{y}=\dfrac{\ln x}{x}\,dx$. Since $\displaystyle\int\frac{\ln x}{x}\,dx=\frac{(\ln x)^{2}}{2}$ (substitution $u=\ln x$), we get $\ln|y|=\dfrac{(\ln x)^{2}}{2}+C_1$, i.e. $y=Ce^{(\ln x)^{2}/2}$. With $y(1)=3$: $\ln 1=0$, so $C=3$. Hence
\[ \boxed{y=3e^{\frac{(\ln x)^{2}}{2}}},\qquad x>0. \]
\end{enumerate}
\end{corrige}

\begin{corrige}[Exercise 8 — Newton's law of cooling]
\begin{enumerate}
\item Let $\theta=T-20$. Then $\dfrac{d\theta}{dt}=\dfrac{dT}{dt}=-k\theta$. Separating: $\dfrac{d\theta}{\theta}=-k\,dt$, so $\ln\theta=-kt+C_1$, i.e. $\theta=Ae^{-kt}$. Hence
\[ T(t)=20+Ae^{-kt}. \]
$T(0)=90$: $A=70$. So $T(t)=20+70e^{-kt}$.
$T(5)=70$: $70=20+70e^{-5k}$, so $e^{-5k}=\dfrac{50}{70}=\dfrac{5}{7}$, giving
\[ k=\frac{1}{5}\ln\frac{7}{5}\approx\frac{0.3365}{5}\approx 0.0673\ \mathrm{min^{-1}}. \]
\item $T(15)=20+70e^{-15k}=20+70\bigl(e^{-5k}\bigr)^{3}=20+70\left(\dfrac{5}{7}\right)^{3}=20+70\times\dfrac{125}{343}\approx 20+25.5\approx \boxed{46^{\circ}\text{C}}$.
\item As $t\to+\infty$, $e^{-kt}\to 0$ (since $k>0$), so $T(t)\to 20^{\circ}$C. The tea tends to the room temperature, which is exactly what experience predicts: cooling never goes below the ambient temperature in this model.
\end{enumerate}
\end{corrige}

\begin{corrige}[Exercise 9 — GCE-style: curve through a point]
\textbf{(a)} The gradient at $(x,y)$ is $\dfrac{dy}{dx}$, so
\[ \frac{dy}{dx}=\frac{2y}{x}= \underbrace{\frac{2}{x}}_{f(x)}\cdot\underbrace{y}_{g(y)}, \]
a product of a function of $x$ alone and a function of $y$ alone: the equation is separable. \hfill \emph{[M1: equation; A1: justification]}

\textbf{(b)} For $y\neq 0$: $\dfrac{dy}{y}=\dfrac{2}{x}\,dx$. Integrating: $\ln|y|=2\ln|x|+C_1=\ln x^{2}+C_1$, so $|y|=e^{C_1}x^{2}$, i.e.
\[ y=Cx^{2}\qquad (C\in\mathbb{R},\ C=0\text{ included}). \]
\emph{[M1: separation; M1: correct integrals; A1: $\ln$ form; A1: explicit form $y=Cx^{2}$]}

\textbf{(c)} Through $(1,2)$: $2=C\cdot 1$, so $C=2$ and $\boxed{y=2x^{2}}$. Since the original equation requires $x\neq 0$ and the curve passes through $x=1$, the solution is valid for $x>0$. \hfill \emph{[M1: substitution of point; A1: $y=2x^{2}$ with domain]}

\textbf{(d)} At $x=2$: $y=2(2)^{2}=8$. Gradient: $\dfrac{dy}{dx}=\dfrac{2y}{x}=\dfrac{16}{2}=8$. Tangent:
\[ y-8=8(x-2)\quad\Longrightarrow\quad \boxed{y=8x-8}. \]
\emph{[M1: point and gradient; A1: equation of tangent]}

\textbf{Examiner's expectations:} clear separation with $dy$ and $dx$ on opposite sides; constant of integration shown; domain stated (a mark is often lost for omitting $x>0$).
\end{corrige}

\begin{corrige}[Exercise 10 — GCE-style: a mixing problem]
\textbf{(a)} Volume in per minute: $3$ L; out: $2$ L; net gain $1$ L/min. Starting from $100$ L: $V(t)=100+t$ litres. \hfill \emph{[B1]}

\textbf{(b)} Salt enters at rate $0.2\times 3=0.6$ kg/min. The mixture leaves at $2$ L/min with concentration $\dfrac{S}{V}=\dfrac{S}{100+t}$ kg/L, so salt leaves at rate $\dfrac{2S}{100+t}$ kg/min. Therefore
\[ \frac{dS}{dt}=0.6-\frac{2S}{100+t}. \]
The right-hand side is a \emph{difference} of a constant and a term involving both $t$ and $S$; it cannot be written as a product $f(t)\,g(S)$, so the equation is \textbf{not separable} as it stands. \hfill \emph{[M1: inflow rate; M1: outflow rate; A1: non-separability justified]}

\textbf{(c)} Outflow now $3$ L/min, volume constant at $100$ L. Salt leaves at $3\times\dfrac{S}{100}=\dfrac{3S}{100}$ kg/min, so
\[ \frac{dS}{dt}=0.6-\frac{3S}{100}=\frac{60-3S}{100}=\underbrace{\frac{3}{100}}_{f(t)\text{ (constant)}}\cdot\underbrace{(20-S)}_{g(S)}, \]
which \emph{is} separable. Separating (for $S\neq 20$):
\[ \frac{dS}{20-S}=\frac{3}{100}\,dt \;\Longrightarrow\; -\ln|20-S|=\frac{3t}{100}+C_1 \;\Longrightarrow\; 20-S=Ae^{-3t/100}. \]
$S(0)=0$: $A=20$. Hence
\[ \boxed{S(t)=20\bigl(1-e^{-0.03t}\bigr)}. \]
\emph{[M1: new equation; A1: separable shown; M1: separation and integration; M1: exponential form; A1: constant from $S(0)=0$]}

\textbf{(d)} As $t\to+\infty$, $e^{-0.03t}\to 0$, so $S(t)\to 20$ kg. Interpretation: the salt content approaches the equilibrium where outflow removes salt as fast as inflow supplies it; the concentration tends to $\dfrac{20}{100}=$ $0.2\ \mathrm{kg/L}$, exactly the concentration of the incoming brine — the tank's content gradually becomes indistinguishable from the inflow. \hfill \emph{[M1: limit; A1: value 20 kg; A1: physical interpretation]}

\textbf{Examiner's expectations:} units used consistently; clear distinction between the non-separable model (b) and the separable one (c); equilibrium interpreted, not merely computed.
\end{corrige}

\begin{corrige}[Exercise 11 — GCE-style essay: beyond separability]
\textbf{(a)} Suppose $\dfrac{x^{2}+y^{2}}{xy}=f(x)g(y)$ for all $x,y>0$. Note $\dfrac{x^{2}+y^{2}}{xy}=\dfrac{x}{y}+\dfrac{y}{x}$. Setting $x=y=1$ gives $f(1)g(1)=2$; setting $x=1,y=2$ and $x=2,y=1$ both give $\dfrac{5}{2}$, so $f(1)g(2)=f(2)g(1)$. Setting $x=y=2$ gives $f(2)g(2)=2$. Then
\[ \bigl(f(1)g(2)\bigr)\bigl(f(2)g(1)\bigr)=\tfrac{25}{4},\qquad \bigl(f(1)g(1)\bigr)\bigl(f(2)g(2)\bigr)=4, \]
but the two left-hand sides are equal (both equal $f(1)f(2)g(1)g(2)$), while $\tfrac{25}{4}\neq 4$: contradiction. Hence no such factorisation exists and the equation is \textbf{not separable}. \hfill \emph{[M1: assumption; M1: test values; A1: contradiction]}

\textbf{(b)} Let $y=vx$. Then $\dfrac{dy}{dx}=v+x\dfrac{dv}{dx}$. Also
\[ \frac{x^{2}+y^{2}}{xy}=\frac{x^{2}+v^{2}x^{2}}{x\cdot vx}=\frac{1+v^{2}}{v}. \]
Therefore
\[ v+x\frac{dv}{dx}=\frac{1+v^{2}}{v}=v+\frac{1}{v}\quad\Longrightarrow\quad \boxed{x\frac{dv}{dx}=\frac{1}{v}}. \]
\emph{[M1: derivative of $vx$; M1: substitution; A1: simplification of RHS; A1: result]}

\textbf{(c)} The new equation is separable: $v\,dv=\dfrac{dx}{x}$. Integrating:
\[ \frac{v^{2}}{2}=\ln x+C_1 \quad (x>0), \qquad\text{so } v^{2}=2\ln x+C. \]
Since $v=\dfrac{y}{x}$, $\dfrac{y^{2}}{x^{2}}=2\ln x+C$, hence
\[ \boxed{y^{2}=2x^{2}\ln x+Cx^{2}}. \]
\emph{[M1: separation; M1: integration; A1: $v^{2}=2\ln x+C$; A1: back-substitution]}

\textbf{(d)} $y(1)=2$: $4=0+C$, so $C=4$ and
\[ y^{2}=2x^{2}\ln x+4x^{2}=2x^{2}(\ln x+2),\qquad y=x\sqrt{2(\ln x+2)} \]
(positive root since $y>0$). Domain: $\ln x+2>0$, i.e. $x>e^{-2}$; the particular solution is valid on $\left]e^{-2},+\infty\right[$. \hfill \emph{[M1: constant; A1: explicit $y$; A1: domain]}

\textbf{Bonus verification.} From $y^{2}=2x^{2}\ln x+Cx^{2}$, differentiate: $2y\dfrac{dy}{dx}=4x\ln x+2x+2Cx$. Divide by $2y$:
\[ \frac{dy}{dx}=\frac{2x\ln x+x+Cx}{y}=\frac{x\bigl(2\ln x+1+C\bigr)}{y}. \]
Now $\dfrac{x^{2}+y^{2}}{xy}=\dfrac{x^{2}+2x^{2}\ln x+Cx^{2}}{xy}=\dfrac{x\bigl(1+2\ln x+C\bigr)}{y}$: identical. \checkmark

\textbf{Examiner's expectations:} rigorous contradiction argument in (a) (a mere "it looks like a sum" earns no marks); correct use of the product rule in (b); the back-substitution $v=y/x$ explicitly written in (c); domain in (d).
\end{corrige}

\newpage

\coursec{Summary sheet}

\begin{popfigure}
\begin{center}\tcbox[colback=popblue,colframe=popblue,coltext=white,arc=9pt,boxsep=4pt,left=6pt,right=6pt]{\bfseries\small First Order Separable DEs}\end{center}
\vspace{-2pt}
\begin{tcbraster}[raster columns=3,raster equal height=rows,raster column skip=6pt,raster row skip=6pt,enhanced,blankest]
\begin{tcolorbox}[enhanced,colback=poporange,colframe=poporange,coltext=white,boxrule=0.9pt,arc=3pt,left=4pt,right=4pt,top=3pt,bottom=3pt,boxsep=1pt,halign=left]\bfseries\scriptsize Meaning \& Classification\end{tcolorbox}
\begin{tcolorbox}[enhanced,colback=popgreen,colframe=popgreen,coltext=white,boxrule=0.9pt,arc=3pt,left=4pt,right=4pt,top=3pt,bottom=3pt,boxsep=1pt,halign=left]\bfseries\scriptsize Recognising Separable Eqs.\end{tcolorbox}
\begin{tcolorbox}[enhanced,colback=poppurple,colframe=poppurple,coltext=white,boxrule=0.9pt,arc=3pt,left=4pt,right=4pt,top=3pt,bottom=3pt,boxsep=1pt,halign=left]\bfseries\scriptsize Separation of Variables\end{tcolorbox}
\begin{tcolorbox}[enhanced,colback=poppink,colframe=poppink,coltext=white,boxrule=0.9pt,arc=3pt,left=4pt,right=4pt,top=3pt,bottom=3pt,boxsep=1pt,halign=left]\bfseries\scriptsize Integration \& General Solution\end{tcolorbox}
\begin{tcolorbox}[enhanced,colback=popgold,colframe=popgold,coltext=white,boxrule=0.9pt,arc=3pt,left=4pt,right=4pt,top=3pt,bottom=3pt,boxsep=1pt,halign=left]\bfseries\scriptsize Particular Solutions (IC)\end{tcolorbox}
\begin{tcolorbox}[enhanced,colback=popdark,colframe=popdark,coltext=white,boxrule=0.9pt,arc=3pt,left=4pt,right=4pt,top=3pt,bottom=3pt,boxsep=1pt,halign=left]\bfseries\scriptsize Applications Growth, Decay, Newton's Law\end{tcolorbox}
\begin{tcolorbox}[enhanced,colback=popblue,colframe=popblue,coltext=white,boxrule=0.9pt,arc=3pt,left=4pt,right=4pt,top=3pt,bottom=3pt,boxsep=1pt,halign=left]\bfseries\scriptsize Exam Techniques\end{tcolorbox}
\end{tcbraster}

\legende{Mind map of Chapter PMS05: First Order Differential Equations With Separable Variables}
\end{popfigure}

\begin{bilan}
\textbf{Key Definitions}\par
\begin{itemize}
\item \cle{Differential Equation (DE)}: an equation involving an unknown function and its derivatives.
\item \cle{Order}: the highest derivative present (here first order).
\item \cle{Degree}: the power of the highest derivative after clearing fractions and radicals.
\item \cle{Separable DE}: can be written as $\frac{dy}{dx}=g(x)h(y)$ or $M(x)dx+N(y)dy=0$.
\item \cle{General Solution}: a family of solutions containing an arbitrary constant $C$.
\item \cle{Particular Solution}: a specific solution obtained by imposing an \cle{initial condition} $y(x_0)=y_0$.
\end{itemize}

\textbf{Fundamental Formula / Law}\par
If $\frac{dy}{dx}=g(x)h(y)$ with $h(y)\neq0$, then
\[
\int\frac{1}{h(y)}\,dy = \int g(x)\,dx + C.
\]
After integration, solve for $y$ explicitly if possible; otherwise keep the implicit relation.

\textbf{Step‑by‑Step Method}\par
\begin{enumerate}
\item \textbf{Identify} whether the DE is separable: write it in the form $\frac{dy}{dx}=g(x)h(y)$.
\item \textbf{Separate} variables: bring all $y$ terms (including $dy$) to one side and all $x$ terms (including $dx$) to the other.
\item \textbf{Integrate} both sides: $\int\frac{dy}{h(y)} = \int g(x)\,dx + C$. Remember the constant of integration on one side only.
\item \textbf{Simplify} the result: combine logarithms, exponentiate, or solve algebraically for $y$.
\item \textbf{Apply initial condition} (if given): substitute $x_0,y_0$ to find $C$, then write the particular solution.
\item \textbf{Check} by differentiating your solution and verifying it satisfies the original DE.
\end{enumerate}

\textbf{Common Mistakes (Traps)}\par
\begin{itemize}
\item Forgetting the constant $C$ after integration.
\item Dividing by $h(y)$ without considering the case $h(y)=0$ (which may give singular solutions).
\item Incorrect separation: e.g. treating $\frac{dy}{dx}=x+y$ as separable (it is not).
\item Misplacing the constant when exponentiating: $e^{\ln|y|+C}=e^C|y|$, not $|y|+C$.
\item Not stating the domain of the solution, especially when logarithms or square roots appear.
\item Failing to simplify the final answer (e.g. leaving $\ln|y|=x^2+C$ instead of $y=\pm e^C e^{x^2}$).
\end{itemize}

\textbf{GCE Exam Tips}\par
\begin{itemize}
\item Show \emph{every} integration step; examiners award marks for correct separation and integration even if the final algebra has a minor slip.
\item When an initial condition is given, substitute \emph{before} solving for $y$ explicitly if it avoids messy algebra.
\item If the DE models a real situation (population growth, radioactive decay, cooling), define variables clearly and state the model equation (e.g. $\frac{dP}{dt}=kP$).
\item Practise past GCE questions: typical tasks are ``Find the general solution'', ``Find the particular solution given $y(0)=2$'', and ``Interpret the constant $k$ in the context''.
\item Use the variation table (sign of derivative) to discuss the behaviour of solutions when asked to sketch or comment on long‑term trends.
\end{itemize}
\end{bilan}

\findecours

\end{document}
